% Introduction to Motion of Connected Particles — Pure Mathematics with Mechanics Upper Sixth — Cours signé OnBuch+
% Official MINESEC syllabus. Compiler avec : tectonic PMM17-introduction-to-motion-of-connected-particles.tex
\newcommand{\DOCMATIERE}{Pure Mathematics with Mechanics}
\newcommand{\DOCNIVEAU}{Upper Sixth}
\newcommand{\DOCMODULE}{Upper Sixth — Pure mathematics with mechanics}
\newcommand{\DOCLECON}{17}
\newcommand{\DOCTITRE}{Introduction to Motion of Connected Particles}
% OnBuch+ — préambule « cours signés » ENRICHI (compilé avec tectonic / XeLaTeX)
% Système visuel « Pop » d'OnBuch+ : fond crème (jamais blanc), bordures encre
% épaisses, ombres « dures » décalées, titres Archivo Black en capitales.
% Version MATHÉMATIQUES : graphes (pgfplots/tikz), boîtes pédagogiques
% (définition, propriété, méthode, attention, le savais-tu, exercice résolu,
% exercice, TP, bilan), page de garde et signature OnBuch+ ancrée en bas.
%
% Macros attendues dans main.tex AVANT \input{preamble} :
%   \DOCMATIERE  \DOCNIVEAU  \DOCMODULE  \DOCLECON (numéro)  \DOCTITRE
\documentclass[11pt]{article}

\usepackage{fontspec}
\usepackage[english]{babel}
\usepackage[a4paper,margin=2cm,top=3.4cm,bottom=2.8cm,headheight=1.4cm,footskip=1.3cm]{geometry}
\usepackage{amsmath,amssymb,amsfonts}
\usepackage{mathtools} % \xmapsto, \coloneqq... souvent utilisés par les modèles
\usepackage{cancel} % simplifications barrées (bilans de forces)
% \abs{z} (module, valeur absolue) et \norm{u} (norme), utilisés par les modèles
\providecommand{\abs}{}\let\abs\relax\DeclarePairedDelimiter{\abs}{\lvert}{\rvert}
\providecommand{\norm}{}\let\norm\relax\DeclarePairedDelimiter{\norm}{\lVert}{\rVert}
\usepackage[table]{xcolor} % [table] : fournit aussi \rowcolors (alternance de lignes)
\usepackage{pagecolor}
\usepackage{tikz}
\usetikzlibrary{arrows.meta,positioning,calc,shapes.geometric,shapes.misc,shapes.arrows,decorations.pathmorphing,decorations.pathreplacing,decorations.markings,patterns,shadows,mindmap,trees,fit,backgrounds,matrix,intersections,angles,quotes}
\usepackage{pgfplots}
\usepackage[version=4]{mhchem} % équations nucléaires \ce{^{235}_{92}U}
\usepackage[european,siunitx]{circuitikz} % schémas électriques
\pgfplotsset{compat=1.18}
\usepgfplotslibrary{fillbetween,groupplots} % aires entre courbes (calcul intégral)
\usepackage{enumitem}
\usepackage{fancyhdr}
% [most] charge toutes les bibliothèques tcolorbox (listings, minted, raster,
% documentation...) et gonfle inutilement la mémoire TeX sur un document long
% (~300 boîtes) : on ne charge que ce qui est réellement utilisé.
\usepackage[skins,breakable]{tcolorbox}
\usepackage{graphicx}
\usepackage{colortbl}
\usepackage{booktabs}
\usepackage{tabularx}
\usepackage{diagbox} % case d'en-tête diagonale des tableaux à double entrée
% Colonnes extensibles centrée / gauche / droite (C, L, R), souvent utilisées par les modèles
\newcolumntype{C}{>{\centering\arraybackslash}X}
\newcolumntype{L}{>{\raggedright\arraybackslash}X}
\newcolumntype{R}{>{\raggedleft\arraybackslash}X}
\usepackage{array}
\usepackage{multicol}
\usepackage{siunitx}
% parse-numbers=false : accepte aussi \SI{2,5 \times 10^{-3}}{...}
\sisetup{locale=UK,output-decimal-marker={.},group-minimum-digits=4,parse-numbers=false}
\DeclareSIUnit\ohm{\ensuremath{\symup{\Omega}}} % U+2126 absent de la police
\DeclareSIUnit\division{div} % réglages d'oscilloscope (ms/div, V/div)
\DeclareSIUnit\tr{tr} % tours (vitesse de rotation, ex. \SI{1500}{\tr\per\minute})
\usepackage{qrcode}
% \degree : commande fréquemment utilisée par les modèles pour un angle ou une
% température en dehors de \SI{}{} (non fournie par les paquets chargés).
\providecommand{\degree}{^{\circ}}
% \qed (fin de démonstration), utilisé par les modèles sans amsthm
\providecommand{\qed}{\hfill$\square$}
% \barre : double barre des valeurs interdites dans un tableau de variations (tkz-tab)
\providecommand{\barre}{\ensuremath{\|}}
% environnement proof (amsthm non chargé) utilisé par les modèles
\ifdefined\proof\else\newenvironment{proof}[1][Proof]{\par\noindent\textit{#1.}\ }{\qed\par}\fi
\usepackage{lastpage}
\usepackage{hyperref}
\hypersetup{hidelinks}

% ============================================================
% Palette « Pop » OnBuch+ — mêmes hex que la classe P (plus_ui.dart)
% ============================================================
\definecolor{poporange}{HTML}{F59321}
\definecolor{popdark}{HTML}{E07A0C}
\definecolor{popcream}{HTML}{FFF6E8}
\definecolor{popink}{HTML}{1C1714}
\definecolor{poppurple}{HTML}{7A5AE0}
\definecolor{popgreen}{HTML}{2E9E6A}
\definecolor{poppink}{HTML}{E0457B}
\definecolor{popblue}{HTML}{2F7DE1}
\definecolor{popgold}{HTML}{C98A12}
\definecolor{popmuted}{HTML}{8A7F76}
\definecolor{popline}{HTML}{EADFD2}
\definecolor{popgray}{HTML}{8A7F76}
\colorlet{popgrey}{popgray}
% Alias tolérants (le modèle nomme parfois les couleurs différemment)
\colorlet{popwhite}{white}
\colorlet{popblack}{popink}
\colorlet{popred}{poppink}
\colorlet{popyellow}{popgold}
% Teintes claires pour fonds de tableaux / zones
\colorlet{popblueL}{popblue!10!white}
\colorlet{poporangeL}{poporange!14!white}
\colorlet{popgreenL}{popgreen!12!white}
\colorlet{poppurpleL}{poppurple!12!white}
\colorlet{poppinkL}{poppink!12!white}
\colorlet{popgoldL}{popgold!14!white}

\pagecolor{popcream}

% ============================================================
% Typographie
% ============================================================
\defaultfontfeatures{Ligatures=TeX}
\setmainfont{PlusJakartaSans-Medium.ttf}[Path=../../../fonts/,BoldFont=PlusJakartaSans-Bold.ttf]
% Maths en sans-serif (Fira Math) avec chiffres et lettres droites de Plus
% Jakarta, pour que les formules se fondent dans le texte.
\usepackage[math-style=ISO,bold-style=ISO]{unicode-math}
\setmathfont{FiraMath-Regular.otf}[Path=../../../fonts/]
\setmathfont{PlusJakartaSans-Medium.ttf}[Path=../../../fonts/,range={up/num,up/{Latin,latin},\mathup}]
% (icomma retiré : décimales avec point en anglais)
% Caractères mathématiques Unicode absents des polices de texte (Plus Jakarta,
% Archivo Black) : rendus par leur équivalent mathématique, en texte comme en
% formule — sinon ils s'affichent en carré vide (ex. « Arithmétique dans ℤ »).
\usepackage{newunicodechar}
\newunicodechar{ℝ}{\ensuremath{\mathbb{R}}}
\newunicodechar{ℤ}{\ensuremath{\mathbb{Z}}}
\newunicodechar{ℂ}{\ensuremath{\mathbb{C}}}
\newunicodechar{ℕ}{\ensuremath{\mathbb{N}}}
\newunicodechar{ℚ}{\ensuremath{\mathbb{Q}}}
\newunicodechar{𝔻}{\ensuremath{\mathbb{D}}}
\newunicodechar{α}{\ensuremath{\alpha}}
\newunicodechar{β}{\ensuremath{\beta}}
\newunicodechar{γ}{\ensuremath{\gamma}}
\newunicodechar{δ}{\ensuremath{\delta}}
\newunicodechar{ε}{\ensuremath{\varepsilon}}
\newunicodechar{θ}{\ensuremath{\theta}}
\newunicodechar{λ}{\ensuremath{\lambda}}
\newunicodechar{μ}{\ensuremath{\mu}}
\newunicodechar{σ}{\ensuremath{\sigma}}
\newunicodechar{τ}{\ensuremath{\tau}}
\newunicodechar{φ}{\ensuremath{\varphi}}
\newunicodechar{ψ}{\ensuremath{\psi}}
\newunicodechar{ω}{\ensuremath{\omega}}
\newunicodechar{Σ}{\ensuremath{\Sigma}}
% majuscules produites par \MakeUppercase dans les titres (α-aminés -> Α)
\newunicodechar{Α}{\ensuremath{\alpha}}
\newunicodechar{Β}{\ensuremath{\beta}}
\newunicodechar{Γ}{\ensuremath{\Gamma}}
\newunicodechar{Δ}{\ensuremath{\Delta}}
\newunicodechar{Ε}{\ensuremath{\varepsilon}}
\newunicodechar{Θ}{\ensuremath{\Theta}}
\newunicodechar{Λ}{\ensuremath{\Lambda}}
\newunicodechar{Μ}{\ensuremath{\mu}}
\newunicodechar{Τ}{\ensuremath{\tau}}
\newunicodechar{Φ}{\ensuremath{\Phi}}
\newunicodechar{Ψ}{\ensuremath{\Psi}}
\newunicodechar{Ω}{\ensuremath{\Omega}}
\newunicodechar{↦}{\ensuremath{\mapsto}}
\newunicodechar{∈}{\ensuremath{\in}}
\newunicodechar{∉}{\ensuremath{\notin}}
\newunicodechar{∧}{\ensuremath{\wedge}}
\newunicodechar{⇔}{\ensuremath{\Leftrightarrow}}
\newunicodechar{⟺}{\ensuremath{\Longleftrightarrow}}
\newunicodechar{⇒}{\ensuremath{\Rightarrow}}
\newunicodechar{⟹}{\ensuremath{\Longrightarrow}}
\newunicodechar{∘}{\ensuremath{\circ}}
\newunicodechar{∪}{\ensuremath{\cup}}
\newunicodechar{∩}{\ensuremath{\cap}}
\newunicodechar{≡}{\ensuremath{\equiv}}
\newunicodechar{⊂}{\ensuremath{\subset}}
\newunicodechar{⊥}{\ensuremath{\perp}}
\newunicodechar{∃}{\ensuremath{\exists}}
\newunicodechar{∀}{\ensuremath{\forall}}
\newunicodechar{∛}{\ensuremath{\sqrt[3]{\,}}}
\newunicodechar{∜}{\ensuremath{\sqrt[4]{\,}}}
\newunicodechar{⃗}{\ensuremath{{}^{\to}}}
\newunicodechar{ⁿ}{\ifmmode^{n}\else\textsuperscript{\ensuremath{n}}\fi}
\newunicodechar{ˣ}{\ifmmode^{x}\else\textsuperscript{\ensuremath{x}}\fi}
\newunicodechar{ᵇ}{\ifmmode^{b}\else\textsuperscript{\ensuremath{b}}\fi}
\newunicodechar{ʳ}{\ifmmode^{r}\else\textsuperscript{\ensuremath{r}}\fi}
\newunicodechar{ᵘ}{\ifmmode^{u}\else\textsuperscript{\ensuremath{u}}\fi}
\newunicodechar{ʸ}{\ifmmode^{y}\else\textsuperscript{\ensuremath{y}}\fi}
\newunicodechar{ᵃ}{\ifmmode^{a}\else\textsuperscript{\ensuremath{a}}\fi}
\newunicodechar{ᵉ}{\ifmmode^{e}\else\textsuperscript{\ensuremath{e}}\fi}
\newunicodechar{ᵏ}{\ifmmode^{k}\else\textsuperscript{\ensuremath{k}}\fi}
\newunicodechar{ᵖ}{\ifmmode^{p}\else\textsuperscript{\ensuremath{p}}\fi}
\newunicodechar{ᴺ}{\ifmmode^{N}\else\textsuperscript{\ensuremath{N}}\fi}
\newunicodechar{ᶜ}{\ifmmode^{c}\else\textsuperscript{\ensuremath{c}}\fi}
\newunicodechar{ʰ}{\ifmmode^{h}\else\textsuperscript{\ensuremath{h}}\fi}
\newunicodechar{ᵠ}{\ifmmode^{q}\else\textsuperscript{\ensuremath{q}}\fi}
\newunicodechar{ₙ}{\ifmmode_{n}\else\textsubscript{\ensuremath{n}}\fi}
\newunicodechar{ₐ}{\ifmmode_{a}\else\textsubscript{\ensuremath{a}}\fi}
\newunicodechar{ᵢ}{\ifmmode_{i}\else\textsubscript{\ensuremath{i}}\fi}
\newunicodechar{ᵣ}{\ifmmode_{r}\else\textsubscript{\ensuremath{r}}\fi}
\newunicodechar{ₖ}{\ifmmode_{k}\else\textsubscript{\ensuremath{k}}\fi}
\newunicodechar{ᵦ}{\ifmmode_{\beta}\else\textsubscript{\ensuremath{\beta}}\fi}
\newunicodechar{ₓ}{\ifmmode_{x}\else\textsubscript{\ensuremath{x}}\fi}
\newfontfamily\popdisplay{ArchivoBlack-Regular.ttf}[Path=../../../fonts/]
\newfontfamily\poplogo{SpaceGrotesk-Bold.ttf}[Path=../../../fonts/]
\newfontfamily\popsemi{PlusJakartaSans-SemiBold.ttf}[Path=../../../fonts/]
\newfontfamily\popxbold{PlusJakartaSans-ExtraBold.ttf}[Path=../../../fonts/]
\newfontfamily\popxboldls{PlusJakartaSans-ExtraBold.ttf}[Path=../../../fonts/,LetterSpace=3.0]

\setlist[enumerate]{leftmargin=1.7em,itemsep=5pt,topsep=4pt}
\setlist[itemize]{leftmargin=1.7em,itemsep=4pt,topsep=4pt,label={\color{poporange}\textbullet}}
\setlength{\parindent}{0pt}
\setlength{\parskip}{5pt}
\renewcommand{\arraystretch}{1.25}

% pgfplots : style Pop par défaut
\pgfplotsset{
  every axis/.append style={axis line style={popink,line width=1pt},tick style={popink},
    grid style={popline,line width=0.5pt},label style={font=\small},tick label style={font=\footnotesize}},
  popcurve/.style={poporange,line width=1.8pt,no markers,smooth},
}
% Même style disponible dans un \draw TikZ simple (hors pgfplots)
\tikzset{popcurve/.style={poporange,line width=1.8pt}}
\tikzset{popgrid/.style={popline,very thin}}
\colorlet{popcurve}{poporange} % le nom du style est parfois utilisé comme couleur

% ============================================================
% Pill / squiggle
% ============================================================
\newcommand{\poppill}[3]{%
  \tikz[baseline=-0.55ex]\node[draw=popink,line width=1.2pt,rounded corners=2.6pt,%
    fill=#2,inner xsep=6.5pt,inner ysep=3pt]{%
    \popxboldls\fontsize{7}{7}\selectfont\color{#3}\MakeUppercase{#1}};%
}
\newcommand{\popsquiggle}[1]{%
  \tikz[baseline=-0.4ex]{
    \draw[#1,line width=2.6pt,line cap=round]
      plot [smooth] coordinates {(0,0.09) (0.34,0.01) (0.68,0.21) (1.02,0.01) (1.36,0.10)};
  }%
}
% Mot-clé surligné (marqueur orange sous le texte)
\newcommand{\cle}[1]{{\popxbold\color{popdark}#1}}

% ============================================================
% En-tête / pied
% ============================================================
\pagestyle{fancy}
\fancyhf{}
\renewcommand{\headrulewidth}{0pt}
\renewcommand{\footrulewidth}{0pt}
\lhead{\raisebox{-0.15ex}{\poplogo\fontsize{13}{13}\selectfont\color{popink}OnBuch\color{poporange}+}}
\rhead{\poppill{\DOCMATIERE\ \textbullet\ \DOCNIVEAU\ \textbullet\ Chapter \DOCLECON}{white}{popink}}
\lfoot{\popxbold\fontsize{7.6}{7.6}\selectfont\color{popmuted}\MakeUppercase{Signed course OnBuch+}\ \textbullet\ {\popsemi\DOCTITRE}}
\rfoot{\tikz[baseline=-0.6ex]\node[rounded corners=8.5pt,fill=popink,minimum height=17pt,minimum width=17pt,inner xsep=5pt,inner sep=0]{\popxbold\fontsize{8}{8}\selectfont\color{popcream}\thepage\,/\,\pageref*{LastPage}};}

% ============================================================
% Titres
% ============================================================
\newcommand{\chaptitre}[2]{%
  \begin{tcolorbox}[enhanced,colback=poporange,colframe=popink,boxrule=2.6pt,arc=11pt,
      drop shadow=popink,left=15pt,right=15pt,top=12pt,bottom=11pt,before skip=3pt,after skip=18pt]
    \poppill{#2}{popink}{popcream}\par\vspace{9pt}
    {\raggedright\hyphenpenalty=10000\exhyphenpenalty=10000\popdisplay\fontsize{22}{24}\selectfont\color{popcream}\MakeUppercase{#1}\par}
  \end{tcolorbox}
}
% Section numérotée : pastille orange avec le numéro + titre Archivo
\newcounter{popsec}
\newcommand{\coursec}[1]{%
  \stepcounter{popsec}\setcounter{popsub}{0}%
  \par\vspace{16pt}\noindent
  \begin{minipage}{\linewidth}
  \tikz[baseline=-0.7ex]\node[draw=popink,line width=1.6pt,fill=poporange,rounded corners=4pt,minimum size=22pt,inner sep=0]{\popdisplay\fontsize{12}{12}\selectfont\color{popcream}\thepopsec};%
  \hspace{8pt}{\popdisplay\fontsize{15}{17}\selectfont\color{popink}\MakeUppercase{#1}}\par
  \vspace{4pt}\noindent{\color{poporange}\rule{36pt}{3.6pt}}
  \end{minipage}\par\nopagebreak\vspace{8pt}
}
\newcounter{popsub}
\newcommand{\courssub}[1]{%
  \stepcounter{popsub}%
  \par\vspace{9pt}\noindent{\popxbold\fontsize{11.5}{13}\selectfont\color{popink}{\color{poporange}\thepopsec.\thepopsub}\hspace{6pt}#1}\par\nopagebreak\vspace{4pt}
}

% ============================================================
% Boîtes pédagogiques — même recette : fond blanc, bordure épaisse colorée,
% ombre dure encre, badge plein. Titre optionnel [#1].
% ============================================================
\tcbset{popbase/.style={breakable,enhanced,colback=white,boxrule=2.3pt,arc=9pt,
  shadow={3pt}{-3pt}{0pt}{popink},left=14pt,right=14pt,top=11pt,bottom=11pt,before skip=11pt,after skip=15pt,
  fontupper=\popsemi\fontsize{10.6}{14.5}\selectfont\color{popink},
  fontlower=\popsemi\fontsize{10.6}{14.5}\selectfont\color{popink}}}
% Coche dessinée (le glyphe ✓ n'existe pas dans les polices du document)
\newcommand{\popcheck}[1]{\tikz[baseline=-0.5ex]\draw[#1,line width=1.6pt,line cap=round,line join=round](-0.13,0.0)--(-0.04,-0.09)--(0.13,0.1);}
\providecommand{\coche}{\popcheck{popgreen}}
\providecommand{\resultat}[1]{\par\smallskip\noindent\fcolorbox{popgreen}{popgreenL}{\parbox{\dimexpr\linewidth-2\fboxsep-2\fboxrule\relax}{\textbf{Result.} #1}}\par\smallskip}
\newcommand{\popboxhead}[3]{\poppill{#1}{#2}{white}\ifx\relax#3\relax\else\hspace{6pt}{\popxbold\fontsize{10.6}{12}\selectfont\color{popink}#3}\fi\par\vspace{7pt}}

\newtcolorbox{aretenir}[1][]{popbase,colframe=popgreen,before upper={\popboxhead{Key points}{popgreen}{#1}}}
\newtcolorbox{exemplebox}[1][]{popbase,colframe=poporange,before upper={\popboxhead{Exemple}{poporange}{#1}}}
\newtcolorbox{definition}[1][]{popbase,colframe=popblue,before upper={\popboxhead{Definition}{popblue}{#1}}}
\newtcolorbox{propriete}[1][]{popbase,colframe=poppurple,before upper={\popboxhead{Property}{poppurple}{#1}}}
\newtcolorbox{methode}[1][]{popbase,colframe=popgold,before upper={\popboxhead{Method}{popgold}{#1}}}
\newtcolorbox{attention}[1][]{popbase,colframe=poppink,before upper={\popboxhead{Warning}{poppink}{#1}}}
\newtcolorbox{experience}[1][]{popbase,colframe=popblue,colback=popblueL,before upper={\popboxhead{Experiment}{popblue}{#1}}}
% « Le savais-tu ? » : carte violette pleine (contexte, culture, Cameroun)
\newtcolorbox{savaistu}[1][]{popbase,colback=poppurple,colframe=popink,
  fontupper=\popsemi\fontsize{10.4}{14.2}\selectfont\color{white},
  before upper={\poppill{Did you know?}{white}{poppurple}\ifx\relax#1\relax\else\hspace{6pt}{\popxbold\color{white}#1}\fi\par\vspace{7pt}}}
% Exercice résolu : énoncé en haut, \tcblower puis la solution sur fond crème
\newcounter{popexo}\newcounter{popexr}
\newtcolorbox{exoresolu}[1][]{popbase,colframe=poporange,colbacklower=popcream,
  segmentation style={popink,line width=1.2pt,dashed},
  before upper={\stepcounter{popexr}\popboxhead{Worked example \thepopexr}{poporange}{#1}},
  before lower={\poppill{Solution}{popink}{popcream}\par\vspace{6pt}}}
% Exercice (non résolu en place) — la correction est dans la partie Corrigés
\newtcolorbox{exercice}[1][]{popbase,colframe=popink,
  before upper={\stepcounter{popexo}\popboxhead{Exercise \thepopexo}{popink}{#1}}}
% Corrigé
\newtcolorbox{corrige}[1][]{popbase,colframe=popgreen,colback=popgreenL,
  before upper={\popboxhead{Solution}{popgreen}{#1}}}
% Objectifs / prérequis / bilan
\newtcolorbox{objectifs}{popbase,colframe=popink,colback=poporangeL,
  before upper={\popboxhead{Chapter objectives}{popink}{}}}
\newtcolorbox{prerequis}{popbase,colframe=popink,colback=popblueL,
  before upper={\popboxhead{Prerequisites}{popblue}{}}}
\newtcolorbox{bilan}{popbase,colframe=popink,colback=white,boxrule=2.8pt,
  before upper={\popboxhead{Fiche bilan}{poporange}{L'essentiel à connaître pour l'examen}}}
% Figure encadrée avec légende
\newtcolorbox{popfigure}[1][]{enhanced,breakable=false,colback=white,colframe=popline,boxrule=1.4pt,arc=8pt,
  left=10pt,right=10pt,top=10pt,bottom=8pt,before skip=10pt,after skip=12pt,halign=center,
  lower separated=false,halign lower=center,
  fontlower=\popsemi\fontsize{9}{11.5}\selectfont\color{popmuted},
  #1}
\newcommand{\legende}[1]{\par\vspace{6pt}{\popsemi\fontsize{9}{11.5}\selectfont\color{popmuted}#1}}

% ============================================================
% Signature OnBuch+
% ============================================================
\newcommand{\poplogoinline}[1]{{\poplogo\fontsize{#1}{#1}\selectfont\color{popink}OnBuch\color{poporange}+}}
\newcommand{\signatureonbuch}{%
  \begin{tcolorbox}[enhanced,colback=popink,colframe=popink,boxrule=2.2pt,arc=10pt,
      drop shadow=poporange,left=14pt,right=14pt,top=10pt,bottom=10pt,before skip=0pt,after skip=0pt]
    \tikz[baseline=-0.7ex]\node[circle,fill=popgreen,draw=popcream,line width=1.3pt,minimum size=22pt,inner sep=0]{\popcheck{white}};%
    \hspace{9pt}\begin{minipage}[c]{0.62\linewidth}
      {\popxbold\fontsize{12}{13}\selectfont\color{popcream}Signed\ \textbullet\ {\poplogo OnBuch\color{poporange}+}}\par\vspace{2pt}
      {\raggedright\popsemi\fontsize{8}{10.5}\selectfont\color{popline}\MakeUppercase{OnBuch+ teaching team}\\\MakeUppercase{Follows the official MINESEC syllabus}\par}
    \end{minipage}\hfill
    {\poplogo\fontsize{22}{22}\selectfont\color{popcream}OnBuch\color{poporange}+}
  \end{tcolorbox}%
}

% ============================================================
% Page de garde
% #1 titre · #2 matière · #3 niveau · #4 module · #5 n° leçon · #6 durée ·
% #7 description / objectifs (texte ou itemize)
% ============================================================
\newcommand{\pagegarde}[7]{%
  \thispagestyle{empty}
  \newgeometry{a4paper,margin=1.7cm,top=1.6cm,bottom=1.4cm}%
  \begin{tikzpicture}[remember picture,overlay]
    \fill[poporange!18] ([xshift=-3.2cm,yshift=-3.6cm]current page.north east) circle (4.6cm);
    \fill[poppurple!14] ([xshift=2.4cm,yshift=7.6cm]current page.south west) circle (3.8cm);
  \end{tikzpicture}%
  {\poplogo\fontsize{30}{30}\selectfont\color{popink}OnBuch\color{poporange}+}\hfill
  \raisebox{0.6ex}{\poppill{Signed course \textbullet\ Premium}{popink}{popcream}}\par
  \vspace{1pt}\popsquiggle{poppurple}\par
  \vspace{8pt}
  \begin{tcolorbox}[enhanced,colback=poporange,colframe=popink,boxrule=2.8pt,arc=13pt,
      drop shadow=popink,left=18pt,right=18pt,top=12pt,bottom=13pt,after skip=10pt]
    \poppill{#2 \textbullet\ #3}{popink}{popcream}\hspace{5pt}\poppill{#4}{white}{popink}\par\vspace{8pt}
    {\popdisplay\fontsize{14}{14}\selectfont\color{popink}CHAPTER #5}\par\vspace{4pt}
    {\raggedright\hyphenpenalty=10000\exhyphenpenalty=10000\popdisplay\fontsize{25}{27}\selectfont\color{popcream}\MakeUppercase{#1}\par}
    \vspace{8pt}
    {\popxbold\fontsize{9.5}{9.5}\selectfont\color{popink}\MakeUppercase{Suggested time: #6}}
  \end{tcolorbox}
  \begin{tcolorbox}[enhanced,colback=white,colframe=popink,boxrule=2.2pt,arc=10pt,
      drop shadow=popink,left=16pt,right=16pt,top=10pt,bottom=10pt,after skip=8pt]
    \poppill{In this chapter}{poppurple}{white}\par\vspace{5pt}
    \popsemi\fontsize{9.3}{12.6}\selectfont\color{popink}#7
  \end{tcolorbox}
  \noindent\begin{minipage}[t]{0.487\linewidth}
    \begin{tcolorbox}[enhanced,colback=popgreen,colframe=popink,boxrule=2.2pt,arc=10pt,
        drop shadow=popink,left=11pt,right=11pt,top=10pt,bottom=10pt,height=3.1cm,valign=center]
      \tcbox[colback=white,colframe=popink,boxrule=1.3pt,arc=5pt,boxsep=0pt,nobeforeafter,
        left=3pt,right=3pt,top=3pt,bottom=3pt,valign=center]{\qrcode[height=1.9cm]{https://play.google.com/store/apps/details?id=com.luvvix.onbuch&hl=en}}%
      \hspace{8pt}\begin{minipage}[c]{0.5\linewidth}
        {\popdisplay\fontsize{11}{12.5}\selectfont\color{white}\MakeUppercase{Download OnBuch}}\par\vspace{4pt}
        {\raggedright\popsemi\fontsize{8.4}{11}\selectfont\color{white}Free \textbullet\ Google Play\\AI tutor, past papers, results\par}
      \end{minipage}
    \end{tcolorbox}
  \end{minipage}\hfill
  \begin{minipage}[t]{0.487\linewidth}
    \begin{tcolorbox}[enhanced,colback=poppurple,colframe=popink,boxrule=2.2pt,arc=10pt,
        drop shadow=popink,left=11pt,right=11pt,top=10pt,bottom=10pt,height=3.1cm,valign=center]
      \tikz[baseline=-0.7ex]\node[circle,fill=white,draw=popink,line width=1.3pt,minimum size=1.6cm,inner sep=0]{\popdisplay\fontsize{24}{24}\selectfont\color{poppurple}+};%
      \hspace{8pt}\begin{minipage}[c]{0.6\linewidth}
        {\popdisplay\fontsize{11}{12.5}\selectfont\color{white}\MakeUppercase{Go OnBuch+}}\par\vspace{4pt}
        {\raggedright\popsemi\fontsize{8.4}{11}\selectfont\color{white}All signed courses, unlimited AI tutor, PDF sheets — OnBuch+ tab in the app\par}
      \end{minipage}
    \end{tcolorbox}
  \end{minipage}
  \vspace{8pt}
  \popcontenu
  \vfill
  \signatureonbuch
  \restoregeometry
  \newpage
}

% Rangée « Ce cours contient » (page de garde)
\newcommand{\poptile}[3]{% #1 couleur #2 gros texte #3 légende
  \begin{tcolorbox}[enhanced,colback=#1,colframe=popink,boxrule=2pt,arc=9pt,shadow={2.5pt}{-2.5pt}{0pt}{popink},
      left=6pt,right=6pt,top=9pt,bottom=9pt,halign=center,valign=center,height=1.75cm,width=\linewidth,nobeforeafter]
    {\popdisplay\fontsize{12.5}{13}\selectfont\color{white}\MakeUppercase{#2}}\par\vspace{3pt}
    {\centering\popsemi\fontsize{7.8}{9.5}\selectfont\color{white}#3\par}
  \end{tcolorbox}}
\newcommand{\popcontenu}{%
  \par\noindent{\popxboldls\fontsize{7.5}{7.5}\selectfont\color{popmuted}THIS SIGNED COURSE CONTAINS}\par\vspace{7pt}
  \noindent\begin{minipage}[t]{0.235\linewidth}\poptile{popblue}{Course}{complete, illustrated, step by step}\end{minipage}\hfill
  \begin{minipage}[t]{0.235\linewidth}\poptile{poporange}{Methods}{and worked examples}\end{minipage}\hfill
  \begin{minipage}[t]{0.235\linewidth}\poptile{poppink}{Exercises}{exam style + solutions}\end{minipage}\hfill
  \begin{minipage}[t]{0.235\linewidth}\poptile{popgreen}{Summary sheet}{the essentials to remember}\end{minipage}\par}

% Sommaire
\newtcolorbox{sommaire}{enhanced,colback=white,colframe=popink,boxrule=2.2pt,arc=10pt,
  drop shadow=popink,left=18pt,right=18pt,top=13pt,bottom=13pt,before skip=6pt,after skip=20pt,
  before upper={\poppill{Contents}{poppurple}{white}\par\vspace{9pt}\popsemi\fontsize{10.6}{14.5}\selectfont\color{popink}}}

% Signature de fin de cours — poussée tout en bas de la dernière page
\newcommand{\findecours}{%
  \par\vfill\nopagebreak
  \begin{center}\popsquiggle{poporange}\end{center}\vspace{4pt}
  \signatureonbuch
}
\AtBeginDocument{\def\llbracket{\mathopen{[\mkern-3.5mu[}}\def\rrbracket{\mathclose{]\mkern-3.5mu]}}} % intervalles d'entiers (glyphe ⟦ absent de la police)
% Glyphes absents de Fira Math : redéfinis à partir de glyphes disponibles
\AtBeginDocument{%
  \def\setminus{\mathbin{\backslash}}\let\smallsetminus\setminus
  \def\vdots{\mathord{\vbox{\baselineskip3.2pt\lineskiplimit0pt\kern4pt\hbox{.}\hbox{.}\hbox{.}}}}%
  \def\ddots{\mathinner{\mkern1mu\raise6.5pt\hbox{.}\mkern2mu\raise3.6pt\hbox{.}\mkern2mu\raise0.7pt\hbox{.}\mkern1mu}}%
  \def\triangle{\mathord{\scalebox{0.9}{$\symup{\Delta}$}}}%
  \def\nabla{\mathord{\raisebox{\depth}{\scalebox{1}[-1]{$\symup{\Delta}$}}}}%
  \def\checkmark{\popcheck{popdark}}%
  \def\star{\mathbin{\ast}}\def\bigstar{\mathord{\ast}}%
  \def\diamond{\mathbin{\scalebox{0.6}{$\Diamond$}}}%
  \def\bigcup{\mathop{\mathchoice{\raisebox{-0.3em}{\scalebox{1.7}{$\cup$}}}{\scalebox{1.25}{$\cup$}}{\cup}{\cup}}\displaylimits}%
  \def\bigcap{\mathop{\mathchoice{\raisebox{-0.3em}{\scalebox{1.7}{$\cap$}}}{\scalebox{1.25}{$\cap$}}{\cap}{\cap}}\displaylimits}%
  \def\lesseqqgtr{\mathrel{\lessgtr}}\def\bigoplus{\mathop{\oplus}\displaylimits}%
}
\providecommand{\ssi}{if and only if} % abréviation fréquente
\AtBeginDocument{\providecommand{\wideparen}{\overparen}} % arc de cercle

%% FIGFIX-BEGIN v3 — correctifs globaux des figures (docs/onbuch-generation/tools/figfix)
\usepackage{adjustbox}\usepackage{etoolbox}
% toute figure TikZ/pgfplots plus large que la ligne est réduite à la largeur disponible
\BeforeBeginEnvironment{tikzpicture}{\begin{adjustbox}{max width=\linewidth}}
\AfterEndEnvironment{tikzpicture}{\end{adjustbox}}
% légendes pgfplots lisibles par-dessus les courbes
\pgfplotsset{every axis/.append style={legend style={fill=white,fill opacity=0.92,text opacity=1,draw=black!15,inner sep=3pt}}}
% étiquettes d'arêtes (syntaxe edge["..."]) avec fond blanc
\tikzset{every edge quotes/.append style={fill=white,inner sep=1.2pt}}
% bulles de cartes mentales : texte à la ligne dans la bulle, bulle un peu plus grande
\tikzset{every concept/.append style={text width=2.0cm,align=center,minimum size=2.3cm,inner sep=2pt}}
%% FIGFIX-END
\begin{document}

\pagegarde{Introduction to Motion of Connected Particles}{Pure Mathematics with Mechanics}{Upper Sixth}{Upper Sixth — Pure mathematics with mechanics}{17}{9 periods}{This chapter introduces the dynamics of two or more bodies linked by strings, rods or contact, moving together under gravity, friction and applied forces. Students learn to model connected systems (smooth and rough surfaces, inclined planes, pulleys and Atwood's machine), to compute accelerations and tensions, and to solve real-life problems using both Newton's laws and energy methods.
\vspace{4pt}
\begin{multicols}{2}\small
\begin{itemize}
\item By the end of this chapter I can model a system of connected particles with clear assumptions (light inextensible string, smooth pulley, particle model).
\item By the end of this chapter I can draw separate force diagrams for each particle of a connected system.
\item By the end of this chapter I can apply Newton's second law to each particle and solve the simultaneous equations for acceleration and tension.
\item By the end of this chapter I can analyse motion of connected particles on smooth horizontal surfaces.
\item By the end of this chapter I can include limiting and non-limiting friction (F = μR) in connected-particle problems on rough surfaces.
\item By the end of this chapter I can resolve forces on an inclined plane and analyse particles connected over a pulley on slopes.
\end{itemize}\end{multicols}}

\begin{sommaire}
\begin{multicols}{2}
\begin{enumerate}
\item Section 1: Modelling Connected Particles and Tension in Strings
\item Section 2: Motion on Smooth and Rough Horizontal Surfaces
\item Section 3: Motion on Inclined Planes
\item Section 4: Atwood's Machine and Multi-Stage Motion
\item Section 5: Pulley Systems with More Than Two Masses and Variable Acceleration
\item Section 6: Energy Methods and Real-Life Problems
\item Integration activity
\item Methods and worked examples
\item Exercises
\item Solutions to the exercises
\item Summary sheet
\end{enumerate}
\end{multicols}
\end{sommaire}

\begin{objectifs}
\begin{itemize}
\item By the end of this chapter I can model a system of connected particles with clear assumptions (light inextensible string, smooth pulley, particle model).
\item By the end of this chapter I can draw separate force diagrams for each particle of a connected system.
\item By the end of this chapter I can apply Newton's second law to each particle and solve the simultaneous equations for acceleration and tension.
\item By the end of this chapter I can analyse motion of connected particles on smooth horizontal surfaces.
\item By the end of this chapter I can include limiting and non-limiting friction (F = μR) in connected-particle problems on rough surfaces.
\item By the end of this chapter I can resolve forces on an inclined plane and analyse particles connected over a pulley on slopes.
\item By the end of this chapter I can solve Atwood's machine problems, including the force on the pulley support.
\item By the end of this chapter I can handle systems with more than two masses and multi-stage motion (string going slack, particle hitting the ground).
\item By the end of this chapter I can use the work–energy principle and conservation of mechanical energy as an alternative method for connected systems.
\item By the end of this chapter I can interpret and solve real-life connected-particle problems in the style of GCE Paper 2 questions.
\end{itemize}
\end{objectifs}

\begin{prerequis}
\begin{itemize}
\item Newton's three laws of motion and F = ma (Lower Sixth mechanics)
\item Resolving forces into components and equilibrium of a particle
\item Friction: coefficient of friction, F ≤ μR, limiting equilibrium
\item Kinematics with constant acceleration: the suvat equations
\item Work, kinetic energy, potential energy and the work–energy principle
\end{itemize}
\end{prerequis}

\newpage

\coursec{Section 1: Modelling Connected Particles and Tension in Strings}

At the Mile 4 motor park in Limbe, a mechanic uses a rope passing over a beam to lift a heavy engine out of a lorry: the engine goes down slowly while a counterweight of scrap metal rises. Why does the engine not crash to the ground, and how heavy must the counterweight be? On the Bamenda escarpment, a pickup tows a broken-down taxi up a steep slope with a chain --- will the chain hold? These everyday situations all involve \cle{connected particles}, and by the end of this chapter you will be able to predict their motion, the forces in the connectors, and the energy involved, exactly as required at the GCE Advanced Level.

In this first section we build the language of the whole chapter: what a connected-particle system is, the modelling assumptions that make it tractable, the nature of \cle{tension}, and a five-step routine that turns any such problem into a pair of simultaneous equations.

\courssub{What is a Connected-Particle System?}

So far in mechanics you have mostly studied a \emph{single} particle moving under given forces. In real life, bodies rarely move alone: a locomotive pulls wagons, a crane cable lifts a container, two climbers are tied together by a rope. Whenever two or more bodies are linked by a string, rope, chain, rod or tow-bar, the motion of one body \emph{constrains} the motion of the others: they cannot move independently.

\begin{definition}[Connected particles]
Two or more bodies are called \cle{connected particles} when they are linked by a connector (string, rope, chain, rod or tow-bar) in such a way that the motion of one body determines, wholly or partly, the motion of the other(s). In the standard model, all particles attached to the same taut, inextensible string move with the \textbf{same speed} and the \textbf{same magnitude of acceleration}.
\end{definition}

Think of the mechanic's rope in Limbe: if the engine descends by $30\ \mathrm{cm}$, the counterweight must rise by exactly $30\ \mathrm{cm}$, because the rope does not stretch. The two motions are locked together --- that is the whole point of the model.

\courssub{Modelling Assumptions}

Real ropes have mass, stretch a little, and rub on rough pulleys. To obtain clean mathematics we idealise. Each assumption below has a precise physical consequence that you must be able to state in an examination.

\begin{aretenir}[The standard modelling assumptions]
\begin{itemize}
\item \textbf{Particle model:} each body is treated as a point mass. Its size, shape and rotation are ignored; all forces act at a single point.
\item \textbf{Light string (massless):} the string has no mass, so no net force is needed to accelerate the string itself. Consequence: the \cle{tension} has the \emph{same value at every point} of the string. (Proof: for a massless string, $F_{\text{net}} = m a = 0$, so the forces at its ends are equal and opposite.)
\item \textbf{Inextensible string:} the string does not stretch. Consequence: the distances moved, the speeds and the magnitudes of the accelerations of all particles attached to the same taut string are \emph{equal}.
\item \textbf{Smooth pulley or peg:} there is no friction in the pulley (and the pulley is light). Consequence: the tension is \emph{the same on both sides} of the pulley; the pulley only changes the \emph{direction} of the force.
\end{itemize}
\end{aretenir}

\begin{attention}[Inextensibility is a kinematic condition]
A common error is to treat the two particles as independent and give them different accelerations. If the string is taut and inextensible, this is impossible: the string would have to stretch or break. Always write one single unknown $a$ for the common magnitude of acceleration, and one single unknown $v$ for the common speed.
\end{attention}

\courssub{Tension in Strings}

Take a taut string and pull it at both ends. Every small piece of the string is pulled by its neighbours; this internal pulling force, transmitted along the string, is the \cle{tension}.

\begin{definition}[Tension]
The \cle{tension} $T$ in a taut string is the pulling force that the string exerts on each body attached to it, directed \emph{along the string, away from the body} (a string can only pull). If the string is \textbf{light} and passes over a \textbf{smooth} pulley, then $T$ has the same value throughout the string.
\end{definition}

Why does a string always pull and never push? Because a flexible string cannot sustain compression: push it and it simply buckles and goes \emph{slack}. When a string goes slack, $T = 0$ and the particles cease to be connected --- this is a classic examination trap.

\begin{propriete}[Newton's third law and tension]
The forces exerted by a taut string on the two bodies at its ends form a Newton third-law pair: they are \textbf{equal in magnitude and opposite in direction}. If the string pulls body $A$ with force $T$ towards the string, it pulls body $B$ with force $T$ towards the string at the other end. (Proof not examinable; it follows directly from the light-string assumption and $F=ma$ applied to the massless string.)
\end{propriete}

\begin{popfigure}
\begin{center}
\begin{tikzpicture}[scale=1.0, >=stealth]
% Pulley
\draw[thick, popink] (3.0,1.5) circle (0.4);
\draw[fill=popblueL] (3.0,1.5) circle (0.1);
% Left particle A
\draw[fill=popblueL, draw=popink] (1.0,0.0) rectangle (2.4,1.0);
\node at (1.7,0.5) {$A$};
\draw[->, very thick, poporange] (2.4,0.5) -- (2.8,0.5) node[midway, above]{$T$};
\draw[->, very thick, popink] (1.7,0) -- (1.7,-1.1) node[midway, left]{$m_A g$};
% String left side
\draw[thick, popdark] (2.4,0.5) -- (2.8,0.5) -- (3.0,1.1);
% String right side
\draw[thick, popdark] (3.0,1.9) -- (3.2,1.5) -- (4.6,0.5);
% Right particle B
\draw[fill=popgreenL, draw=popink] (4.6,0.0) rectangle (6.0,1.0);
\node at (5.3,0.5) {$B$};
\draw[->, very thick, poporange] (4.6,0.5) -- (4.2,0.5) node[midway, above]{$T$};
\draw[->, very thick, popink] (5.3,0) -- (5.3,-1.1) node[midway, right]{$m_B g$};
\end{tikzpicture}
\end{center}
\legende{Free-body diagrams of two particles connected by a light inextensible string over a smooth pulley. The tension $T$ has the same magnitude on both sides.}
\end{popfigure}

\textbf{Strings versus rods and tow-bars.}\par
A string, rope or chain can only \emph{pull}: the force in it is always a tension ($T \geq 0$), and it may go slack. A rigid \cle{rod} or \cle{tow-bar} can also \emph{push}: the force in it can be a tension (pulling) or a \cle{thrust} (pushing, sometimes written $T<0$ or stated separately as a thrust). When a car tows a trailer with a rigid tow-bar and then brakes hard, the bar may be in thrust, pushing the trailer back. Always read the question to see which connector is used.

\courssub{The Five-Step Routine}

Every connected-particle problem in this chapter is solved by the same disciplined routine.

\begin{methode}[Solving a connected-particle problem]
\begin{enumerate}
\item \textbf{State the assumptions:} particles, light inextensible string (or light rigid rod), smooth pulley (or smooth/rough surface, as given).
\item \textbf{Draw a separate free-body diagram for each particle}, showing only the forces acting \emph{on that particle}: weight, normal reaction, tension(s), applied forces, friction.
\item \textbf{Choose the direction of motion} for each particle (consistent with the string) and mark the common acceleration $a$.
\item \textbf{Write Newton's second law} $F = ma$ for each particle, resolved along its own direction of motion. (Alternatively, write the equation for the whole system to find $a$ directly, then use one particle's equation to find $T$.)
\item \textbf{Solve the simultaneous equations} for $a$ and $T$ (and check the sign of $T$: for a string, $T \geq 0$; if $T<0$, the string is slack and the model breaks down).
\end{enumerate}
\end{methode}

\begin{attention}[One diagram per particle!]
The most frequent mistake at GCE is to draw both particles on a single diagram and to mix \emph{internal} forces (the tensions) with \emph{external} forces (weights, driving forces, resistances). For the \emph{whole system} the internal tensions cancel out, but for \emph{each particle separately} the tension is essential. Always draw two distinct diagrams and write two distinct equations.
\end{attention}

\courssub{First Example: Two Particles Over a Smooth Fixed Pulley}

Consider two particles of masses $m_1$ and $m_2$ ($m_1 > m_2$) tied to the ends of a light inextensible string passing over a smooth fixed pulley --- the mechanic's situation in Limbe, idealised. The heavier mass $m_1$ descends, the lighter mass $m_2$ rises, both with acceleration of magnitude $a$. Applying the routine:

\begin{itemize}
\item Particle 1 (moving down, take down as positive): $m_1 g - T = m_1 a$.
\item Particle 2 (moving up, take up as positive): $T - m_2 g = m_2 a$.
\end{itemize}

Adding eliminates $T$:
\[
a = \frac{(m_1 - m_2)g}{m_1 + m_2}, \qquad
T = \frac{2 m_1 m_2 g}{m_1 + m_2}.
\]
Note that $a < g$: this is why the engine does not crash to the ground --- the counterweight, through the tension, holds it back. The full quantitative study of this \emph{Atwood machine} comes in Section 4.

\courssub{Worked Example: A Car Towing a Trailer}

\begin{popfigure}
\begin{center}
\begin{tikzpicture}[scale=1.0, >=stealth]
% road
\draw[thick, popink] (-0.5,0) -- (10.5,0);
% car
\draw[fill=popblueL, draw=popink] (0,0) rectangle (2.6,1.1);
\node at (1.3,0.55) {Car};
\draw[fill=white, draw=popink] (0.5,0) circle (0.28);
\draw[fill=white, draw=popink] (2.1,0) circle (0.28);
% tow bar
\draw[very thick, popdark] (2.6,0.45) -- (4.2,0.45);
\node at (3.4,0.8) {\small tow-bar};
% trailer
\draw[fill=popgreenL, draw=popink] (4.2,0) rectangle (6.6,1.1);
\node at (5.4,0.55) {Trailer};
\draw[fill=white, draw=popink] (4.7,0) circle (0.28);
\draw[fill=white, draw=popink] (6.1,0) circle (0.28);
% forces on car
\draw[->, very thick, popgreen] (1.3,1.1) -- (3.3,1.1) node[midway, above]{$D$};
\draw[->, very thick, poppink] (1.3,-0.35) -- (0.1,-0.35) node[midway, below]{$R_1$};
% forces on trailer
\draw[->, very thick, poppink] (5.4,-0.35) -- (4.4,-0.35) node[midway, below]{$R_2$};
% acceleration
\draw[->, very thick, poporange] (8.0,0.55) -- (9.6,0.55) node[midway, above]{$a$};
% tension arrows on each body
\draw[->, very thick, poppurple] (2.6,1.35) -- (1.8,1.35) node[midway, above]{$T$};
\draw[->, very thick, poppurple] (4.2,1.35) -- (5.0,1.35) node[midway, above]{$T$};
\end{tikzpicture}
\end{center}
\legende{Car towing a trailer: driving force $D$ and resistance $R_1$ on the car, resistance $R_2$ on the trailer, and the tow-bar force $T$ pulling the trailer forward and the car backward.}
\end{popfigure}

\begin{exoresolu}[Car towing a trailer on a horizontal road]
A car of mass $1200\ \mathrm{kg}$ tows a trailer of mass $400\ \mathrm{kg}$ along a straight horizontal road by means of a rigid tow-bar. The car's engine produces a driving force of $2400\ \mathrm{N}$. Resistance forces of $600\ \mathrm{N}$ act on the car and $200\ \mathrm{N}$ on the trailer. Find (a) the acceleration of the system, (b) the force in the tow-bar, stating whether it is a tension or a thrust.
\tcblower
\textbf{Solution.} \emph{Step 1 -- Assumptions:} car and trailer are particles; the tow-bar is light and rigid; the road is horizontal.

\emph{Step 2 -- Diagrams:} separate diagrams for car and trailer (see figure above).

\emph{Step 3 -- Direction of motion:} forwards, common acceleration $a$.

\emph{Step 4 -- Newton's second law.}

For the trailer ($400\ \mathrm{kg}$), the only horizontal forces are the tow-bar pull $T$ and the resistance:
\[
T - 200 = 400a. \tag{1}
\]
For the car ($1200\ \mathrm{kg}$), the driving force $D=2400$ N acts forwards, the resistance $600$ N and the tow-bar reaction $T$ act backwards:
\[
2400 - 600 - T = 1200a. \tag{2}
\]

\emph{Step 5 -- Solve.} Adding (1) and (2) gives the system equation:
\[
2400 - 600 - 200 = 1600a \;\Longrightarrow\; a = \frac{1600}{1600} = 1\ \mathrm{m\,s^{-2}}.
\]
From (1): $T = 200 + 400(1) = 600\ \mathrm{N}$.

\textbf{Answers:} (a) $a = 1\ \mathrm{m\,s^{-2}}$; (b) the tow-bar is in \textbf{tension} of $600\ \mathrm{N}$ (it pulls the trailer forward and drags the car backward). Since $T>0$, the bar is pulling, not pushing.
\end{exoresolu}

\begin{exercice}[A pickup and a taxi]
A pickup of mass $1500\ \mathrm{kg}$ tows a broken-down taxi of mass $900\ \mathrm{kg}$ along a horizontal road with a light chain. The driving force is $3600\ \mathrm{N}$; resistances are $750\ \mathrm{N}$ on the pickup and $450\ \mathrm{N}$ on the taxi. Find the acceleration and the tension in the chain. What happens to the tension if the pickup brakes gently?
\end{exercice}

\begin{corrige}[Exercise 1]
Newton's second law for the taxi: $T - 450 = 900a$. For the pickup: $3600 - 750 - T = 1500a$. Adding: $2400 = 2400a$, so $a = 1\ \mathrm{m\,s^{-2}}$ and $T = 450 + 900 = 1350\ \mathrm{N}$. If the pickup brakes, the taxi pushes on the pickup; a chain cannot push, so it goes \emph{slack} ($T = 0$) and the taxi runs free until the chain tightens again.
\end{corrige}

\begin{savaistu}[Why tow-bars instead of ropes?]
Road-safety rules in many countries require a rigid tow-bar rather than a rope for towing on fast roads: with a rope, the moment the towing vehicle slows down, the towed vehicle keeps going, the rope goes slack, and a collision follows. A rigid bar transmits thrust as well as tension, so the towing vehicle can help brake the towed one --- exactly the string-versus-rod distinction of this section.
\end{savaistu}

\begin{aretenir}[Key points of Section 1]
\begin{itemize}
\item Connected particles move with the \textbf{same speed} and \textbf{same magnitude of acceleration} when linked by a taut inextensible string.
\item A \textbf{light} string over a \textbf{smooth} pulley has the \textbf{same tension throughout}; the pulley changes only the direction of the force.
\item A string \textbf{pulls only} and may go slack ($T=0$); a rod or tow-bar can also \textbf{push} (thrust).
\item Always draw \textbf{one free-body diagram per particle} and write $F=ma$ for each, then solve simultaneously.
\item Tension forces on the two bodies form a Newton third-law pair: equal magnitude, opposite direction.
\item Check that $T \ge 0$ for a string; if $T<0$ the string is slack and the particles move independently.
\end{itemize}
\end{aretenir}

\coursec{Section 2: Motion on Smooth and Rough Horizontal Surfaces}

In Section 1 we learnt to model a system of connected particles: light inextensible strings, smooth pulleys, and the key fact that the \cle{tension} is the same throughout a single string while both particles share the same \cle{acceleration} magnitude. We now put these ideas to work in the most classic configuration of the GCE Advanced Level examination: one particle resting on a \cle{horizontal table}, joined by a string passing over a pulley at the edge to a second particle hanging freely. We first treat the smooth table, then introduce friction, and finish with towing problems and the use of the constant-acceleration (\emph{suvat}) formulae.

\courssub{The Smooth Table with a Hanging Particle}

\begin{experience}[Setting up the situation]
Place a block of mass $m$ on a polished table. Tie a light string to it, pass the string over a small pulley fixed at the edge of the table, and attach a mass $M$ hanging freely. Release the system. The hanging mass falls, the block slides towards the edge, and both move together with the same acceleration because the string is inextensible. Our task is to predict that acceleration and the tension in the string.
\end{experience}

\begin{popfigure}
\begin{center}
\begin{tikzpicture}[scale=1.0, >=stealth]
% table
\draw[fill=popblueL, draw=popdark] (-0.2,0) rectangle (5.4,-0.25);
\draw[popdark] (0.3,-0.25) -- (0.3,-1.4);
\draw[popdark] (4.9,-0.25) -- (4.9,-1.4);
% block on table
\draw[fill=popgreenL, draw=popdark] (2.0,0) rectangle (3.2,0.9);
\node at (2.6,0.45) {$m$};
% pulley
\draw[fill=white, draw=popdark] (5.4,0.15) circle (0.18);
\fill[popdark] (5.4,0.15) circle (0.03);
% string
\draw[poporange, thick] (3.2,0.45) -- (5.4,0.45) arc (90:0:0.18) -- (5.58,-1.5);
% hanging mass
\draw[fill=poppink, draw=popdark] (5.28,-1.5) rectangle (5.88,-2.3);
\node at (5.58,-1.9) {$M$};
% forces on m
\draw[->, popblue, thick] (2.6,0.45) -- (4.0,0.45) node[midway, above]{$T$};
\draw[->, popblue, thick] (2.6,0.45) -- (2.6,-0.75) node[midway, left]{$mg$};
\draw[->, popblue, thick] (2.6,0.45) -- (2.6,1.55) node[midway, left]{$R$};
% forces on M
\draw[->, poppurple, thick] (5.58,-1.9) -- (5.58,-3.0) node[midway, right]{$Mg$};
\draw[->, poppurple, thick] (5.58,-1.9) -- (5.58,-0.85) node[midway, right]{$T$};
% acceleration arrows
\draw[->, popgreen, very thick] (1.2,0.45) -- (1.8,0.45) node[midway, above]{$a$};
\draw[->, popgreen, very thick] (6.25,-1.2) -- (6.25,-1.9) node[midway, right]{$a$};
\end{tikzpicture}
\end{center}
\legende{Block of mass $m$ on a smooth table, connected over a smooth pulley at the edge to a hanging mass $M$. Forces on each particle are shown.}
\end{popfigure}

\textbf{Analysing each particle separately.}\par
For the block on the table there is no vertical motion, so the vertical forces balance: $R = mg$. Horizontally the only force is the tension, so Newton's second law gives
\[ T = ma. \tag{1} \]
For the hanging mass, taking the downward direction (the direction of its motion) as positive:
\[ Mg - T = Ma. \tag{2} \]
Adding (1) and (2) eliminates $T$:
\[ Mg = (M + m)a \quad\Longrightarrow\quad a = \frac{Mg}{M+m}, \qquad T = ma = \frac{Mmg}{M+m}. \]

\begin{methode}[Smooth table: direct formulae]
For a hanging mass $M$ driving a mass $m$ on a \textbf{smooth} horizontal table:
\begin{enumerate}
\item Acceleration: $a = \dfrac{Mg}{M+m}$ — driving weight divided by total mass.
\item Tension: $T = \dfrac{Mmg}{M+m}$, obtained by substituting $a$ into $T = ma$.
\item Always check: $T < Mg$ (the hanging mass is accelerating downwards, so its weight must exceed the tension).
\end{enumerate}
This \emph{whole-system shortcut} is valid whenever the string stays taut, the pulley is smooth, and no other horizontal force (such as friction) acts on the table particle.
\end{methode}

\begin{exoresolu}[Smooth table: finding $a$ and $T$]
A particle $P$ of mass $2\ \mathrm{kg}$ rests on a smooth horizontal table. It is connected by a light inextensible string passing over a smooth pulley at the edge of the table to a particle $Q$ of mass $3\ \mathrm{kg}$ hanging freely. The system is released from rest. Find the acceleration of the system and the tension in the string. Take $g = 9.8\ \mathrm{m\,s^{-2}}$.
\tcblower
\textbf{Solution.} For $P$ (horizontal): $T = 2a$. For $Q$ (downwards positive): $3g - T = 3a$. Adding:
\[ 3g = 5a \implies a = \frac{3 \times 9.8}{5} = 5.88\ \mathrm{m\,s^{-2}}. \]
Then $T = 2a = 11.76\ \mathrm{N}$. Check: $T = 11.76 < 3g = 29.4$ \checkmark, and $3g - T = 17.64 = 3 \times 5.88$ \checkmark.
\end{exoresolu}

\begin{attention}[The shortcut gives only $a$]
A very common error is to write $a = \dfrac{Mg}{M+m}$ and then claim $T = Mg$ or $T = Ma$. The whole-system equation contains no tension at all — the internal forces cancel. You \textbf{must} return to the equation of one particle ($T = ma$ or $Mg - T = Ma$) to find $T$. Notice also that $T = \dfrac{Mmg}{M+m}$ is \emph{less} than both $Mg$ and (usually) what candidates guess.
\end{attention}

\courssub{Introducing Friction: the Rough Table}

A real table — a wooden desk in a Yaoundé classroom, for instance — is not smooth. When the block slides (or tends to slide), the surface exerts a frictional force opposing the motion.

\begin{propriete}[Normal reaction and the friction model]
For a particle on a \textbf{horizontal} surface with no vertical acceleration, resolving vertically gives
\[ R = mg. \]
The frictional force satisfies $F \le \mu R$, where $\mu$ is the coefficient of friction. When the body is \textbf{moving} (or in \emph{limiting equilibrium}, on the point of moving),
\[ F = \mu R, \]
directed \textbf{opposite to the motion} (or to the tendency of motion). If the body is at rest and not on the point of moving, $F$ takes whatever value below $\mu R$ is needed to maintain equilibrium.
\end{propriete}

\begin{methode}[Test first: will the system move at all?]
Before writing any equation of motion on a rough table:
\begin{enumerate}
\item Compute the maximum available friction $F_{\max} = \mu R = \mu mg$.
\item Compute the driving force, here the hanging weight $Mg$.
\item If $Mg \le \mu mg$, the system \textbf{does not move}: friction adjusts to $F = Mg$ and the string tension is $T = Mg$ with $a = 0$.
\item If $Mg > \mu mg$, motion occurs; write the equations with $F = \mu R$ opposing the motion.
\end{enumerate}
\end{methode}

\begin{attention}[Do not write $F = \mu R$ blindly]
Two mistakes are frequent. First, writing $F = \mu R$ for a \emph{smooth} surface (where $F = 0$). Second, writing $F = \mu R$ when the body is at rest and the driving force is small: then $F < \mu R$, and assuming $F = \mu R$ can produce a \emph{negative} acceleration — a contradiction telling you that your assumption of motion was wrong. The direction of $F$ also matters: it always opposes the actual or impending motion, so it points \emph{away} from the pulley here.
\end{attention}

\begin{popfigure}
\begin{center}
\begin{tikzpicture}[scale=1.0, >=stealth]
% table (hatched to suggest roughness)
\draw[fill=popblueL, draw=popdark] (-0.2,0) rectangle (5.4,-0.25);
\foreach \x in {0.2,0.7,1.2,1.7,2.2,2.7,3.2,3.7,4.2,4.7,5.2}{\draw[popmuted] (\x,0) -- (\x+0.18,-0.25);}
\draw[popdark] (0.3,-0.25) -- (0.3,-1.4);
\draw[popdark] (4.9,-0.25) -- (4.9,-1.4);
% block
\draw[fill=popgreenL, draw=popdark] (2.0,0) rectangle (3.2,0.9);
\node at (2.6,0.45) {$m$};
% pulley and string
\draw[fill=white, draw=popdark] (5.4,0.15) circle (0.18);
\draw[poporange, thick] (3.2,0.45) -- (5.4,0.45) arc (90:0:0.18) -- (5.58,-1.4);
\draw[fill=poppink, draw=popdark] (5.28,-1.4) rectangle (5.88,-2.2);
\node at (5.58,-1.8) {$M$};
% forces on m
\draw[->, popblue, thick] (2.6,0.45) -- (4.1,0.45) node[midway, above]{$T$};
\draw[->, poppink, very thick] (2.6,0.2) -- (1.2,0.2) node[midway, above]{$F=\mu R$};
\draw[->, popblue, thick] (2.6,0.45) -- (2.6,-0.85) node[midway, left]{$mg$};
\draw[->, popblue, thick] (2.6,0.45) -- (2.6,1.65) node[midway, left]{$R=mg$};
% forces on M
\draw[->, poppurple, thick] (5.58,-1.8) -- (5.58,-2.9) node[midway, right]{$Mg$};
\draw[->, poppurple, thick] (5.58,-1.8) -- (5.58,-0.75) node[midway, right]{$T$};
\draw[->, popgreen, very thick] (6.25,-1.1) -- (6.25,-1.8) node[midway, right]{$a$};
\end{tikzpicture}
\end{center}
\legende{Rough table: friction $F = \mu R$ opposes the motion of the block towards the pulley, while vertically $R = mg$.}
\end{popfigure}

\textbf{Equations of motion on the rough table.}\par
Assuming motion occurs (block moving towards the pulley):
\[ \text{Block: } T - F = ma, \quad F = \mu R = \mu mg; \qquad \text{Hanging mass: } Mg - T = Ma. \]
Adding eliminates $T$:
\[ Mg - \mu mg = (M+m)a \quad\Longrightarrow\quad a = \frac{(M - \mu m)g}{M+m}, \qquad T = M(g - a) = \frac{Mm(1+\mu)g}{M+m}. \]

\begin{exoresolu}[Rough table: full analysis]
A block of mass $4\ \mathrm{kg}$ lies on a rough horizontal table with coefficient of friction $\mu = 0.25$. It is joined by a light inextensible string over a smooth edge pulley to a hanging mass of $6\ \mathrm{kg}$. The system is released from rest. Take $g = 9.8\ \mathrm{m\,s^{-2}}$.
\begin{enumerate}
\item Show that motion occurs.
\item Find the acceleration and the tension.
\item Find the speed of the hanging mass after it has fallen $0.8\ \mathrm{m}$.
\end{enumerate}
\tcblower
\textbf{Solution.}
\textbf{(1)} Maximum friction: $F_{\max} = \mu mg = 0.25 \times 4 \times 9.8 = 9.8\ \mathrm{N}$. Driving weight: $Mg = 6 \times 9.8 = 58.8\ \mathrm{N}$. Since $58.8 > 9.8$, motion occurs.
\textbf{(2)} Block: $T - 9.8 = 4a$. Hanging mass: $58.8 - T = 6a$. Adding: $49 = 10a$, so $a = 4.9\ \mathrm{m\,s^{-2}} > 0$ — consistent with our assumption of motion \checkmark. Then $T = 9.8 + 4(4.9) = 29.4\ \mathrm{N}$.
\textbf{(3)} Using $v^2 = u^2 + 2as$ with $u = 0$, $a = 4.9$, $s = 0.8$:
\[ v^2 = 2 \times 4.9 \times 0.8 = 7.84 \implies v = 2.8\ \mathrm{m\,s^{-1}}. \]
\end{exoresolu}

\begin{exoresolu}[When the system does not move]
The block of the previous example is replaced by one of mass $10\ \mathrm{kg}$, with $\mu = 0.7$, the hanging mass remaining $6\ \mathrm{kg}$. Describe the motion.
\tcblower
\textbf{Solution.} $F_{\max} = \mu mg = 0.7 \times 10 \times 9.8 = 68.6\ \mathrm{N}$, while the driving weight is only $Mg = 58.8\ \mathrm{N}$. Since $58.8 < 68.6$, friction can balance the pull: the system \textbf{remains at rest}. The actual friction is $F = 58.8\ \mathrm{N} < \mu R$, and the tension is $T = Mg = 58.8\ \mathrm{N}$ (the hanging mass is in equilibrium). Writing $a = \dfrac{(M-\mu m)g}{M+m}$ would give a negative value — the signal that the motion assumption fails.
\end{exoresolu}

\courssub{Using the suvat Equations}

Once the (constant) acceleration is known, all the constant-acceleration formulae apply to \emph{each} particle, since both share $a$:
\[ v = u + at, \qquad s = ut + \tfrac{1}{2}at^2, \qquad v^2 = u^2 + 2as, \qquad s = \tfrac{1}{2}(u+v)t. \]
Typical examination questions ask for the time taken by the table particle to reach the pulley, or the speed of the hanging mass when it hits the floor. Remember that the two particles describe the \emph{same} distance along the string while it remains taut.

\begin{exoresolu}[Time to reach the edge]
In the first worked example (smooth table, $m = 2$, $M = 3$, $a = 5.88\ \mathrm{m\,s^{-2}}$), the block starts $1.5\ \mathrm{m}$ from the pulley. Find the time taken to reach the pulley and its speed on arrival.
\tcblower
\textbf{Solution.} From $s = ut + \tfrac{1}{2}at^2$ with $u = 0$:
\[ 1.5 = \tfrac{1}{2}(5.88)t^2 \implies t^2 = \frac{3}{5.88} = 0.5102 \implies t \approx 0.714\ \mathrm{s}. \]
Speed: $v = at = 5.88 \times 0.714 \approx 4.2\ \mathrm{m\,s^{-1}}$ (check with $v^2 = 2as = 17.64$, $v = 4.2$ \checkmark).
\end{exoresolu}

\courssub{Extension: Towing Problems with Resistance}

The same logic governs a \cle{car towing a trailer}, with the tow-bar playing the role of the string. The difference is that a tow-bar can \emph{push} as well as pull: when the car brakes, the bar is in \cle{compression} and exerts a \cle{thrust}. Resistance forces (air resistance, rolling resistance) act on each vehicle, constant or proportional to speed; at GCE level they are usually given as constants.

\begin{methode}[Car and trailer]
\begin{enumerate}
\item Draw separate force diagrams for the car and the trailer; mark driving force, resistances, and the force $T$ in the tow-bar. Assume $T$ is a \textbf{tension} (pull) — i.e. the bar pulls the trailer forward and the car backward.
\item Write $F = ma$ for each vehicle along the direction of motion (take the direction of motion as positive).
\item Add the equations to get $a$ for the whole system, then back-substitute into one equation to find $T$.
\item If $T$ comes out \emph{negative}, the bar is in \textbf{thrust} (compression) of magnitude $|T|$ — typical under braking.
\end{enumerate}
\end{methode}

\begin{exoresolu}[Braking car and trailer: thrust in the tow-bar]
A car of mass $1200\ \mathrm{kg}$ tows a trailer of mass $400\ \mathrm{kg}$ along a horizontal road. The car brakes, producing a braking force of $6000\ \mathrm{N}$; resistances of $200\ \mathrm{N}$ on the car and $100\ \mathrm{N}$ on the trailer also act. Find the deceleration and the force in the tow-bar, stating whether it is a tension or a thrust.
\tcblower
\textbf{Solution.} Take the direction of motion as positive. All resistive forces are negative.
Whole system:
\[ -(6000 + 200 + 100) = (1200 + 400)a \implies -6300 = 1600a \implies a = -\frac{6300}{1600} = -3.9375\ \mathrm{m\,s^{-2}}. \]
The deceleration is $3.9375\ \mathrm{m\,s^{-2}}$.
Now consider the trailer alone. Assume the tow-bar is in tension, so it pulls the trailer forward with force $T$ (positive). The resistance on the trailer is $100\ \mathrm{N}$ backward.
\[ T - 100 = 400 \times (-3.9375) = -1575 \implies T = -1475\ \mathrm{N}. \]
Since $T$ is negative, our assumption of tension is wrong; the bar is in \textbf{compression} (thrust) of magnitude $1475\ \mathrm{N}$. The bar pushes the trailer backward (and the car forward) to slow the trailer down.
\end{exoresolu}

\begin{attention}[Sign conventions in towing]
Decide at the start which way is positive and assume the tow-bar force is a tension (pull). Keep that convention for \emph{both} vehicles: the bar exerts equal and opposite forces on car and trailer (Newton's third law). A negative result for $T$ means the bar is in thrust. Mixing conventions is the most common source of sign errors in braking questions.
\end{attention}

\begin{savaistu}[Friction in daily life]
The friction model $F = \mu R$ is an approximation (Amontons–Coulomb law). Real friction depends on contact area, temperature, and speed. In Cameroon, think of a taxi braking on a wet road in Douala: $\mu$ drops dramatically, increasing stopping distance. Engineers design brake systems to maximise $\mu$ while avoiding wheel lock (which switches from static to kinetic friction, reducing $\mu$). This is why anti-lock braking systems (ABS) are vital for safety.
\end{savaistu}

\begin{aretenir}[Key points of Section 2]
\begin{itemize}
\item Smooth table with hanging mass $M$ driving mass $m$: $a = \dfrac{Mg}{M+m}$, $T = \dfrac{Mmg}{M+m}$.
\item Rough table: first test $Mg$ against $\mu mg$. Motion occurs only if $Mg > \mu mg$; then $a = \dfrac{(M-\mu m)g}{M+m}$ and $F = \mu R$ opposes the motion.
\item On a horizontal surface with no vertical acceleration, $R = mg$.
\item The whole-system equation yields $a$ only; tension requires a separate single-particle equation.
\item With constant $a$ found, use the suvat formulae for times, speeds and distances.
\item In towing problems, assume tension in the tow-bar; a negative result signals a thrust (compression), typical under braking.
\end{itemize}
\end{aretenir}

\begin{exercice}[Practice]
\begin{enumerate}
\item A mass of $5\ \mathrm{kg}$ on a smooth table is connected over a smooth edge pulley to a hanging mass of $3\ \mathrm{kg}$. Find the acceleration, the tension, and the speed of the system after $0.5\ \mathrm{s}$ from rest. ($g = 9.8\ \mathrm{m\,s^{-2}}$)
\item A block of mass $8\ \mathrm{kg}$ rests on a rough horizontal table, $\mu = 0.3$, connected to a hanging mass of $2\ \mathrm{kg}$. Determine whether the system moves; if it does, find $a$ and $T$; if not, find the friction and the tension.
\item Repeat question 2 with a hanging mass of $5\ \mathrm{kg}$, and find how long the block takes to travel $1.2\ \mathrm{m}$ from rest.
\item A car of mass $1000\ \mathrm{kg}$ towing a trailer of mass $250\ \mathrm{kg}$ accelerates with driving force $3200\ \mathrm{N}$ against total resistances of $450\ \mathrm{N}$ on the car and $150\ \mathrm{N}$ on the trailer. Find the acceleration and the tension in the tow-bar.
\end{enumerate}
\end{exercice}

\begin{corrige}[Exercise 1]
$a = \dfrac{3g}{8} = 3.675\ \mathrm{m\,s^{-2}}$; $T = 5a = 18.375\ \mathrm{N}$; $v = at = 3.675 \times 0.5 \approx 1.84\ \mathrm{m\,s^{-1}}$.
\end{corrige}

\begin{corrige}[Exercise 2]
$F_{\max} = 0.3 \times 8 \times 9.8 = 23.52\ \mathrm{N}$; driving weight $= 2g = 19.6\ \mathrm{N} < 23.52\ \mathrm{N}$. No motion: $a = 0$, $F = 19.6\ \mathrm{N}$, $T = 19.6\ \mathrm{N}$.
\end{corrige}

\begin{corrige}[Exercise 3]
$5g = 49 > 23.52$, so motion occurs. $a = \dfrac{(5 - 0.3 \times 8) \times 9.8}{13} = \dfrac{2.6 \times 9.8}{13} = 1.96\ \mathrm{m\,s^{-2}}$; $T = 5(g - a) = 5 \times 7.84 = 39.2\ \mathrm{N}$. From $s = \tfrac{1}{2}at^2$: $t = \sqrt{\dfrac{2 \times 1.2}{1.96}} = \sqrt{1.2245} \approx 1.11\ \mathrm{s}$.
\end{corrige}

\begin{corrige}[Exercise 4]
Whole system: $3200 - 450 - 150 = 1250a$, so $a = \dfrac{2600}{1250} = 2.08\ \mathrm{m\,s^{-2}}$. Trailer: $T - 150 = 250 \times 2.08 = 520$, hence $T = 670\ \mathrm{N}$ (tension, as expected under acceleration).
\end{corrige}

\coursec{Section 3: Motion on Inclined Planes}

In Section 2 we studied connected particles on \emph{horizontal} surfaces, where the weight of a block was completely balanced by the normal reaction ($R = mg$) and only the tension and friction acted along the direction of motion. Life in Cameroon, however, is rarely horizontal: think of a truck climbing the steep road to Dschang, a wheelbarrow pushed up a plank onto a lorry in a Douala warehouse, or a crate sliding down a ramp. On a slope, the weight itself has a component \emph{along} the surface, and this changes everything: the normal reaction, the friction, and the driving force of the motion. This section teaches you to handle such problems with complete confidence.

\courssub{Resolving the Weight on an Inclined Plane}

\textbf{Observation.} Place a book on a flat table and tilt the table slowly. At first nothing happens; then, beyond a certain angle, the book slides. Gravity has not changed — what has changed is the \emph{direction} in which gravity can push the book relative to the surface. To understand the motion we must split the weight into two perpendicular parts: one pulling the object \emph{along} the plane, and one pressing it \emph{into} the plane.

Consider a particle of mass $m$ on a plane inclined at an angle $\alpha$ to the horizontal. The weight $mg$ acts vertically downwards. We resolve it into:
\begin{itemize}
\item a component \textbf{parallel to the plane}, pointing down the slope;
\item a component \textbf{perpendicular to the plane}, pressing into the surface.
\end{itemize}

\begin{popfigure}
\begin{center}
\begin{tikzpicture}[scale=1.05, >=stealth]
% inclined plane
\coordinate (O) at (0,0);
\coordinate (T) at (7,3.06);
\draw[thick, popdark] (O) -- (T);
\draw[thick, popdark] (O) -- (7.6,0);
\draw[popmuted] (1.2,0) arc (0:23.6:1.2);
\node[popmuted] at (1.7,0.28) {$\alpha$};
% block
\begin{scope}[rotate around={23.6:(3.5,1.53)}]
\draw[fill=popblueL, draw=popblue, thick] (3.1,1.53) rectangle (3.9,2.13);
\end{scope}
\coordinate (C) at (3.5,2.05);
% weight
\draw[->, very thick, poppink] (C) -- (3.5,0.35) node[below] {$mg$};
% components
\draw[->, very thick, poporange] (C) -- (2.05,1.42) node[below left, poporange] {$mg\sin\alpha$};
\draw[->, very thick, popgreen] (C) -- (4.22,0.72) node[below right, popgreen] {$mg\cos\alpha$};
% dashed construction
\draw[dashed, popmuted] (3.5,0.35) -- (2.05,1.42);
\draw[dashed, popmuted] (3.5,0.35) -- (4.22,0.72);
% normal reaction
\draw[->, very thick, poppurple] (C) -- (2.78,3.38) node[above left, poppurple] {$R$};
\end{tikzpicture}
\end{center}
\legende{Resolving the weight on a plane inclined at $\alpha$: the component $mg\sin\alpha$ acts down the plane, $mg\cos\alpha$ presses into the plane, and the normal reaction $R$ acts perpendicular to the surface.}
\end{popfigure}

The geometry is the key point. The angle between the vertical (direction of $mg$) and the normal to the plane is exactly $\alpha$, because the normal is tilted by $\alpha$ from the vertical just as the plane is tilted by $\alpha$ from the horizontal. Hence, in the right-angled triangle formed by the weight and its two components:
\[
\text{component down the plane} = mg\sin\alpha, \qquad \text{component into the plane} = mg\cos\alpha.
\]

\begin{propriete}[Weight components and normal reaction on an inclined plane]
For a particle of mass $m$ resting on a plane inclined at angle $\alpha$ to the horizontal:
\begin{itemize}
\item the component of weight \textbf{down the plane} is $mg\sin\alpha$;
\item the component of weight \textbf{perpendicular into the plane} is $mg\cos\alpha$;
\item since there is no acceleration perpendicular to the plane (the particle neither sinks into nor lifts off the surface), resolving perpendicular to the plane gives
\[
R = mg\cos\alpha.
\]
\end{itemize}
This result is examinable and must be justified by resolving perpendicular to the plane.
\end{propriete}

\begin{methode}[Resolving parallel and perpendicular to the plane]
\begin{enumerate}
\item Draw a free-body diagram: weight $mg$ vertically down, normal reaction $R$ perpendicular to the plane, friction (if any) along the plane opposing motion, tension (if any) along the string.
\item Choose axes \textbf{parallel} and \textbf{perpendicular} to the plane.
\item Replace $mg$ by its two components $mg\sin\alpha$ (down the plane) and $mg\cos\alpha$ (into the plane).
\item Write $\sum F = ma$ \textbf{along} the plane, and $\sum F = 0$ \textbf{perpendicular} to the plane.
\item \textbf{Never mix} horizontal/vertical components with parallel/perpendicular components in the same equation. Choose one pair of axes and stay with it.
\end{enumerate}
\end{methode}

\begin{attention}[The most common mistake on inclined planes]
Many candidates write $R = mg$ out of habit from horizontal problems. On an inclined plane this is \textbf{wrong}: the correct reaction is
\[
R = mg\cos\alpha.
\]
Since friction is $F = \mu R$, using $R = mg$ also corrupts the friction force and destroys the whole solution. Always recompute $R$ by resolving perpendicular to the plane.
\end{attention}

\courssub{A Single Particle Sliding on an Inclined Plane}

\textbf{Smooth plane.} With no friction, the only force along the plane is the component $mg\sin\alpha$. Newton's second law down the plane gives
\[
mg\sin\alpha = ma \quad\Longrightarrow\quad \boxed{a = g\sin\alpha.}
\]
Notice that the acceleration is independent of the mass — just as in free fall, but reduced by the factor $\sin\alpha$.

\textbf{Rough plane, particle sliding down.} Friction now acts \emph{up} the plane (opposing the motion), with limiting value $F = \mu R = \mu mg\cos\alpha$. Newton's second law down the plane:
\[
mg\sin\alpha - \mu mg\cos\alpha = ma
\quad\Longrightarrow\quad
\[\boxed{a = g(\sin\alpha - \mu\cos\alpha).}\]
\]

\textbf{When does the particle stay at rest?} For the particle to remain in equilibrium, friction must be able to balance the driving component: we need $mg\sin\alpha \le \mu mg\cos\alpha$, i.e.
\[
\tan\alpha \le \mu.
\]

\begin{definition}[Angle of friction]
The \cle{angle of friction} $\lambda$ is defined by
\[
\tan\lambda = \mu.
\]
A particle placed on a rough plane remains at rest if and only if the inclination satisfies $\alpha \le \lambda$, equivalently $\tan\alpha \le \mu$. If the plane is tilted beyond $\lambda$, the particle must slide. (This is a useful GCE extra: it connects the equilibrium condition directly to the coefficient of friction.)
\end{definition}

\begin{exemplebox}[Checking equilibrium on a rough plank]
A crate of mass $25$ kg rests on a plank inclined at $20^{\circ}$ to the horizontal, with $\mu = 0.4$. Does it stay at rest? Take $g = 9.8\ \mathrm{m\,s^{-2}}$.

\emph{Solution.} $\tan 20^{\circ} \approx 0.364 < 0.4 = \mu$, so $\alpha < \lambda$ and the crate \textbf{remains at rest}. The friction needed is only $mg\sin 20^{\circ} \approx 25 \times 9.8 \times 0.342 \approx 83.8$ N, while the maximum available is $\mu mg\cos 20^{\circ} \approx 0.4 \times 25 \times 9.8 \times 0.940 \approx 92.1$ N. Friction is not limiting — the crate is safe.
\end{exemplebox}

\courssub{Particle on a Plane Connected to a Hanging Particle}

The classic GCE configuration: a block $A$ of mass $m_1$ rests on a plane inclined at $\alpha$, connected by a light inextensible string \emph{parallel to the plane}, passing over a smooth pulley fixed at the top, to a particle $B$ of mass $m_2$ hanging freely.

\begin{popfigure}
\begin{center}
\begin{tikzpicture}[scale=1.0, >=stealth]
% plane
\coordinate (O) at (0,0);
\coordinate (P) at (6.4,2.77);
\draw[thick, popdark] (O) -- (P);
\draw[thick, popdark] (O) -- (7,0);
\draw[popmuted] (1.1,0) arc (0:23.4:1.1);
\node[popmuted] at (1.55,0.26) {$\alpha$};
% pulley
\draw[fill=white, draw=popdark, thick] (6.4,3.15) circle (0.38);
\fill[popdark] (6.4,3.15) circle (0.05);
% block A on plane
\begin{scope}[rotate around={23.4:(3.2,1.39)}]
\draw[fill=popblueL, draw=popblue, thick] (2.7,1.39) rectangle (3.7,2.09);
\node at (3.2,1.74) {$m_1$};
\end{scope}
% string
\draw[thick, poporange] (3.62,1.98) -- (6.12,3.42);
\draw[thick, poporange] (6.4,3.53) -- (6.4,4.9);
% hanging mass
\draw[fill=popgreenL, draw=popgreen, thick] (6.05,4.9) rectangle (6.75,5.6);
\node at (6.4,5.25) {$m_2$};
% forces on m1
\draw[->, very thick, poppink] (3.2,1.95) -- (3.2,0.5) node[below] {$m_1 g$};
\draw[->, very thick, poppurple] (3.2,1.95) -- (2.5,3.1) node[above left] {$R$};
\draw[->, very thick, poporange] (3.6,2.15) -- (4.7,2.63) node[above] {$T$};
\draw[->, very thick, popdark] (2.85,1.75) -- (1.85,1.32) node[below left] {$F$};
% forces on m2
\draw[->, very thick, poporange] (6.4,5.6) -- (6.4,6.35) node[right] {$T$};
\draw[->, very thick, poppink] (6.4,4.9) -- (6.4,4.05) node[right] {$m_2 g$};
% acceleration hint
\draw[->, very thick, popblue] (7.15,5.3) -- (7.15,4.4) node[midway, right] {$a$};
\end{tikzpicture}
\end{center}
\legende{Block $m_1$ on a rough inclined plane connected over a smooth pulley to hanging mass $m_2$. Free-body elements: on the slope, $T$ up the plane, friction $F$ opposing motion, $R$ perpendicular, weight resolved; on the hanging mass, $T$ up and $m_2 g$ down.}
\end{popfigure}

\textbf{Setting up the equations.} Suppose the system moves with $m_2$ descending (so $m_1$ moves \emph{up} the plane) with common acceleration $a$, and the plane is rough with friction on $m_1$ acting \emph{down} the plane (opposing its motion). Then:

\begin{itemize}
\item Perpendicular to the plane (no acceleration): $R = m_1 g\cos\alpha$.
\item Along the plane, for $m_1$ (taking up the plane as positive):
\[
T - m_1 g\sin\alpha - \mu m_1 g\cos\alpha = m_1 a. \tag{1}
\]
\item Vertically, for the hanging mass $m_2$ (taking downwards as positive):
\[
m_2 g - T = m_2 a. \tag{2}
\]
\end{itemize}

Adding (1) and (2) eliminates $T$:
\[
a = \frac{m_2 g - m_1 g\sin\alpha - \mu m_1 g\cos\alpha}{m_1 + m_2},
\qquad
T = m_2(g - a).
\]

\begin{attention}[Assuming the direction of motion]
When the direction of motion is not obvious, \textbf{assume} one, write the equations consistently (friction opposing the assumed motion), and solve. If the computed $a$ is \textbf{negative}, the actual motion is in the opposite direction — but beware: you must then redo the problem with friction reversed, because friction always opposes the \emph{real} motion. If the driving force is smaller than the maximum available friction plus the opposing weight component, the system may simply \textbf{not move at all}; a negative acceleration under a limiting-friction assumption can mean equilibrium, not reversed motion.
\end{attention}

\begin{exoresolu}[Rough plane with hanging mass: finding $a$ and $T$]
A block $A$ of mass $4$ kg lies on a rough plane inclined at $30^{\circ}$ to the horizontal, with coefficient of friction $\mu = 0.2$. It is connected by a light inextensible string, parallel to the plane and passing over a smooth pulley at the top, to a particle $B$ of mass $6$ kg hanging freely. The system is released from rest. Take $g = 9.8\ \mathrm{m\,s^{-2}}$. Find the acceleration of the system and the tension in the string, and check consistency.
\tcblower
\textbf{Step 1 — Guess the direction.} Driving force if $B$ descends: $m_2 g = 6 \times 9.8 = 58.8$ N. Resisting: component of $A$'s weight down the plane, $m_1 g\sin 30^{\circ} = 4 \times 9.8 \times 0.5 = 19.6$ N, plus friction $\mu m_1 g\cos 30^{\circ} = 0.2 \times 4 \times 9.8 \times 0.866 \approx 6.79$ N. Total resistance $\approx 26.4$ N $< 58.8$ N, so $B$ does descend and $A$ moves up the plane.

\textbf{Step 2 — Normal reaction.} Perpendicular to the plane:
\[
R = m_1 g\cos 30^{\circ} = 4 \times 9.8 \times 0.866 \approx 33.95\ \mathrm{N}.
\]

\textbf{Step 3 — Equations of motion.} For $A$ up the plane, with friction $F = \mu R \approx 6.79$ N down the plane:
\[
T - 19.6 - 6.79 = 4a. \tag{1}
\]
For $B$ descending:
\[
58.8 - T = 6a. \tag{2}
\]

\textbf{Step 4 — Solve.} Adding (1) and (2):
\[
58.8 - 26.39 = 10a \quad\Longrightarrow\quad a = \frac{32.41}{10} \approx 3.24\ \mathrm{m\,s^{-2}}.
\]
From (2): $T = 58.8 - 6(3.24) \approx 58.8 - 19.45 \approx 39.4$ N.

\textbf{Step 5 — Consistency check.} $a > 0$, confirming our assumed direction. Also $T = 39.4$ N lies between $m_1 g\sin\alpha + F \approx 26.4$ N and $m_2 g = 58.8$ N, exactly as expected: the tension must both drag $A$ up and brake $B$'s fall. $\boxed{a \approx 3.24\ \mathrm{m\,s^{-2}},\ T \approx 39.4\ \mathrm{N}}$
\end{exoresolu}

\begin{exoresolu}[Smooth plane and suvat follow-up: speed at the bottom]
A particle of mass $2$ kg is released from rest at the top of a \emph{smooth} plane of length $5$ m inclined at $30^{\circ}$. Take $g = 9.8\ \mathrm{m\,s^{-2}}$. Find its acceleration and its speed at the bottom of the plane.
\tcblower
\textbf{Acceleration.} On a smooth plane, $a = g\sin\alpha = 9.8 \times \sin 30^{\circ} = 4.9\ \mathrm{m\,s^{-2}}$ down the plane.

\textbf{Speed at the bottom.} Using $v^2 = u^2 + 2as$ with $u = 0$, $a = 4.9$, $s = 5$:
\[
v^2 = 0 + 2 \times 4.9 \times 5 = 49 \quad\Longrightarrow\quad v = 7\ \mathrm{m\,s^{-1}}.
\]
Note the elegant check: $v^2 = 2g\sin\alpha \cdot s = 2g h$, where $h = s\sin\alpha$ is the vertical drop — exactly the free-fall result, foreshadowing the energy methods of Section 6.
\end{exoresolu}

\courssub{When the String Goes Slack: Motion Up the Plane}

A favourite GCE twist: the hanging mass hits the ground (or the string breaks) while the block on the plane is still moving \emph{up}. From that instant the tension vanishes, and the block decelerates under the component of its own weight (and friction, if rough) until it stops momentarily.

\begin{exoresolu}[Distance travelled up the plane after the string goes slack]
In the worked example above (rough plane, $\mu = 0.2$, $\alpha = 30^{\circ}$, $m_1 = 4$ kg), suppose the string breaks when block $A$ is moving up the plane at $3\ \mathrm{m\,s^{-1}}$. How much further does $A$ travel up the plane before coming to instantaneous rest? Take $g = 9.8\ \mathrm{m\,s^{-2}}$.
\tcblower
\textbf{New equation of motion.} With the string slack, the forces on $A$ along the plane are the weight component $m_1 g\sin\alpha$ and friction $\mu m_1 g\cos\alpha$, \emph{both} acting down the plane (friction opposes the upward motion):
\[
-m_1 g\sin\alpha - \mu m_1 g\cos\alpha = m_1 a
\quad\Longrightarrow\quad
a = -g(\sin\alpha + \mu\cos\alpha).
\]
Numerically: $a = -9.8(0.5 + 0.2 \times 0.866) = -9.8 \times 0.6732 \approx -6.60\ \mathrm{m\,s^{-2}}$.

\textbf{Suvat.} With $u = 3$, $v = 0$, $a = -6.60$:
\[
v^2 = u^2 + 2as \;\Longrightarrow\; 0 = 9 - 2 \times 6.60 \times s
\;\Longrightarrow\; s = \frac{9}{13.2} \approx 0.68\ \mathrm{m}.
\]
The block travels about $0.68$ m further up before stopping. (Will it then slide back? Yes, since $\tan 30^{\circ} \approx 0.577 > \mu = 0.2$: the rest condition fails, so it accelerates back down with $a = g(\sin\alpha - \mu\cos\alpha) \approx 3.20\ \mathrm{m\,s^{-2}}$.)
\end{exoresolu}

\begin{aretenir}[Key points of Section 3]
\begin{itemize}
\item On a plane inclined at $\alpha$: weight splits into $mg\sin\alpha$ down the plane and $mg\cos\alpha$ into the plane; hence $R = mg\cos\alpha$, \textbf{never} $mg$.
\item Smooth plane: $a = g\sin\alpha$. Rough plane sliding down: $a = g(\sin\alpha - \mu\cos\alpha)$.
\item Rest is possible if and only if $\tan\alpha \le \mu$, i.e. $\alpha \le \lambda$ where $\tan\lambda = \mu$ is the angle of friction.
\item For a block on a plane connected to a hanging mass: write one equation along the plane for the block, one vertical equation for the hanging mass, then add to eliminate $T$.
\item If the assumed direction of motion gives $a < 0$, reverse the direction \emph{and} the friction, and solve again — or conclude the system stays at rest.
\item After the string goes slack, recompute the acceleration with $T = 0$ before applying suvat.
\end{itemize}
\end{aretenir}

\begin{exercice}[Practice]
A block of mass $5$ kg rests on a rough plane inclined at $25^{\circ}$ to the horizontal, with $\mu = 0.3$. It is connected by a light inextensible string parallel to the plane, over a smooth pulley at the top, to a hanging particle of mass $3$ kg. Take $g = 9.8\ \mathrm{m\,s^{-2}}$.
\begin{enumerate}
\item Show that the system does move, and state which way.
\item Find the acceleration and the tension in the string.
\item The hanging mass hits the ground after descending $1.5$ m. Find the speed of the block at that instant.
\item After the string goes slack, find the further distance the block moves up the plane before momentary rest.
\end{enumerate}
\end{exercice}

\begin{corrige}[Exercise 1]
\textbf{1.} Driving force (hanging side): $3 \times 9.8 = 29.4$ N down. Resisting if $A$ moves up: $5 \times 9.8 \times \sin 25^{\circ} \approx 20.7$ N plus maximum friction $0.3 \times 5 \times 9.8 \times \cos 25^{\circ} \approx 13.3$ N, total $\approx 34.0$ N $> 29.4$ N. Check the other direction: driving force down the plane is $20.7$ N against $29.4$ N (weight of $B$) plus friction up the plane (max $13.3$ N) — impossible. Since $20.7 + 13.3 = 34.0 > 29.4 > 20.7 - 13.3 = 7.4$, friction can hold the system: \textbf{it does not move} (static equilibrium, friction $= 29.4 - 20.7 = 8.7$ N up the plane, below the limiting $13.3$ N). 

\textbf{2–4.} Because the system remains at rest, parts 2, 3 and 4 have no physical solution with the given data. \emph{If} the plane were smoother (e.g. $\mu = 0.1$), then motion would occur with the $3$ kg mass descending: $a = \dfrac{29.4 - 20.7 - 4.44}{8} \approx 0.53\ \mathrm{m\,s^{-2}}$, $T = 3(9.8 - 0.53) \approx 27.8$ N; after $1.5$ m, $v = \sqrt{2 \times 0.53 \times 1.5} \approx 1.26\ \mathrm{m\,s^{-1}}$; after the string goes slack, deceleration $g(\sin 25^{\circ} + 0.1\cos 25^{\circ}) \approx 5.03\ \mathrm{m\,s^{-2}}$, giving further distance $s = \dfrac{1.26^2}{2 \times 5.03} \approx 0.16$ m. The essential lesson: \emph{always test whether the system moves before computing an acceleration}.
\end{corrige}

\begin{savaistu}[Inclined planes in daily life and in history]
Galileo Galilei used inclined planes precisely because $a = g\sin\alpha$ is a \emph{diluted} version of free fall: with the crude clocks of his time, a ball rolling slowly down a gentle slope allowed him to measure that distances grow as the square of time — the birth of kinematics. Closer to home, every ramp at a Douala building site, every switchback on the Bamenda escarpment road, is an applied lesson in $mg\sin\alpha$: engineers choose the gradient so that vehicles can climb without the driving wheels losing grip, i.e. so that the required friction stays below $\mu R = \mu mg\cos\alpha$.
\end{savaistu}

\coursec{Section 4: Atwood's Machine and Multi-Stage Motion}

In Section 3 we analysed particles connected across smooth and rough inclined planes. The core strategy was to write Newton's second law \emph{separately} for each particle, then couple the equations through the common tension and the common magnitude of acceleration. We now apply this same rigorous method to the classic \cle{Atwood's machine}: two masses hanging vertically from a light inextensible string passing over a smooth fixed pulley. This system is a favourite of GCE examiners because it yields clean algebra, illustrates the string constraint perfectly, and leads naturally to \emph{multi-stage} problems where the forces change abruptly during the motion.

\courssub{The Atwood's Machine: Setting Up the Model}

\begin{experience}[Observing an Atwood's machine]
Hang a light string over a smooth fixed pulley (a well-oiled bicycle wheel works well in the school laboratory) and attach a mass of $200\ \mathrm{g}$ to one end and $100\ \mathrm{g}$ to the other. Release the system from rest. The heavier mass descends, the lighter one rises, and --- this is the key observation --- both move \emph{slowly}, much more slowly than a freely falling stone. The system acts as a ``diluted gravity'' device: by choosing the masses appropriately, we can make the acceleration as small as we wish. This is precisely why George Atwood invented it in 1784: to study accelerated motion with the crude timing devices available at the time.
\end{experience}

\begin{definition}[Atwood's machine]
An \cle{Atwood's machine} consists of two particles of masses $m$ and $M$ (with $M>m$) attached to the ends of a \textbf{light inextensible string} passing over a \textbf{smooth fixed pulley}, and released from rest. The modelling assumptions have precise consequences:
\begin{itemize}
\item \textbf{Light string:} its mass is negligible, so the tension $T$ is the same throughout the string.
\item \textbf{Smooth pulley:} no friction at the axle, so the tension is the same on both sides of the pulley.
\item \textbf{Inextensible string:} its length does not change, therefore the two particles always have accelerations of the same magnitude $a$, but in opposite vertical directions (one up, one down).
\end{itemize}
\end{definition}

\begin{popfigure}
\begin{center}
\begin{tikzpicture}[scale=1.0,>=stealth]
% support
\fill[poppurple!25] (-1.2,3.4) rectangle (1.2,3.7);
\draw[thick,popdark] (-1.2,3.4) -- (1.2,3.4);
% pulley
\draw[thick,popdark,fill=popblueL] (0,2.6) circle (0.55);
\fill[popdark] (0,2.6) circle (0.05);
\draw[thick,popdark] (0,2.6) -- (0,3.4);
% string
\draw[thick,poporange] (-1.4,0.4) -- (-1.4,2.6) arc[start angle=180,end angle=0,radius=0.55] (1.4,2.6) -- (1.4,1.2);
% masses
\draw[thick,popdark,fill=popgreenL] (-1.75,-0.1) rectangle (-1.05,0.4);
\node[popdark] at (-1.4,0.15) {$m$};
\draw[thick,popdark,fill=poppink!40] (0.95,0.6) rectangle (1.85,1.2);
\node[popdark] at (1.4,0.9) {$M$};
% tension arrows
\draw[->,thick,popblue] (-1.4,0.4) -- (-1.4,1.3) node[midway,left]{$T$};
\draw[->,thick,popblue] (1.4,1.2) -- (1.4,2.1) node[midway,right]{$T$};
% weight arrows
\draw[->,thick,popdark] (-1.4,-0.1) -- (-1.4,-1.0) node[midway,left]{$mg$};
\draw[->,thick,popdark] (1.4,0.6) -- (1.4,-0.4) node[midway,right]{$Mg$};
% acceleration arrows
\draw[->,very thick,popgreen] (-2.3,0.15) -- (-2.3,1.05) node[midway,left]{$a$};
\draw[->,very thick,popgreen] (2.4,0.9) -- (2.4,0.0) node[midway,right]{$a$};
\end{tikzpicture}
\end{center}
\legende{Atwood's machine: forces on each mass and the common magnitude $a$ of the accelerations.}
\end{popfigure}

\textbf{Equations of motion.} Since $M>m$, mass $M$ accelerates \emph{downwards} and mass $m$ accelerates \emph{upwards}. The most efficient sign convention is to choose the positive direction \emph{for each particle separately} to be its actual direction of motion. This automatically respects the string constraint (equal magnitude $a$) and avoids sign errors.

\begin{itemize}
\item For $M$ (positive direction \textbf{downwards}): the resultant force is $Mg - T$. Newton's second law gives
\[ Mg - T = Ma. \tag{1} \]
\item For $m$ (positive direction \textbf{upwards}): the resultant force is $T - mg$. Newton's second law gives
\[ T - mg = ma. \tag{2} \]
\end{itemize}

Adding equations (1) and (2) eliminates the unknown tension $T$:
\[ Mg - mg = Ma + ma \quad\Longrightarrow\quad (M-m)g = (M+m)a. \]

\begin{propriete}[Standard results for the Atwood's machine]
For an Atwood's machine with $M>m$, released from rest:
\[ a = \frac{(M-m)g}{M+m}, \qquad T = \frac{2Mmg}{M+m}. \]
The tension always satisfies $mg < T < Mg$: it exceeds the weight of the rising mass (it must pull it up) but is less than the weight of the falling mass (which must have a net downward force). The derivation above \textbf{is examinable} --- do not quote the formulae without justification unless the question explicitly says ``you may assume''.
\end{propriete}

\begin{exemplebox}[Quick check of the bounds]
Take $M = 5\ \mathrm{kg}$, $m = 3\ \mathrm{kg}$, $g = 9.8\ \mathrm{m\,s^{-2}}$. Then
\[ a = \frac{(5-3) \times 9.8}{5+3} = \frac{2 \times 9.8}{8} = 2.45\ \mathrm{m\,s^{-2}}, \qquad T = \frac{2 \times 5 \times 3 \times 9.8}{8} = 36.75\ \mathrm{N}. \]
Check: $mg = 29.4\ \mathrm{N} < 36.75\ \mathrm{N} < 49\ \mathrm{N} = Mg$. \checkmark
\end{exemplebox}

\begin{attention}[Sign errors with opposite directions]
The two masses move in \textbf{opposite} directions. A classic error is to write both equations with a single global ``up positive'' convention, obtaining $Mg - T = Ma$ and $T - mg = -ma$, and then adding to get $Mg - mg = 0$ (nonsense). \textbf{Fix the positive direction for each particle separately}, in its own direction of motion, and state it clearly on your diagram. The string constraint then simply reads: accelerations equal in magnitude, opposite in direction.
\end{attention}

\courssub{Force on the Pulley Support}

The pulley has two strands of string pulling down on it, each with tension $T$. Since the pulley is smooth and light, it merely redirects the tensions; it does not alter their magnitude.

\begin{popfigure}
\begin{center}
\begin{tikzpicture}[scale=1.0,>=stealth]
\fill[poppurple!25] (-1.2,2.2) rectangle (1.2,2.5);
\draw[thick,popdark] (-1.2,2.2) -- (1.2,2.2);
\draw[thick,popdark,fill=popblueL] (0,1.4) circle (0.55);
\fill[popdark] (0,1.4) circle (0.05);
\draw[thick,popdark] (0,1.95) -- (0,2.2);
\draw[thick,poporange] (-0.55,1.4) -- (-0.55,0.2);
\draw[thick,poporange] (0.55,1.4) -- (0.55,0.2);
\draw[->,very thick,popblue] (-0.55,0.2) -- (-0.55,-0.8) node[midway,left]{$T$};
\draw[->,very thick,popblue] (0.55,0.2) -- (0.55,-0.8) node[midway,right]{$T$};
\draw[->,very thick,popgreen] (1.6,1.4) -- (1.6,0.3) node[midway,right]{$2T$};
\draw[->,very thick,popdark] (0,2.5) -- (0,3.3) node[midway,right]{$R$};
\node[popdark,align=center] at (0,-1.3) {Support reaction $R = 2T$ (light pulley)};
\end{tikzpicture}
\end{center}
\legende{The pulley support carries the downward force $2T$ from the two strands of the string.}
\end{popfigure}

\begin{methode}[Force on the pulley: $R = 2T$]
\begin{enumerate}
\item The pulley feels two downward pulls of magnitude $T$, one from each vertical strand of the string. The total downward force on the pulley is therefore $2T$.
\item For a light pulley in vertical equilibrium, the upward reaction force $R$ from the support equals this total downward force: $R = 2T$.
\item Compare $R$ with the total weight $(M+m)g$. Using $T = \dfrac{2Mmg}{M+m}$,
\[ R = 2T = \frac{4Mmg}{M+m}. \]
The difference between the total weight and the support reaction is
\[ (M+m)g - R = (M+m)g - \frac{4Mmg}{M+m} = \frac{(M+m)^2g - 4Mmg}{M+m} = \frac{(M-m)^2g}{M+m}. \]
Notice that $(M-m)g = (M+m)a$, so $(M-m)^2g/(M+m) = (M-m)a$. Hence
\[ R = (M+m)g - (M-m)a. \]
\item The support carries \emph{less} than the combined weight $(M+m)g$. The ``missing weight'' is $(M-m)a$, which is exactly the net force acting on the two-particle system (mass $M$ accelerating down, mass $m$ accelerating up).
\end{enumerate}
\end{methode}

\begin{savaistu}[Weighing an accelerating system]
If the whole Atwood's machine (pulley, string, and masses) were placed on a very sensitive scale, the scale would read $R = 2T$ (plus the pulley's own weight if not light), not $(M+m)g$. The same principle explains why a lift carrying accelerating passengers exerts a different force on its cable than the static weight --- a theme we revisit in real-life problems later in this chapter.
\end{savaistu}

\courssub{Kinematics of the Machine}

Because the acceleration $a$ is constant, all the \emph{suvat} equations apply to each particle individually. The string constraint guarantees that at every instant:
\begin{itemize}
\item the two masses have \textbf{equal speeds} (but opposite directions);
\item the distance descended by $M$ equals the distance ascended by $m$.
\end{itemize}

\begin{exoresolu}[Basic Atwood's machine kinematics]
Two particles of masses $2\ \mathrm{kg}$ and $5\ \mathrm{kg}$ are connected by a light inextensible string passing over a smooth fixed pulley. The system is released from rest with the $5\ \mathrm{kg}$ mass $1.4\ \mathrm{m}$ above the ground. Take $g = 9.8\ \mathrm{m\,s^{-2}}$. Find (a) the acceleration, (b) the tension, (c) the speed with which the $5\ \mathrm{kg}$ mass hits the ground, (d) the time taken.
\tcblower
\textbf{Solution.} Let $M = 5\ \mathrm{kg}$ (moving down), $m = 2\ \mathrm{kg}$ (moving up). Positive direction: down for $M$, up for $m$.
\[ 5g - T = 5a \quad (1), \qquad T - 2g = 2a \quad (2). \]
(a) Adding (1) and (2): $3g = 7a$, so $a = \dfrac{3 \times 9.8}{7} = 4.2\ \mathrm{m\,s^{-2}}$.
(b) From (2): $T = 2g + 2a = 19.6 + 8.4 = 28\ \mathrm{N}$.
(c) For the $5\ \mathrm{kg}$ mass: $u = 0$, $s = 1.4\ \mathrm{m}$, $a = 4.2\ \mathrm{m\,s^{-2}}$. Using $v^2 = u^2 + 2as$:
\[ v^2 = 2 \times 4.2 \times 1.4 = 11.76 \quad\Longrightarrow\quad v = \sqrt{11.76} = 3.43\ \mathrm{m\,s^{-1}}\ (3\ \mathrm{s.f.}). \]
(d) Using $v = u + at$: $t = \dfrac{3.43}{4.2} = 0.817\ \mathrm{s}$ (3 s.f.).
\end{exoresolu}

\courssub{Multi-Stage Motion: When the String Goes Slack}

GCE questions often involve a scenario where the heavier mass starts above the ground and the lighter mass starts on (or near) the ground. Three distinct stages occur:

\textbf{Stage 1: Connected motion.} Both masses move together with acceleration $a = \dfrac{(M-m)g}{M+m}$ until $M$ hits the ground (or the string breaks). At that instant both masses have the same speed $v$.

\textbf{Stage 2: Free projectile motion (string slack).} Once $M$ lands, the string goes \emph{slack}: the tension drops instantly to zero. Mass $m$ is now moving \emph{upwards} with speed $v$, acted on by gravity alone. It continues to rise, decelerating at $g$, until it comes momentarily to rest. Using $v^2 = u^2 + 2as$ with final speed $0$ and acceleration $-g$ (taking upward as positive), the extra height gained is
\[ h = \frac{v^2}{2g}, \qquad \text{time to momentary rest: } t = \frac{v}{g}. \]

\textbf{Stage 3: The string becomes taut again (qualitative only).} After reaching its highest point, $m$ falls back down. The slack string gradually straightens, and when it becomes taut there is a sudden \emph{jerk}: an impulsive tension that almost instantaneously changes the velocities of both particles (and may set $M$ moving again if it can leave the ground). The analysis of this jerk requires \emph{impulse and momentum}, which is treated in a later chapter --- at GCE you will only ever be asked for a \emph{qualitative} comment here.

\begin{popfigure}
\begin{center}
\begin{tikzpicture}[scale=0.95,>=stealth]
% ground
\fill[popgreenL] (-3.2,-0.35) rectangle (3.2,0);
\draw[thick,popdark] (-3.2,0) -- (3.2,0);
\node[popmuted] at (2.6,-0.2) {\footnotesize ground};
% pulley
\draw[thick,popdark] (0,4.6) -- (0,4.9);
\fill[poppurple!25] (-0.8,4.9) rectangle (0.8,5.1);
\draw[thick,popdark,fill=popblueL] (0,4.1) circle (0.5);
% M on ground
\draw[thick,popdark,fill=poppink!40] (-1.75,0) rectangle (-0.95,0.6);
\node[popdark] at (-1.35,0.3) {$M$};
% slack string (curved)
\draw[thick,poporange,dashed] (-1.35,0.6) .. controls (-1.9,2.2) and (-1.1,3.4) .. (-0.5,4.1) arc[start angle=180,end angle=0,radius=0.5] (0.5,4.1) .. controls (1.3,3.2) and (1.1,2.6) .. (1.35,2.4);
% m rising
\draw[thick,popdark,fill=popgreenL] (1.0,2.4) rectangle (1.7,2.9);
\node[popdark] at (1.35,2.65) {$m$};
\draw[->,very thick,popgreen] (2.2,2.65) -- (2.2,3.75) node[midway,right]{$v$};
\draw[->,thick,popdark] (1.35,2.4) -- (1.35,1.5) node[midway,right]{$mg$};
\node[popdark,align=center] at (-1.35,-0.75) {\footnotesize $M$ at rest on ground};
\node[popdark,align=center] at (1.6,4.6) {\footnotesize slack string: $T = 0$};
\end{tikzpicture}
\end{center}
\legende{Stage 2: $M$ has landed, the string is slack, and $m$ rises freely under gravity alone.}
\end{popfigure}

\begin{methode}[Solving multi-stage problems]
\begin{enumerate}
\item \textbf{Draw a separate force diagram for each stage.} Forces change between stages --- never recycle a diagram.
\item \textbf{Stage 1 (connected):} Apply Newton's second law to each particle, solve for $a$ and $T$, then use \emph{suvat} to find the speed $v$ at the moment the stage ends (when a mass hits the ground or the string breaks).
\item \textbf{Carry over the velocity:} The final velocity of one stage is the initial velocity of the next. Both masses have the same speed $v$ at the transition.
\item \textbf{Stage 2 (slack string):} The free particle moves under gravity alone ($a = \pm g$). Use $v^2 = u^2 - 2gh$ for the extra height and $t = u/g$ for the time to momentary rest.
\item \textbf{Assemble} total times and total heights from the stages, and comment qualitatively on any jerk if asked.
\end{enumerate}
\end{methode}

\begin{attention}[Do not keep the tension after the string goes slack]
The single most common error in these questions is to continue using $T = \dfrac{2Mmg}{M+m}$ --- or indeed any tension at all --- after $M$ has landed. Once the string is slack, $T = 0$ and the rising mass is a \textbf{projectile} with acceleration $g$ downwards. Its deceleration is $g$, not the Stage-1 value $a$. Similarly, never use the Stage-1 acceleration to compute the extra height or the time to rest.
\end{attention}

\begin{exoresolu}[GCE-style multi-stage problem]
Particles $P$ and $Q$, of masses $0.3\ \mathrm{kg}$ and $0.5\ \mathrm{kg}$ respectively, are attached to the ends of a light inextensible string which passes over a smooth fixed pulley. The system is released from rest with both particles $0.8\ \mathrm{m}$ above horizontal ground. Take $g = 9.8\ \mathrm{m\,s^{-2}}$.
(a) Find the acceleration of the system and the tension in the string.
(b) Find the speed of $Q$ when it reaches the ground.
(c) After $Q$ hits the ground (and does not rebound), find the greatest height above the ground reached by $P$.
(d) Find the total time from release until $P$ reaches its greatest height.
(e) State, with a reason, what happens to the tension in the string at the instant $Q$ lands.
\tcblower
\textbf{Solution.} (a) Let $M = 0.5\ \mathrm{kg}$ (down positive) and $m = 0.3\ \mathrm{kg}$ (up positive).
\[ 0.5g - T = 0.5a \quad (1), \qquad T - 0.3g = 0.3a \quad (2). \]
Adding: $0.2g = 0.8a$, so $a = \dfrac{0.2 \times 9.8}{0.8} = 2.45\ \mathrm{m\,s^{-2}}$.
From (2): $T = 0.3(g + a) = 0.3 \times 12.25 = 3.675\ \mathrm{N}$.

(b) Stage 1 ends when $Q$ has fallen $0.8\ \mathrm{m}$. With $u = 0$:
\[ v^2 = 2 \times 2.45 \times 0.8 = 3.92 \quad\Longrightarrow\quad v = \sqrt{3.92} = 1.98\ \mathrm{m\,s^{-1}}\ (3\ \mathrm{s.f.}). \]
At this instant $P$ is $0.8 + 0.8 = 1.6\ \mathrm{m}$ above the ground, moving up at $1.98\ \mathrm{m\,s^{-1}}$.

(c) Stage 2: string slack, $P$ rises under gravity alone. Extra height:
\[ h = \frac{v^2}{2g} = \frac{3.92}{19.6} = 0.2\ \mathrm{m}. \]
Greatest height above ground $= 1.6 + 0.2 = \mathbf{1.8\ m}$.

(d) Stage 1 time: $t_1 = \dfrac{v}{a} = \dfrac{1.98}{2.45} = 0.808\ \mathrm{s}$.
Stage 2 time: $t_2 = \dfrac{v}{g} = \dfrac{1.98}{9.8} = 0.202\ \mathrm{s}$.
Total time $= t_1 + t_2 = 1.01\ \mathrm{s}$ (3 s.f.).

(e) When $Q$ lands, the string goes slack and the tension falls instantaneously to zero, because there is no longer anything pulling the string taut on $Q$'s side. $P$ then moves under gravity alone until the string becomes taut again, at which point a jerk (impulsive tension) occurs --- its analysis requires impulse and momentum, which is beyond this chapter.
\end{exoresolu}

\begin{exercice}[Practice: full multi-stage problem]
Two particles of masses $1\ \mathrm{kg}$ and $3\ \mathrm{kg}$ hang from a light inextensible string over a smooth fixed pulley. The system is released from rest with the $3\ \mathrm{kg}$ mass $2\ \mathrm{m}$ above the ground and the $1\ \mathrm{kg}$ mass on the ground. Take $g = 9.8\ \mathrm{m\,s^{-2}}$.
(a) Show that the acceleration is $4.9\ \mathrm{m\,s^{-2}}$ and find the tension.
(b) Find the reaction at the pulley support while the system is moving, and verify that $R = (M+m)g - (M-m)a$.
(c) Find the speed of the $3\ \mathrm{kg}$ mass as it hits the ground.
(d) Find the greatest height reached by the $1\ \mathrm{kg}$ mass and the total time taken to reach it.
\end{exercice}

\begin{corrige}[Exercise 1]
(a) Let $M = 3\ \mathrm{kg}$ (down positive), $m = 1\ \mathrm{kg}$ (up positive).
\[ 3g - T = 3a \quad (1), \qquad T - g = a \quad (2). \]
Adding: $2g = 4a \implies a = \frac{g}{2} = 4.9\ \mathrm{m\,s^{-2}}$. From (2): $T = g + a = 9.8 + 4.9 = 14.7\ \mathrm{N}$.

(b) $R = 2T = 29.4\ \mathrm{N}$.
Check: $(M+m)g - (M-m)a = 4 \times 9.8 - 2 \times 4.9 = 39.2 - 9.8 = 29.4\ \mathrm{N} = R$. \checkmark

(c) $v^2 = 2as = 2 \times 4.9 \times 2 = 19.6$, so $v = \sqrt{19.6} = 4.43\ \mathrm{m\,s^{-1}}$ (3 s.f.).

(d) At the moment $M$ hits the ground, $m$ is $2\ \mathrm{m}$ above the ground (it has risen the same distance $M$ fell), moving up at $4.43\ \mathrm{m\,s^{-1}}$.
Extra height: $h = \dfrac{v^2}{2g} = \dfrac{19.6}{19.6} = 1\ \mathrm{m}$.
Greatest height $= 2 + 1 = \mathbf{3\ m}$.
Times: $t_1 = \dfrac{v}{a} = \dfrac{4.43}{4.9} = 0.904\ \mathrm{s}$, $t_2 = \dfrac{v}{g} = \dfrac{4.43}{9.8} = 0.452\ \mathrm{s}$.
Total time $= 1.36\ \mathrm{s}$ (3 s.f.).
\end{corrige}

\begin{aretenir}[Key points of Section 4]
\begin{itemize}
\item Atwood's machine ($M>m$): $Mg - T = Ma$, $T - mg = ma$, giving $a = \dfrac{(M-m)g}{M+m}$ and $T = \dfrac{2Mmg}{M+m}$, with $mg < T < Mg$.
\item The pulley support carries $R = 2T$, which is \emph{less} than $(M+m)g$; the deficit is $(M-m)a$, the net force accelerating the system.
\item The two masses always have equal speeds in opposite directions; distances risen and fallen are equal.
\item Multi-stage motion: when a mass lands or the string breaks, the string goes slack ($T = 0$) and the free mass continues as a projectile under gravity alone; the final velocity of each stage seeds the next.
\item The jerk when the string re-tightens involves impulsive tension --- qualitative comment only at this stage.
\end{itemize}
\end{aretenir}

\coursec{Section 5: Pulley Systems with More Than Two Masses and Variable Acceleration}

In Section 4 we mastered Atwood's machine and multi-stage motion where the acceleration was constant. We now extend our toolkit to systems with \textbf{three or more particles} and to situations where the \textbf{driving force changes with time or position}, leading to variable acceleration. The fundamental principles remain Newton's second law and the geometric constraints imposed by inextensible strings, but the algebra becomes richer and the integration techniques more varied.

\courssub{Systems of Three or More Particles on Horizontal Surfaces}

Consider a light inextensible string passing over a smooth pulley at the edge of a smooth horizontal table. Two particles of masses $m_1$ and $m_2$ lie on the table, connected in series, and a third particle of mass $m_3$ hangs vertically. The string is taut throughout. Because the string is inextensible, all three particles move with the \textbf{same magnitude of acceleration} $a$. However, the string is now divided into \textbf{two distinct segments}: one between $m_1$ and $m_2$ (tension $T_1$) and another between $m_2$ and the pulley (tension $T_2$). The hanging mass feels tension $T_2$ as well (same segment).

\begin{popfigure}
\begin{tikzpicture}[scale=0.9, every node/.style={font=\small}]
% Table
\draw[popdark, thick] (-3,0) -- (3,0);
\draw[popdark, thick] (-3,-0.2) -- (3,-0.2);
\draw[popdark] (-3,0) -- (-3,-0.2);
\draw[popdark] (3,0) -- (3,-0.2);
% Pulley at right edge
\draw[popdark, thick] (3,0) arc (0:90:0.5) -- (-0.5,0.5);
\draw[popdark, fill=popblueL] (3,0.5) circle (0.15);
% Mass m1 on table
\draw[poporange, fill=poporange!20] (-2.5,-0.1) rectangle (-1.5,0.5);
\node at (-2,0.2) {$m_1$};
% Mass m2 on table
\draw[poporange, fill=poporange!20] (-0.5,-0.1) rectangle (0.5,0.5);
\node at (0,0.2) {$m_2$};
% String on table: from m1 to m2
\draw[popblue, thick] (-1.5,0.2) -- (-0.5,0.2);
% String from m2 to pulley
\draw[popblue, thick] (0.5,0.2) -- (2.5,0.2);
\draw[popblue, thick] (2.5,0.2) -- (2.5,0.5) arc (0:90:0.5);
% Hanging mass m3
\draw[poporange, fill=poporange!20] (3.2,0.5) -- (3.2,-1.5) -- (3.8,-1.5) -- (3.8,0.5) -- cycle;
\node at (3.5,-0.5) {$m_3$};
% Tension labels
\node[popblue] at (-1,0.55) {$T_1$};
\node[popblue] at (1.5,0.55) {$T_2$};
\node[popblue] at (3.5,0.8) {$T_2$};
% Acceleration arrows
\draw[->, popgreen, thick] (-2,0.7) -- (-1,0.7) node[midway,above] {$a$};
\draw[->, popgreen, thick] (0,0.7) -- (1,0.7) node[midway,above] {$a$};
\draw[->, popgreen, thick] (3.5,-0.8) -- (3.5,-1.3) node[midway,right] {$a$};
% Weight arrows
\draw[->, popdark] (-2,-0.4) -- (-2,-0.8) node[below] {$m_1g$};
\draw[->, popdark] (0,-0.4) -- (0,-0.8) node[below] {$m_2g$};
\draw[->, popdark] (3.5,-1.5) -- (3.5,-2) node[below] {$m_3g$};
% Normal reactions
\draw[->, poppurple] (-2,-0.1) -- (-2,0.5) node[above] {$R_1$};
\draw[->, poppurple] (0,-0.1) -- (0,0.5) node[above] {$R_2$};
\end{tikzpicture}
\legende{Three-mass system: two masses on a smooth table linked in series to a hanging mass. Note the two different tensions $T_1$ and $T_2$.}
\end{popfigure}

\begin{methode}[One unknown tension per string segment]
When a system contains several particles connected by strings, \textbf{each continuous segment of string that is separated by a particle or a pulley generally carries a different tension}. Label them $T_1$, $T_2$, \dots\ and write Newton's second law for each particle individually. The constraint of equal acceleration magnitude (for particles on the same taut string) links the equations.
\end{methode}

\begin{exoresolu}[Three masses on a smooth table and a hanging mass]
A particle $A$ of mass $2\,\mathrm{kg}$ and a particle $B$ of mass $3\,\mathrm{kg}$ rest on a smooth horizontal table, connected by a light inextensible string. The string passes over a smooth pulley at the edge of the table and carries a particle $C$ of mass $5\,\mathrm{kg}$ hanging freely. The system is released from rest. Find the acceleration of the system and the tensions $T_1$ (between $A$ and $B$) and $T_2$ (between $B$ and $C$).
\tcblower
\textbf{Solution.}
Let the acceleration be $a$ (to the right for $A$ and $B$, downward for $C$).
\begin{itemize}
\item For $A$ (only horizontal force is $T_1$): $T_1 = 2a$. \hfill (1)
\item For $B$ (horizontal forces: $T_2$ to the right, $T_1$ to the left): $T_2 - T_1 = 3a$. \hfill (2)
\item For $C$ (vertical forces: weight $5g$ down, tension $T_2$ up): $5g - T_2 = 5a$. \hfill (3)
\end{itemize}
Add (1), (2) and (3): $5g = (2+3+5)a = 10a \;\Rightarrow\; a = \frac{g}{2} = 4.9\,\mathrm{m\,s^{-2}}$.
Substitute into (1): $T_1 = 2 \times 4.9 = 9.8\,\mathrm{N}$.
Substitute into (3): $T_2 = 5g - 5a = 5(g - a) = 5(9.8 - 4.9) = 24.5\,\mathrm{N}$.
\emph{Check:} Equation (2) gives $24.5 - 9.8 = 14.7 = 3 \times 4.9$. Consistent.
\end{exoresolu}

\begin{propriete}[Whole-system equation for connected particles]
For any system of $n$ particles connected by light inextensible strings passing over smooth pulleys, if all particles move with the \textbf{same magnitude of acceleration} $a$, then
\[
a = \frac{\text{sum of external driving forces in the direction of motion}}{\text{total mass of the system}}.
\]
Internal tensions cancel out when the equations are summed. This gives a quick way to find $a$ before solving for individual tensions.
\end{propriete}

\courssub{Movable Pulleys (Simple Cases)}

A \textbf{movable pulley} is a pulley that can accelerate. Consider a light pulley of negligible mass carrying a particle of mass $M$, with a light inextensible string passing over it. One end of the string is fixed, the other is pulled with acceleration $a_1$, while the pulley (and mass $M$) accelerates with $a_2$. The geometry of the string gives a constraint: if the string length is constant, then $2a_2 = a_1$ (or more generally $2a_{\text{pulley}} = a_1 + a_2$ for two free ends). In the GCE Advanced Level syllabus, only simple configurations where the constraint is linear are examined.

\begin{exemplebox}[Mass on a movable pulley]
A light smooth pulley carries a mass $M = 4\,\mathrm{kg}$. A light inextensible string passes over the pulley; one end is attached to a fixed support, the other end is pulled vertically upward with a force $F = 60\,\mathrm{N}$. Find the acceleration of the mass $M$.
\textbf{Solution.} Let $a$ be the upward acceleration of the mass $M$ (and the pulley). The free end of the string has acceleration $2a$ upward (constraint: if the pulley rises by $y$, the free end must be pulled up by $2y$ to keep the string length constant). The tension $T$ is the same throughout the string (light pulley, smooth). For the mass $M$: $2T - Mg = Ma$. For the free end (massless): $F - T = 0 \Rightarrow T = F = 60\,\mathrm{N}$. Then $2(60) - 4g = 4a \Rightarrow 120 - 39.2 = 4a \Rightarrow a = 20.2\,\mathrm{m\,s^{-2}}$ upward.
\end{exemplebox}

\courssub{Variable Acceleration: Chain Sliding off a Table}

Many real systems involve a driving force that changes as the configuration changes. A classic example is a \textbf{uniform chain} initially coiled at the edge of a smooth horizontal table, with a small length $x_0$ hanging over the edge. As the chain slides, the hanging length $x$ increases, so the driving weight $\lambda x g$ (where $\lambda$ is the mass per unit length) increases. The acceleration is no longer constant.

\begin{popfigure}
\begin{tikzpicture}[scale=0.8, every node/.style={font=\small}]
% Table surface
\draw[popdark, thick] (-4,0) -- (4,0);
\draw[popdark, thick] (-4,-0.3) -- (4,-0.3);
\draw[popdark] (-4,0) -- (-4,-0.3);
\draw[popdark] (4,0) -- (4,-0.3);
% Chain on table (coiled representation)
\foreach \i in {-3.5,-3,-2.5,-2,-1.5,-1,-0.5,0} {
  \draw[poporange, thick] (\i, -0.1) -- (\i, 0.1);
}
% Hanging part
\draw[poporange, thick] (0,0) -- (0,-2.5);
\draw[poporange, thick] (0.2,0) -- (0.2,-2.5);
\draw[poporange, fill=poporange!20] (0,-2.5) rectangle (0.2,-2.7);
% Length labels
\draw[<->, popblue] (-3.5,0.4) -- (0,0.4) node[midway,above] {$L-x$ (on table)};
\draw[<->, popblue] (0.5,0) -- (0.5,-2.5) node[midway,right] {$x$ (hanging)};
% Forces
\draw[->, popgreen, thick] (0.1,-1.25) -- (1.5,-1.25) node[right] {Driving weight $\lambda x g$};
\draw[->, poppurple, thick] (-2,0.4) -- (-2,1) node[above] {No friction (smooth)};
% Mass per unit length label
\node[popdark] at (-3.5,-0.6) {$\lambda = M/L$};
\end{tikzpicture}
\legende{Uniform chain of length $L$ and mass $M$ sliding off a smooth table. The hanging length $x$ varies with time.}
\end{popfigure}

Let the total mass be $M$, total length $L$, so linear density $\lambda = M/L$. At a given instant, hanging length $= x$, mass on table $= \lambda(L-x)$, hanging mass $= \lambda x$. The only horizontal external force on the whole chain is the weight of the hanging part (the table is smooth, so no friction). Applying the whole-system equation:
\[
a = \frac{\text{net driving force}}{\text{total mass}} = \frac{\lambda x g}{M} = \frac{\lambda x g}{\lambda L} = \frac{g}{L}\,x.
\]
Thus acceleration is \textbf{proportional to the hanging length} $x$: $a = \frac{g}{L}x$. This is a variable acceleration problem where $a$ is a function of displacement $x$.

\begin{popfigure}
\begin{tikzpicture}
\begin{axis}[
  width=10cm, height=6cm,
  axis lines=middle,
  xlabel={$x$ (hanging length)}, ylabel={$a$ (acceleration)},
  xmin=0, xmax=5, ymin=0, ymax=10,
  xtick={0,2.5,5}, xticklabels={$0$, $L/2$, $L$},
  ytick={0,5,10}, yticklabels={$0$, $g/2$, $g$},
  domain=0:5, samples=50,
  tick label style={font=\small},
  label style={font=\small},
  grid=major, grid style={popline},
]
\addplot[popblue, thick] {2*x}; % g/L = 2 for L=4.9, g=9.8
\node[popblue, anchor=south west] at (axis cs:4,8) {$a = \frac{g}{L}x$};
\end{axis}
\end{tikzpicture}
\legende{Acceleration $a$ as a linear function of hanging length $x$ for a uniform chain on a smooth table.}
\end{popfigure}

\begin{methode}[Solving variable acceleration problems]
When acceleration is not constant, the \textbf{suvat equations do not apply}. Instead:
\begin{enumerate}
\item Express acceleration as a function of the relevant variable: $a = f(x)$ or $a = f(t)$.
\item If $a = f(x)$: use $a = v\,\frac{dv}{dx}$. Separate variables: $v\,dv = f(x)\,dx$ and integrate.
\item If $a = f(t)$: use $a = \frac{dv}{dt}$. Separate variables: $dv = f(t)\,dt$ and integrate.
\item Apply initial conditions to determine constants of integration.
\item If needed, find $x(t)$ by integrating $v = dx/dt$.
\end{enumerate}
\end{methode}

\begin{exoresolu}[Speed of the chain when it leaves the table]
A uniform chain of length $L = 2\,\mathrm{m}$ and mass $M = 0.5\,\mathrm{kg}$ is placed on a smooth horizontal table with $x_0 = 0.2\,\mathrm{m}$ hanging over the edge. The chain is released from rest. Find the speed of the chain at the instant the last link leaves the table.
\tcblower
\textbf{Solution.}
Linear density $\lambda = M/L = 0.25\,\mathrm{kg/m}$.
Equation of motion for the whole chain: $a = \frac{g}{L}x = \frac{9.8}{2}x = 4.9x$.
Since $a = f(x)$, use $a = v\frac{dv}{dx}$:
\[
v\frac{dv}{dx} = 4.9x \quad\Rightarrow\quad v\,dv = 4.9x\,dx.
\]
Integrate from initial state ($x=0.2$, $v=0$) to final state ($x=L=2$, $v=v_f$):
\[
\int_0^{v_f} v\,dv = 4.9\int_{0.2}^{2} x\,dx
\;\Rightarrow\; \frac{1}{2}v_f^2 = 4.9\left[\frac{x^2}{2}\right]_{0.2}^{2}
= 2.45\,(4 - 0.04) = 2.45 \times 3.96 = 9.702.
\]
Thus $v_f^2 = 19.404 \;\Rightarrow\; v_f = \sqrt{19.404} \approx 4.40\,\mathrm{m\,s^{-1}}$.
\emph{Alternative energy method:} The work done by gravity on the hanging part equals the gain in kinetic energy of the whole chain. The weight of the hanging part is $\lambda g x$, and it acts through a distance $dx$ as $x$ increases. Total work $= \int_{x_0}^L \lambda g x\,dx = \frac{1}{2}\lambda g (L^2 - x_0^2)$. Equating to $\frac{1}{2}M v^2 = \frac{1}{2}\lambda L v^2$ gives $v = \sqrt{\frac{g}{L}(L^2 - x_0^2)} = \sqrt{gL - \frac{g x_0^2}{L}}$, which matches the integration result.
\end{exoresolu}

\courssub{Variable Acceleration with Time-Dependent Force}

Another common case is when the net force varies explicitly with time.

\begin{exoresolu}[Particle pulled by a time-varying force]
A particle of mass $2\,\mathrm{kg}$ rests on a smooth horizontal table. It is pulled by a horizontal force $F(t) = 6t\,\mathrm{N}$ (where $t$ is in seconds) applied via a light inextensible string. The particle starts from rest at $t=0$. Find its velocity and displacement at $t=3\,\mathrm{s}$.
\tcblower
\textbf{Solution.}
Newton's second law: $F = ma \;\Rightarrow\; 6t = 2a \;\Rightarrow\; a = 3t$.
Since $a = f(t)$, use $a = \frac{dv}{dt}$:
\[
\frac{dv}{dt} = 3t \quad\Rightarrow\quad dv = 3t\,dt.
\]
Integrate from $t=0$ ($v=0$) to $t$:
\[
\int_0^v dv = 3\int_0^t t\,dt \;\Rightarrow\; v = \frac{3}{2}t^2.
\]
At $t=3\,\mathrm{s}$: $v = \frac{3}{2} \times 9 = 13.5\,\mathrm{m\,s^{-1}}$.
Now find displacement: $v = \frac{dx}{dt} = \frac{3}{2}t^2 \;\Rightarrow\; dx = \frac{3}{2}t^2\,dt$.
Integrate from $t=0$ ($x=0$) to $t=3$:
\[
x = \frac{3}{2}\int_0^3 t^2\,dt = \frac{3}{2}\left[\frac{t^3}{3}\right]_0^3 = \frac{1}{2} \times 27 = 13.5\,\mathrm{m}.
\]
\end{exoresolu}

\courssub{Further Examples of Variable Acceleration}

The chain problem and the time-varying force are prototypes. Other situations include:
\begin{itemize}
\item A rope being pulled up with a force that varies with the length already lifted.
\item A particle attached to a spring on a rough surface where friction depends on position (though spring force is usually treated in the oscillations chapter).
\item Systems where mass is added or removed (e.g., sand falling onto a moving conveyor belt) — treated via momentum, not $F=ma$ directly.
\end{itemize}
In all cases, the strategy is: \textbf{write Newton's second law for the system at a generic instant}, identify the variable acceleration, and integrate appropriately.

\begin{attention}[Common mistake: using suvat for variable acceleration]
\textbf{Never} use $v^2 = u^2 + 2as$, $s = ut + \frac{1}{2}at^2$, etc., when acceleration is not constant. These equations are derived assuming constant $a$. If $a$ varies with $x$ or $t$, you \textbf{must} integrate using $a = v\,dv/dx$ or $a = dv/dt$.
\end{attention}

\begin{attention}[Common mistake: assuming a single tension throughout]
In a system with \textbf{multiple string segments} (separated by particles or pulleys), the tensions are generally \textbf{different}. Always label them $T_1$, $T_2$, \dots\ and write separate equations for each particle. Only if the string is continuous over a smooth pulley with no mass attached at the junction does the tension remain the same on both sides.
\end{attention}

\begin{savaistu}[Real-life application: Cable deployment]
The chain sliding off a table models the deployment of a cable from a drum, the unwinding of a hose, or the motion of a flexible rope in a pulley system. Engineers use the variable acceleration analysis to design braking systems that prevent excessive speeds when the driving weight increases. In Cameroon, similar principles apply to the design of cable-car systems in mountainous regions like the Mandara Mountains.
\end{savaistu}

\begin{aretenir}[Key points for Section 5]
\begin{itemize}
\item For systems with $\ge 3$ particles: write one $F=ma$ equation per particle; use the constraint of equal acceleration magnitude along each taut string segment.
\item Different string segments $\rightarrow$ different tensions ($T_1$, $T_2$, \dots).
\item Whole-system shortcut: $a = \frac{\sum \text{external driving forces}}{\sum \text{masses}}$ (valid when all parts share the same $|a|$).
\item Movable pulleys introduce geometric constraints (e.g., $2a_{\text{pulley}} = a_1 + a_2$); treat only simple cases.
\item Variable acceleration: if $a = f(x)$, use $v\,dv = f(x)\,dx$; if $a = f(t)$, use $dv = f(t)\,dt$. Integrate with initial conditions.
\item \textbf{Do not use suvat} when $a$ is not constant.
\end{itemize}
\end{aretenir}

\begin{exercice}[Practice problems]
\begin{enumerate}
\item Three particles of masses $2\,\mathrm{kg}$, $3\,\mathrm{kg}$ and $5\,\mathrm{kg}$ are connected by light inextensible strings. The $2\,\mathrm{kg}$ and $3\,\mathrm{kg}$ masses lie on a smooth horizontal table, the $5\,\mathrm{kg}$ mass hangs over a smooth pulley at the edge. Find the acceleration and the two tensions.
\item A uniform chain of length $3\,\mathrm{m}$ and mass $1.2\,\mathrm{kg}$ rests on a smooth table with $0.5\,\mathrm{m}$ hanging over the edge. It is released from rest. Find the speed when the chain leaves the table.
\item A light movable pulley carries a mass of $6\,\mathrm{kg}$. A light string passes over it; one end is fixed, the other is pulled upward with a constant force of $80\,\mathrm{N}$. Find the acceleration of the $6\,\mathrm{kg}$ mass.
\item A particle of mass $m$ on a smooth horizontal table is connected by a string passing over a pulley at the edge to a particle of mass $2m$ hanging vertically. A second string connects the hanging particle to a third particle of mass $3m$ hanging below it. Find the acceleration of the system and the tensions in both strings.
\item A particle of mass $3\,\mathrm{kg}$ on a smooth horizontal plane is acted upon by a force $F = (4t + 2)\,\mathrm{N}$ in a fixed horizontal direction. If the particle starts from rest, find its velocity and displacement after $4\,\mathrm{s}$.
\end{enumerate}
\end{exercice}

\begin{corrige}[Exercise 1]
$m_1=2$, $m_2=3$ on table, $m_3=5$ hanging. Whole system: $a = \frac{5g}{2+3+5} = \frac{g}{2}=4.9\,\mathrm{m/s^2}$.
Tensions: For $m_1$: $T_1 = 2a = 9.8\,\mathrm{N}$. For $m_3$: $5g - T_2 = 5a \Rightarrow T_2 = 5(g-a)=24.5\,\mathrm{N}$.
\end{corrige}

\begin{corrige}[Exercise 2]
$\lambda = 0.4\,\mathrm{kg/m}$, $L=3$, $x_0=0.5$. $a = \frac{g}{L}x = \frac{9.8}{3}x$.
$v\,dv = \frac{9.8}{3}x\,dx$. Integrate: $\frac{1}{2}v^2 = \frac{9.8}{6}(x^2 - 0.5^2)$.
At $x=3$: $v^2 = \frac{9.8}{3}(9 - 0.25) = \frac{9.8}{3}\times 8.75 = 28.583\dots \Rightarrow v \approx 5.35\,\mathrm{m/s}$.
\end{corrige}

\begin{corrige}[Exercise 3]
Constraint: $a_{\text{mass}} = a$, $a_{\text{free end}} = 2a$. Tension $T = F = 80\,\mathrm{N}$ (massless free end).
For mass $6\,\mathrm{kg}$: $2T - 6g = 6a \Rightarrow 160 - 58.8 = 6a \Rightarrow a = 16.87\,\mathrm{m/s^2}$ upward.
\end{corrige}

\begin{corrige}[Exercise 4]
Masses: $m$ (table), $2m$ (upper hanging), $3m$ (lower hanging). Strings: $T_1$ between $m$ and $2m$, $T_2$ between $2m$ and $3m$.
Equations:
$m$: $T_1 = ma$.
$2m$: $T_2 + 2mg - T_1 = 2ma$.
$3m$: $3mg - T_2 = 3ma$.
Add: $5mg = 6ma \Rightarrow a = \frac{5g}{6}$.
$T_1 = m\frac{5g}{6} = \frac{5mg}{6}$.
$T_2 = 3mg - 3ma = 3mg - 3m\frac{5g}{6} = \frac{mg}{2}$.
\end{corrige}

\begin{corrige}[Exercise 5]
$m=3\,\mathrm{kg}$, $F=4t+2$. $a = \frac{F}{m} = \frac{4t+2}{3}$.
$dv = \frac{4t+2}{3}dt$. Integrate $0$ to $t$: $v = \frac{1}{3}(2t^2 + 2t) = \frac{2}{3}t(t+1)$.
At $t=4$: $v = \frac{2}{3}\times 4 \times 5 = \frac{40}{3} \approx 13.3\,\mathrm{m/s}$.
$x = \int_0^4 v\,dt = \frac{2}{3}\int_0^4 (t^2+t)\,dt = \frac{2}{3}\left[\frac{t^3}{3}+\frac{t^2}{2}\right]_0^4 = \frac{2}{3}\left(\frac{64}{3}+8\right) = \frac{2}{3}\times \frac{88}{3} = \frac{176}{9} \approx 19.6\,\mathrm{m}$.
\end{corrige}

\coursec{Section 6: Energy Methods and Real-Life Problems}

In the previous section we analysed pulley systems with more than two masses and even cases where the acceleration varies with time.  We solved those problems by writing Newton's second law for each particle, eliminating the tensions, and integrating when necessary.  That approach is powerful because it gives the acceleration and the tension at every instant.  However, when the question only asks for the \emph{speed after a given displacement} (or the distance needed to reach a certain speed), there is a much shorter route: the \textbf{work–energy principle}.  In this section we develop energy methods for connected particles, show how they bypass the internal tensions entirely, and apply them to realistic engineering situations such as lifts, cable cars and towed vehicles.

\courssub{Work, Kinetic Energy and Potential Energy in Connected Systems}

\begin{definition}[Work–Energy Principle for a System of Particles]
For any system of particles, the change in the total kinetic energy equals the work done by \emph{all} forces (external and internal) acting on the system:
\[
\Delta K_{\text{total}} = W_{\text{external}} + W_{\text{internal}}.
\]
If the internal forces are conservative (e.g. gravity within the system) we can move their work to the left‑hand side as a change in potential energy.  For a system subject to gravity and possibly friction, the principle becomes
\[
\Delta K + \Delta U_{\text{grav}} = W_{\text{friction}} + W_{\text{other external}}.
\]
\end{definition}

In a connected system with a \cle{light inextensible string} passing over a \cle{smooth pulley}, the tensions are internal forces.  Because the string is inextensible, the displacement of each particle along the string is the same in magnitude.  The tension on one particle does positive work while the tension on the other does an equal amount of negative work.  Hence:

\begin{propriete}[Tensions Do No Net Work in a Light Inextensible String]
In a connected system where the string is light and inextensible and the pulley is smooth, the total work done by the tension forces is zero.  Tensions merely \emph{transfer} energy between the particles without dissipating or creating it.
\end{propriete}

This property is the key to the energy method: we can write a single energy equation for the \emph{whole system} without ever calculating the tension.

\courssub{Conservation of Mechanical Energy (Smooth Contacts)}

When all contacts are smooth (no friction) and the pulley is smooth, the only external forces doing work are gravity (conservative) and possibly a driving force.  If no non‑conservative external forces act, mechanical energy is conserved.

\begin{aretenir}[Energy Conservation for a Two‑Mass Atwood‑Type System]
Consider two masses $M$ and $m$ ($M>m$) connected by a light inextensible string over a smooth pulley.  If the system is released from rest and the heavier mass descends a distance $h$, then:
\begin{itemize}
\item Loss of gravitational PE of $M$: $Mgh$.
\item Gain of gravitational PE of $m$: $mgh$.
\item Gain of kinetic energy of \emph{both} masses: $\frac{1}{2}(M+m)v^2$.
\end{itemize}
Conservation of mechanical energy gives:
\[
Mgh = mgh + \frac{1}{2}(M+m)v^2 \quad\Longrightarrow\quad (M-m)gh = \frac{1}{2}(M+m)v^2.
\]
\end{aretenir}

\begin{popfigure}
\begin{tikzpicture}[scale=1.1, every node/.style={font=\small}]
% Left mass M (descending)
\draw[popblue, thick] (-2.5,2) rectangle (-1.5,1) node[midway] {$M$};
\draw[popblue, ->, thick] (-2,1) -- (-2,0) node[midway, right] {$h$};
\draw[popblue, thick] (-2.5,2) -- (-2.5,2.5) -- (-1.5,2.5) -- (-1.5,2);
% Right mass m (rising)
\draw[poporange, thick] (1.5,0) rectangle (2.5,1) node[midway] {$m$};
\draw[poporange, ->, thick] (2,1) -- (2,2) node[midway, right] {$h$};
\draw[poporange, thick] (1.5,0) -- (1.5,-0.5) -- (2.5,-0.5) -- (2.5,0);
% Pulley
\draw[popdark, thick] (0,2.5) circle (0.5);
\draw[popdark, thick] (0,2.5) -- (0,3.2) node[above] {Smooth pulley};
% String
\draw[popmuted, thick] (-2,1.5) .. controls (-2,2) and (-0.5,3) .. (0,3) .. controls (0.5,3) and (2,2) .. (2,1.5);
% Energy flow arrows
\draw[popgreen, ->, thick] (-2.5,2.7) -- (-0.5,2.7) node[midway, above] {PE lost by $M$};
\draw[popgreen, ->, thick] (0.5,2.7) -- (2.5,2.7) node[midway, above] {PE gained by $m$};
\draw[popred, ->, thick] (-2.5,0.5) -- (-0.5,0.5) node[midway, below] {KE gained by $M$};
\draw[popred, ->, thick] (0.5,0.5) -- (2.5,0.5) node[midway, below] {KE gained by $m$};
% Labels
\node[popblue] at (-2,3.2) {$Mgh$};
\node[poporange] at (2,3.2) {$mgh$};
\node[popred] at (-2,0) {$\frac12 Mv^2$};
\node[popred] at (2,0) {$\frac12 mv^2$};
\end{tikzpicture}
\legende{Energy‑flow diagram for an Atwood's machine: the potential energy lost by the descending mass $M$ is converted into kinetic energy of both masses and potential energy gained by the rising mass $m$.}
\end{popfigure}

\begin{attention}[Common Mistake: Forgetting the Kinetic Energy of the Rising Mass]
In a connected system \emph{every moving particle} carries kinetic energy.  A frequent error is to write $Mgh = \frac{1}{2}Mv^2$, ignoring the fact that the lighter mass $m$ also moves and therefore has kinetic energy $\frac{1}{2}mv^2$.  Always include the KE of \emph{all} particles that are in motion.
\end{attention}

\courssub{Work–Energy Principle with Friction (Rough Surfaces)}

When one or more particles move on rough surfaces, friction does negative work.  Mechanical energy is no longer conserved; we must use the full work–energy principle.

\begin{propriete}[Work–Energy Principle with Friction]
For a connected system where a particle of mass $m$ moves a distance $s$ on a rough horizontal surface (coefficient of friction $\mu$), the work done against friction is $W_{\text{fric}} = -\mu mg s$.  The energy equation becomes:
\[
\text{(Loss of PE of descending masses)} = \text{(Gain of KE of all masses)} + \text{(Gain of PE of rising masses)} + \mu mg s.
\]
Equivalently:
\[
\Delta K + \Delta U = W_{\text{friction}} \quad (\text{where } W_{\text{friction}} < 0).
\]
\end{propriete}

\begin{methode}[Energy Route for Connected Particles]
\begin{enumerate}
\item \textbf{Identify the initial and final states.}  Choose a reference level for gravitational potential energy (usually the lowest point reached by any particle).
\item \textbf{Compute the total change in gravitational potential energy $\Delta U$ for \emph{all} particles.}  $\Delta U = \sum m_i g \Delta h_i$ (positive if the particle rises).
\item \textbf{Compute the total change in kinetic energy $\Delta K$ for \emph{all} particles.}  If starting from rest, $\Delta K = \frac{1}{2}\sum m_i v^2$.
\item \textbf{Account for work done against friction.}  For each particle on a rough surface, $W_{\text{fric}} = -\mu R \times \text{distance moved}$.
\item \textbf{Write the work–energy equation:} $\Delta K = -\Delta U + W_{\text{fric}} + W_{\text{other external}}$.
\item \textbf{Solve for the unknown} (usually $v$ or the distance $s$).
\end{enumerate}
\end{methode}

\begin{exoresolu}[Block on Rough Table Connected to Hanging Mass]
A block of mass $4\,\mathrm{kg}$ rests on a rough horizontal table ($\mu = 0.2$).  It is connected by a light inextensible string passing over a smooth pulley at the edge of the table to a hanging mass of $2\,\mathrm{kg}$.  The system is released from rest.  Find the speed of the masses after the hanging mass has descended $1.5\,\mathrm{m}$.
\tcblower
\textbf{Solution by Energy Method.}
\begin{itemize}
\item Initial state: system at rest, hanging mass at reference level (PE = 0).  Block on table: PE constant (horizontal motion).
\item Final state: hanging mass descended $h = 1.5\,\mathrm{m}$ (PE lost = $mgh = 2 \times 9.8 \times 1.5 = 29.4\,\mathrm{J}$).  Block moved $1.5\,\mathrm{m}$ on table.
\item Friction on block: $R = 4g = 39.2\,\mathrm{N}$, $F_{\text{fric}} = \mu R = 0.2 \times 39.2 = 7.84\,\mathrm{N}$.  Work done against friction $W_{\text{fric}} = -7.84 \times 1.5 = -11.76\,\mathrm{J}$.
\item Both masses move with same speed $v$.  Total KE gain = $\frac{1}{2}(4+2)v^2 = 3v^2$.
\item Energy equation: $\Delta K = -\Delta U + W_{\text{fric}}$.
  \[
  3v^2 = 29.4 - 11.76 = 17.64 \quad\Rightarrow\quad v^2 = 5.88 \quad\Rightarrow\quad v = \sqrt{5.88} \approx 2.42\,\mathrm{m\,s^{-1}}.
  \]
\end{itemize}
\textbf{Check with Newton's Laws (optional).}
For hanging mass: $2g - T = 2a$.  For block: $T - \mu(4g) = 4a$.  Adding: $2g - 0.2(4g) = 6a \Rightarrow a = \frac{19.6 - 7.84}{6} = 1.96\,\mathrm{m\,s^{-2}}$.  Then $v^2 = 2as = 2 \times 1.96 \times 1.5 = 5.88$, same result.  The energy method avoided finding $T$ entirely.
\end{exoresolu}

\courssub{Comparing Newton's Laws and Energy Methods}

\begin{center}
\begin{tabularx}{\linewidth}{|l|X|X|}\hline
\textbf{Aspect} & \textbf{Newton's Second Law ($F=ma$)} & \textbf{Work–Energy Principle} \\ \hline
\textbf{Unknowns} & Gives acceleration $a$ and tension $T$ at every instant. & Gives speed $v$ after a given displacement $s$ (or $s$ for a given $v$). \\ \hline
\textbf{Forces needed} & Must identify \emph{all} forces on \emph{each} particle, including internal tensions. & Only \emph{external} forces that do work (gravity, friction, driving forces). Tensions cancel out. \\ \hline
\textbf{Variable acceleration} & Handles naturally via $a = v\frac{dv}{dx}$ or integration. & Handles naturally if work integrals can be evaluated. \\ \hline
\textbf{When to prefer} & \begin{itemize}\item Need tension or acceleration at a specific instant.\item Forces change direction (e.g. particle leaves surface).\item Motion on curved paths where normal reaction does no work but affects constraints.\end{itemize} & \begin{itemize}\item Only speed after a given displacement is asked.\item System has many particles; writing many $F=ma$ equations is tedious.\item Friction is present but constant (work = force $\times$ distance).\end{itemize} \\ \hline
\end{tabularx}
\end{center}

\begin{aretenir}[Rule of Thumb]
\begin{itemize}
\item \textbf{``Find the acceleration / tension''} $\rightarrow$ Newton's laws.
\item \textbf{``Find the speed after moving $x$ metres''} $\rightarrow$ Energy method (usually faster).
\item \textbf{``Find the distance to come to rest''} $\rightarrow$ Energy method.
\end{itemize}
\end{aretenir}

\courssub{Power in Connected Systems (GCE Extension)}

Power is the rate at which work is done.  For a constant force $F$ acting on a particle moving with velocity $v$ in the direction of the force, $P = Fv$.  In a connected system, the power delivered by a driving force (e.g. a motor pulling a cable) equals the rate of increase of the total mechanical energy of the system plus the rate of working against friction.

\begin{exemplebox}[Power of a Winch]
A winch pulls a loaded cage of mass $500\,\mathrm{kg}$ up a rough inclined plane at constant speed $2\,\mathrm{m\,s^{-1}}$.  The plane is at $30^\circ$ to the horizontal and $\mu = 0.15$.  Find the power output of the winch.
\tcblower
At constant speed, $\Delta K = 0$.  In $1\,\mathrm{s}$ the cage moves $2\,\mathrm{m}$ up the plane.
\begin{itemize}
\item Gain in PE: $mgh = 500 \times 9.8 \times (2\sin30^\circ) = 4900\,\mathrm{J}$.
\item Work against friction: $F_{\text{fric}} = \mu mg\cos30^\circ = 0.15 \times 500 \times 9.8 \times \frac{\sqrt{3}}{2} \approx 636.5\,\mathrm{N}$.  Work in $1\,\mathrm{s}$: $636.5 \times 2 = 1273\,\mathrm{J}$.
\item Total work per second (power) $P = 4900 + 1273 = 6173\,\mathrm{W} \approx 6.17\,\mathrm{kW}$.
\end{itemize}
Alternatively, tension in cable $T = mg\sin30^\circ + \mu mg\cos30^\circ \approx 3086.5\,\mathrm{N}$, then $P = Tv = 3086.5 \times 2 = 6173\,\mathrm{W}$.
\end{exemplebox}

\courssub{Real-Life Modelling Problems}

The GCE examination often presents a practical situation (a lift, a cable car, a towed vehicle) and expects you to go through the full \textbf{modelling cycle}.

\begin{methode}[Modelling Cycle for Real-Life Problems]
\begin{enumerate}
\item \textbf{State assumptions explicitly.}  Typical assumptions: light inextensible strings/cables, smooth pulleys, particles (no rotation), constant $g = 9.8\,\mathrm{m\,s^{-2}}$, no air resistance, constant coefficient of friction.
\item \textbf{Draw a clear diagram} showing all forces, displacements, and coordinate directions.
\item \textbf{Formulate equations} using either Newton's laws or energy principles (choose the most efficient).
\item \textbf{Solve mathematically} for the required quantities.
\item \textbf{Interpret the answer} in the context of the problem (units, direction, physical plausibility).
\item \textbf{Critique the model:} How would the answer change if we relaxed an assumption?  (e.g. massive cable $\rightarrow$ tension varies along cable; pulley friction $\rightarrow$ tensions on either side differ; air resistance $\rightarrow$ terminal velocity).
\end{enumerate}
\end{methode}

\begin{exoresolu}[Lift with Counterweight — A Classic Application]
A lift of mass $1200\,\mathrm{kg}$ is connected by a light inextensible cable passing over a smooth pulley to a counterweight of mass $950\,\mathrm{kg}$.  The lift carries passengers of total mass $350\,\mathrm{kg}$.  The system starts from rest.  Find:
\begin{enumerate}
\item The acceleration of the lift.
\item The tension in the cable.
\item The speed of the lift after it has moved $10\,\mathrm{m}$.
\item The power required from the motor to maintain a constant upward speed of $2\,\mathrm{m\,s^{-1}}$ when the lift is fully loaded (assume the motor drives the pulley).
\end{enumerate}
\tcblower
\textbf{Assumptions:} Cable light and inextensible, pulley smooth, lift and counterweight move vertically as particles, $g=9.8\,\mathrm{m\,s^{-2}}$.

\textbf{Diagram:} Lift (total mass $M = 1200+350 = 1550\,\mathrm{kg}$) on left, counterweight $m = 950\,\mathrm{kg}$ on right.  Lift moves up, counterweight moves down.

\textbf{(i) Acceleration (Newton's Laws).}
Let upward be positive for \emph{both} masses.
Lift: $T - Mg = M a$.
Counterweight: $T - mg = m(-a) \Rightarrow mg - T = ma$.
Add: $(m-M)g = (M+m)a \;\Rightarrow\; a = \frac{(950-1550)\times9.8}{2500} = -2.352\,\mathrm{m\,s^{-2}}$.
The negative sign means the lift accelerates \emph{downwards} (the counterweight is lighter than the loaded lift).  Magnitude $a = 2.35\,\mathrm{m\,s^{-2}}$ downwards.

\textbf{(ii) Tension.}
$T = M(g+a) = 1550(9.8 - 2.352) = 1550 \times 7.448 = 11544.4\,\mathrm{N} \approx 11.5\,\mathrm{kN}$.
(Check with counterweight: $T = m(g-a) = 950(9.8 + 2.352) = 950 \times 12.152 = 11544.4\,\mathrm{N}$.  Consistent.)

\textbf{(iii) Speed after $10\,\mathrm{m}$ (Energy Method).}
Lift moves down $10\,\mathrm{m}$, counterweight moves up $10\,\mathrm{m}$.
Loss of PE of lift: $Mgh = 1550 \times 9.8 \times 10 = 151900\,\mathrm{J}$.
Gain of PE of counterweight: $mgh = 950 \times 9.8 \times 10 = 93100\,\mathrm{J}$.
Net loss of PE = $151900 - 93100 = 58800\,\mathrm{J}$.
This becomes KE of both: $\frac{1}{2}(M+m)v^2 = \frac{1}{2}(2500)v^2 = 1250 v^2$.
$1250 v^2 = 58800 \Rightarrow v^2 = 47.04 \Rightarrow v = 6.86\,\mathrm{m\,s^{-1}}$.

\textbf{(iv) Power for Constant Upward Speed $2\,\mathrm{m\,s^{-1}}$.}
At constant speed, $a=0$.  The motor must overcome the weight difference.
Net force the motor must overcome: $Mg - mg = (1550-950)\times9.8 = 5880\,\mathrm{N}$.
Power $P = F v = 5880 \times 2 = 11760\,\mathrm{W} = 11.76\,\mathrm{kW}$.
(Alternatively: rate of gain of PE of lift $= Mgv = 15190 \times 2 = 30380\,\mathrm{W}$; rate of loss of PE of counterweight $= mgv = 9310 \times 2 = 18620\,\mathrm{W}$; net power needed $= 30380 - 18620 = 11760\,\mathrm{W}$.)

\textbf{Critique of Model:}
\begin{itemize}
\item If the cable has mass, tension varies along its length; the motor would need extra power to lift the cable itself.
\item If the pulley has friction, the tensions on either side differ even without a motor; the motor must overcome both weight difference and pulley friction.
\item Air resistance on the lift would increase the required power at high speeds.
\end{itemize}
\end{exoresolu}

\begin{popfigure}
\begin{tikzpicture}[scale=0.9, every node/.style={font=\small}]
% Lift shaft
\draw[popdark, thick] (-3,0) -- (-3,6);
\draw[popdark, thick] (-1,0) -- (-1,6);
% Counterweight shaft
\draw[popdark, thick] (3,0) -- (3,6);
\draw[popdark, thick] (5,0) -- (5,6);
% Pulley at top
\draw[popdark, thick] (-2,6) circle (0.6);
\draw[popdark, thick] (4,6) circle (0.6);
% Cable
\draw[popmuted, thick] (-2,5.4) -- (-2,6.6) .. controls (-2,7) and (4,7) .. (4,6.6) -- (4,5.4);
% Lift
\draw[popblue, thick, fill=popblueL] (-2.8,2) rectangle (-1.2,3.5) node[midway, align=center] {Lift\\$M=1550\,\mathrm{kg}$};
\draw[popblue, ->, thick] (-2,2) -- (-2,1) node[midway, right] {$10\,\mathrm{m}$};
% Counterweight
\draw[poporange, thick, fill=poporange!30] (3.2,4) rectangle (4.8,5) node[midway, align=center] {Counterweight\\$m=950\,\mathrm{kg}$};
\draw[poporange, ->, thick] (4,5) -- (4,6) node[midway, right] {$10\,\mathrm{m}$};
% Forces
\draw[popred, ->, thick] (-2,3.5) -- (-2,4.5) node[midway, right] {$T$};
\draw[popred, ->, thick] (-2,2) -- (-2,1.2) node[midway, right] {$Mg$};
\draw[popred, ->, thick] (4,4) -- (4,3.2) node[midway, right] {$T$};
\draw[popred, ->, thick] (4,5) -- (4,5.8) node[midway, right] {$mg$};
% Motor
\draw[popgreen, thick, fill=popgreenL] (0,6.5) rectangle (2,7.5) node[midway, align=center] {Motor\\$P=11.76\,\mathrm{kW}$};
\draw[popgreen, ->, thick] (1,6.5) -- (1,6.2);
\end{tikzpicture}
\legende{Lift with counterweight: forces, displacements and motor power.  The motor supplies the power to overcome the weight difference when moving at constant speed.}
\end{popfigure}

\courssub{Variable Acceleration and Energy: $v^2$ vs Displacement Graph}

For a system with constant acceleration (like the Atwood machine with constant masses), the energy equation $(M-m)gh = \frac{1}{2}(M+m)v^2$ gives $v^2 = \frac{2(M-m)g}{M+m}\,h$.  This is a straight line through the origin when $v^2$ is plotted against $h$.

\begin{popfigure}
\begin{tikzpicture}
\begin{axis}[
    width=10cm, height=6cm,
    xlabel={Displacement $h$ (m)},
    ylabel={$v^2$ ($\mathrm{m^2\,s^{-2}}$)},
    xmin=0, xmax=5,
    ymin=0, ymax=50,
    grid=major,
    grid style={popline},
    axis lines=middle,
    tick label style={font=\small},
    label style={font=\small},
]
\addplot[popblue, thick, domain=0:5, samples=2] {4.704*x};
\addlegendentry{$v^2 = 2ah$ (constant $a$)}
\addplot[poporange, thick, domain=0:5, samples=50] {4.704*x + 0.5*x^2};
\addlegendentry{Variable acceleration (e.g. mass leaking)}
\end{axis}
\end{tikzpicture}
\legende{Graph of $v^2$ against displacement $h$.  For constant acceleration (smooth Atwood machine) the graph is a straight line through the origin with slope $2a$.  Curvature indicates variable acceleration (e.g. mass changing, friction varying).}
\end{popfigure}

\begin{attention}[Common Mistake: Using Conservation of Energy When Friction Acts]
If any non‑conservative force (friction, air resistance, a motor doing work) acts on the system, \textbf{mechanical energy is not conserved}.  You must use the work–energy principle:
\[
\Delta K + \Delta U = W_{\text{non-conservative}}.
\]
Writing ``Loss of PE = Gain of KE'' when friction is present will give an answer that is too large (it predicts a higher speed than reality).
\end{attention}

\begin{exoresolu}[Towed Vehicle on a Hill]
A car of mass $1200\,\mathrm{kg}$ tows a trailer of mass $800\,\mathrm{kg}$ up a hill inclined at $\sin^{-1}(0.1)$ to the horizontal.  The resistance to motion (constant) is $400\,\mathrm{N}$ on the car and $250\,\mathrm{N}$ on the trailer.  The driving force of the car's engine is $5000\,\mathrm{N}$.  The tow‑bar is light and inextensible.  Find:
\begin{enumerate}
\item The acceleration of the system.
\item The tension in the tow‑bar.
\item The distance travelled from rest to reach a speed of $15\,\mathrm{m\,s^{-1}}$.
\end{enumerate}
\tcblower
\textbf{Assumptions:} Car and trailer as particles, tow‑bar light and inextensible, constant resistances, $g=9.8$.

\textbf{Diagram:} Car (front) pulling trailer.  Incline angle $\theta$, $\sin\theta=0.1$, $\cos\theta=\sqrt{0.99}\approx0.995$.

\textbf{(i) Acceleration (Newton's Laws — needed for tension).}
System as a whole (mass $2000\,\mathrm{kg}$):
Driving force $-$ total resistance $-$ component of weight down plane $= Ma$.
$5000 - (400+250) - 2000g\sin\theta = 2000a$.
$5000 - 650 - 2000\times9.8\times0.1 = 2000a$.
$5000 - 650 - 1960 = 2000a \Rightarrow 2390 = 2000a \Rightarrow a = 1.195\,\mathrm{m\,s^{-2}}$.

\textbf{(ii) Tension in tow‑bar.}
Consider trailer only (mass $800\,\mathrm{kg}$):
$T - 250 - 800g\sin\theta = 800a$.
$T = 250 + 800\times9.8\times0.1 + 800\times1.195 = 250 + 784 + 956 = 1990\,\mathrm{N}$.

\textbf{(iii) Distance to reach $15\,\mathrm{m\,s^{-1}}$ (Energy Method — faster).}
Work done by engine $= 5000 \times s$.
Work against resistances $= -(400+250)s = -650s$.
Gain in PE $= (1200+800)g (s\sin\theta) = 2000 \times 9.8 \times 0.1 s = 1960s$.
Gain in KE $= \frac{1}{2}(2000)(15^2) = 225000\,\mathrm{J}$.
Work–energy: $\Delta K = W_{\text{engine}} + W_{\text{resist}} + W_{\text{weight}}$.
$225000 = 5000s - 650s - 1960s = 2390s$.
$s = \frac{225000}{2390} \approx 94.1\,\mathrm{m}$.
(Check: $v^2 = 2as = 2\times1.195\times94.1 \approx 225$, correct.)
\end{exoresolu}

\courssub{Further Real-Life Examples}

\begin{exemplebox}[Cable Car]
A cable car of mass $2000\,\mathrm{kg}$ carries $20$ passengers (average mass $70\,\mathrm{kg}$).  It is pulled up a slope of $30^\circ$ by a cable driven by a motor at the top.  The cable is light and inextensible.  The resistance to motion is $1500\,\mathrm{N}$.  Find the power required to maintain a constant speed of $6\,\mathrm{m\,s^{-1}}$.
\tcblower
Total mass $M = 2000 + 20\times70 = 3400\,\mathrm{kg}$.
At constant speed, tension $T = Mg\sin30^\circ + 1500 = 3400\times9.8\times0.5 + 1500 = 16660 + 1500 = 18160\,\mathrm{N}$.
Power $P = Tv = 18160 \times 6 = 108960\,\mathrm{W} \approx 109\,\mathrm{kW}$.
\end{exemplebox}

\begin{exemplebox}[Bucket Raised from a Well — Variable Mass]
A bucket of mass $5\,\mathrm{kg}$ (including water) is raised vertically from a well $20\,\mathrm{m}$ deep by a rope wound on a winch.  The rope is light and inextensible.  Water leaks out at a constant rate so that the bucket is empty at the top.  The mass of the empty bucket is $2\,\mathrm{kg}$.  Find the work done by the winch if the bucket is raised at constant speed.
\tcblower
Mass varies linearly from $5\,\mathrm{kg}$ to $2\,\mathrm{kg}$ over $20\,\mathrm{m}$.
Average mass $= \frac{5+2}{2} = 3.5\,\mathrm{kg}$.
Work done against gravity $= (\text{average weight}) \times \text{distance} = 3.5g \times 20 = 3.5 \times 9.8 \times 20 = 686\,\mathrm{J}$.
(Alternatively, integrate: $m(x) = 5 - \frac{3}{20}x$, $W = \int_0^{20} m(x)g\,dx = g[5x - \frac{3}{40}x^2]_0^{20} = 9.8(100 - 30) = 686\,\mathrm{J}$.)
This is an example of \cle{variable acceleration} (if speed were not constant) handled by energy integration.
\end{exemplebox}

\begin{exemplebox}[Load Sliding Down a Rough Chute]
A package of mass $10\,\mathrm{kg}$ slides down a rough chute inclined at $25^\circ$ to the horizontal.  The coefficient of friction is $0.3$.  The package starts from rest at the top and reaches the bottom with speed $8\,\mathrm{m\,s^{-1}}$.  Find the length of the chute.
\tcblower
Energy method: Loss of PE $=$ Gain in KE $+$ Work against friction.
$mgh = \frac{1}{2}mv^2 + \mu mg\cos\theta \times s$, where $h = s\sin\theta$.
$mg s\sin\theta = \frac{1}{2}mv^2 + \mu mg\cos\theta\, s$.
Cancel $m$: $g s\sin\theta = \frac{1}{2}v^2 + \mu g\cos\theta\, s$.
$s(g\sin\theta - \mu g\cos\theta) = \frac{1}{2}v^2$.
$s = \frac{v^2}{2g(\sin\theta - \mu\cos\theta)} = \frac{64}{2\times9.8(\sin25^\circ - 0.3\cos25^\circ)}$.
$\sin25^\circ\approx0.423$, $\cos25^\circ\approx0.906$.
Denominator $= 19.6(0.423 - 0.272) = 19.6\times0.151 = 2.96$.
$s \approx 21.6\,\mathrm{m}$.
\end{exemplebox}

\courssub{Summary and Examination Tips}

\begin{aretenir}[Key Points for the GCE Examination]
\begin{itemize}
\item \textbf{Energy conservation} applies only when \emph{all} contacts are smooth and no external non‑conservative forces do work.  Equation: $\Delta K + \Delta U = 0$.
\item \textbf{Work–energy principle} is the general tool: $\Delta K + \Delta U = W_{\text{friction}} + W_{\text{driving}}$.
\item \textbf{Tensions in a light inextensible string do no net work} — they are internal forces that cancel out in the system energy equation.
\item \textbf{Always include the kinetic energy of \emph{every} moving particle.}
\item \textbf{Choose the method wisely:} Newton's laws for acceleration/tension; energy for speed/distance.
\item \textbf{Modelling cycle:} Assumptions $\rightarrow$ Diagram $\rightarrow$ Equations $\rightarrow$ Solution $\rightarrow$ Interpretation $\rightarrow$ Critique.
\item \textbf{Power} $P = Fv$ for constant force and velocity; for variable force, $P = \frac{dW}{dt}$.
\end{itemize}
\end{aretenir}

\begin{exercice}[Practice Problems]
\begin{itemize}
\item Two masses $3\,\mathrm{kg}$ and $5\,\mathrm{kg}$ are connected by a light inextensible string over a smooth pulley.  The $5\,\mathrm{kg}$ mass is held at rest on a rough horizontal table ($\mu=0.4$) and the $3\,\mathrm{kg}$ mass hangs freely.  The system is released.  Use energy methods to find the speed of the masses after the hanging mass has fallen $1.2\,\mathrm{m}$.
\item A train of mass $400\,\mathrm{tonnes}$ climbs a slope of $1$ in $100$ at a constant speed of $20\,\mathrm{m\,s^{-1}}$.  The resistance to motion is $50\,\mathrm{N}$ per tonne.  Find the power of the engine in kW.
\item A skier of mass $75\,\mathrm{kg}$ starts from rest at the top of a rough slope of length $200\,\mathrm{m}$ inclined at $15^\circ$.  The coefficient of friction is $0.08$.  Find the skier's speed at the bottom.  If the skier then travels on horizontal snow with the same coefficient of friction, how far does she travel before stopping?
\item A lift of mass $800\,\mathrm{kg}$ is counterbalanced by a mass of $600\,\mathrm{kg}$.  The lift carries $10$ people of average mass $70\,\mathrm{kg}$.  The system starts from rest.  (a) Find the acceleration. (b) Find the tension in the cable. (c) Use energy to find the speed after the lift has descended $15\,\mathrm{m}$. (d) What power must the motor provide to raise the loaded lift at a constant $2\,\mathrm{m\,s^{-1}}$?
\end{itemize}
\end{exercice}

\begin{corrige}[Exercise 1]
Loss of PE of $3\,\mathrm{kg}$ mass $= 3g(1.2) = 35.28\,\mathrm{J}$.
Work against friction on $5\,\mathrm{kg}$ mass $= \mu mg s = 0.4 \times 5 \times 9.8 \times 1.2 = 23.52\,\mathrm{J}$.
Gain in KE of both masses $= \frac{1}{2}(3+5)v^2 = 4v^2$.
Energy: $35.28 = 4v^2 + 23.52 \Rightarrow 4v^2 = 11.76 \Rightarrow v^2 = 2.94 \Rightarrow v = 1.71\,\mathrm{m\,s^{-1}}$.
\end{corrige}

\begin{corrige}[Exercise 2]
Mass $= 400\,000\,\mathrm{kg}$.  Slope $\sin\theta = 1/100 = 0.01$.
Component of weight down slope $= Mg\sin\theta = 400\,000 \times 9.8 \times 0.01 = 39\,200\,\mathrm{N}$.
Resistance $= 50 \times 400 = 20\,000\,\mathrm{N}$.
Total force to overcome at constant speed $= 39\,200 + 20\,000 = 59\,200\,\mathrm{N}$.
Power $P = Fv = 59\,200 \times 20 = 1\,184\,000\,\mathrm{W} = 1184\,\mathrm{kW}$.
\end{corrige}

\begin{corrige}[Exercise 3]
Slope: loss of PE $= mgh = 75 \times 9.8 \times (200\sin15^\circ) = 75 \times 9.8 \times 51.76 = 38\,041\,\mathrm{J}$.
Work against friction on slope $= \mu mg\cos15^\circ \times 200 = 0.08 \times 75 \times 9.8 \times 0.966 \times 200 = 11\,370\,\mathrm{J}$.
KE at bottom $= 38\,041 - 11\,370 = 26\,671\,\mathrm{J}$.
$v = \sqrt{2\times26\,671/75} = 26.7\,\mathrm{m\,s^{-1}}$.
Horizontal: initial KE $= 26\,671\,\mathrm{J}$.  Work against friction $= \mu mg d = 0.08 \times 75 \times 9.8 \times d = 58.8 d$.
$58.8 d = 26\,671 \Rightarrow d = 453\,\mathrm{m}$.
\end{corrige}

\begin{corrige}[Exercise 4]
(a) Total lift mass $M = 800 + 700 = 1500\,\mathrm{kg}$.  Counterweight $m = 600\,\mathrm{kg}$.
$a = \frac{(M-m)g}{M+m} = \frac{900 \times 9.8}{2100} = 4.2\,\mathrm{m\,s^{-2}}$ (lift down).
(b) $T = M(g-a) = 1500(9.8-4.2) = 8400\,\mathrm{N}$.
(c) Lift descends $15\,\mathrm{m}$, counterweight rises $15\,\mathrm{m}$.
Net PE loss $= (M-m)gh = 900 \times 9.8 \times 15 = 132\,300\,\mathrm{J}$.
KE gain $= \frac{1}{2}(2100)v^2 = 1050 v^2$.
$1050 v^2 = 132\,300 \Rightarrow v^2 = 126 \Rightarrow v = 11.2\,\mathrm{m\,s^{-1}}$.
(d) Constant speed up: motor must overcome weight difference $= (M-m)g = 8820\,\mathrm{N}$.
$P = Fv = 8820 \times 2 = 17\,640\,\mathrm{W} = 17.6\,\mathrm{kW}$.
\end{corrige}

\begin{savaistu}[Historical Note: The Atwood Machine]
The Atwood machine was invented in 1784 by the English mathematician \textbf{George Atwood} (1745–1807) as a laboratory device to verify Newton's laws of motion with reduced acceleration (so that motion could be timed accurately with the clocks of the era).  By choosing two nearly equal masses, the acceleration $a = \frac{M-m}{M+m}g$ can be made very small, allowing precise measurement of $g$.  Modern physics laboratories still use a version of Atwood's machine, often with a photogate timer, to teach dynamics and energy conservation.
\end{savaistu}

\begin{savaistu}[Cameroon Context: The Edea–Douala Cable Car Project]
Feasibility studies have been conducted for a cable car system linking Douala and Edea to ease traffic congestion on the national highway.  Such a system would involve cabins of mass $\sim 2000\,\mathrm{kg}$ carrying up to $30$ passengers, moving on cables with a driving motor at one station.  The power calculations we have just performed are exactly what mechanical engineers carry out to size the motors and cables, taking into account the weight difference between ascending and descending cabins, friction in the pulleys, wind loads, and safety factors.  The ``modelling cycle'' you learn here is the same cycle used in real engineering design offices in Yaoundé and Douala.
\end{savaistu}

\newpage

\coursec{Integration activity}

\courssub{Aims of the activity}

This integration activity closes the chapter by bringing together, in one realistic modelling task, every skill developed in Sections 1 to 6:

\begin{itemize}
\item stating and justifying \cle{modelling assumptions} (light inextensible rope, smooth pulley, particle model);
\item drawing correct \cle{free-body diagrams} for each body of a connected system;
\item applying \cle{Newton's second law} to each particle and solving the resulting simultaneous equations for the \cle{acceleration} and the \cle{tension};
\item using \cle{suvat equations} for constant-acceleration stages;
\item analysing \cle{multi-stage motion} (what happens when one body suddenly stops or the rope goes slack);
\item applying the \cle{work--energy principle} as an independent check of kinematic results;
\item resolving forces on a \cle{rough inclined plane} and deciding whether motion is possible;
\item \cle{critiquing a model}: discussing the effect of rope mass, pulley friction and impulsive jerks.
\end{itemize}

By the end, each group will have carried out a complete \emph{modelling cycle}: real situation $\rightarrow$ assumptions $\rightarrow$ equations $\rightarrow$ solution $\rightarrow$ interpretation $\rightarrow$ critique.

\courssub{Situation}

\textbf{The Limbe Engine Hoist.}\par

Mr.~Ekane runs a small vehicle-repair workshop in Limbe, at the foot of Mount Cameroon. He must remove a car engine of mass $M = 120$~kg and lower it gently to the workshop floor. Having no crane, he passes a strong rope over a smooth horizontal steel beam fixed \SI{4.0}{m} above the floor. To one end of the rope he attaches the engine; to the other end he attaches a \cle{counterweight} made of scrap metal of mass $m = 90$~kg. He releases the system from rest, with the engine initially \SI{1.5}{m} above the floor.

The rope is rated by its manufacturer for a maximum safe tension of \SI{1500}{N}. Take $g = \SI{9.8}{m.s^{-2}}$ throughout.

\begin{popfigure}
\begin{center}
\begin{tikzpicture}[scale=0.9, every node/.style={font=\small}]
% beam
\draw[line width=2.5pt, popdark] (-2.2,3.4) -- (2.2,3.4);
\foreach \x in {-2.0,-1.6,...,2.0}{\draw[popdark] (\x,3.4) -- (\x+0.25,3.65);}
% pulley
\draw[fill=popblueL, draw=popblue, line width=1pt] (0,3.0) circle (0.4);
\fill[popblue] (0,3.0) circle (0.05);
% ropes
\draw[line width=1.2pt, popink] (-0.4,3.0) -- (-0.4,1.2);
\draw[line width=1.2pt, popink] (0.4,3.0) -- (0.4,0.6);
% engine
\draw[fill=poporange!70, draw=popdark] (-0.95,0.5) rectangle (0.15,1.2);
\node at (-0.4,0.85) {$M$};
\node[below] at (-0.4,0.45) {engine $120$ kg};
% counterweight
\draw[fill=popgreenL, draw=popgreen!60!black] (0.0,0.0) rectangle (0.8,0.6);
\node at (0.4,0.3) {$m$};
\node[below] at (0.4,-0.05) {scrap $90$ kg};
% floor
\draw[line width=1.5pt, popdark] (-2.2,-0.6) -- (2.2,-0.6);
\foreach \x in {-2.1,-1.7,...,2.1}{\draw[popdark] (\x,-0.6) -- (\x-0.2,-0.85);}
% height marker
\draw[<->, poppurple] (-1.6,-0.6) -- (-1.6,0.5) node[midway, left]{\SI{1.5}{m}};
\end{tikzpicture}
\end{center}
\legende{The engine hoist: engine $M$ and scrap counterweight $m$ connected by a rope passing over a smooth beam.}
\end{popfigure}

In a \textbf{second stage} of the job, Mr.~Ekane wonders whether the same counterweight could later drag the engine back \emph{up} a rough concrete ramp inclined at $20^\circ$ to the horizontal, for which the coefficient of friction between engine and ramp is $\mu = 0.3$ (the rope still passes over the smooth beam, with the counterweight hanging vertically).

Working in groups, you are the consulting team. Produce a full written report answering the tasks below, ending with a critical discussion of the model.

\courssub{Tasks}

\begin{enumerate}
\item \textbf{Modelling.} State at least four modelling assumptions needed to treat this as a standard connected-particles problem. Draw the free-body diagram of the engine and that of the counterweight, showing all forces.
\item \textbf{First stage of motion.} Taking the engine and the counterweight separately, apply Newton's second law and hence find:
\begin{enumerate}
\item the acceleration of the system;
\item the tension $T$ in the rope.
\end{enumerate}
Is the rope safe? Justify with an inequality.
\item \textbf{Speed after \SI{1.5}{m}.} The engine descends \SI{1.5}{m} from rest before touching the floor.
\begin{enumerate}
\item Find its speed on reaching the floor using a suvat equation.
\item Find the same speed using the work--energy principle applied to the \emph{whole system} (loss of gravitational potential energy $=$ gain of kinetic energy). Compare the two answers.
\end{enumerate}
\item \textbf{Second stage: rope goes slack.} Once the engine rests on the floor, the counterweight is still moving upwards. Explain why the rope becomes slack, state the counterweight's subsequent acceleration, and find the \emph{extra} height it rises before momentary rest. Hence find its greatest height above its position at the moment of release.
\item \textbf{Extension: the rough ramp.} For the return job, the engine lies on a ramp inclined at $20^\circ$ ($\mu = 0.3$) and is attached to the same $90$~kg counterweight hanging vertically over the beam.
\begin{enumerate}
\item Determine, with full working, whether the $90$~kg counterweight can pull the engine up the ramp; if it can, find the acceleration of the system.
\item Find the minimum counterweight mass $m_{\min}$ that would just set the engine moving up the ramp.
\end{enumerate}
\item \textbf{Critique.} Discuss how each of the following would change your numerical answers: the rope has mass; the beam is rough (friction at the axle); the engine is snatched suddenly rather than released smoothly. Would the true safe working tension of the rope need to be higher or lower than your computed $T$?
\end{enumerate}

\courssub{Model answer}

\textbf{1. Modelling assumptions and free-body diagrams.}

Assumptions: the rope is \cle{light} (mass negligible) and \cle{inextensible} (both bodies have the same magnitude of acceleration); the beam acts as a \cle{smooth pulley} (same tension on both sides); the bodies are modelled as \cle{particles}; air resistance is neglected; $g$ is constant.

Forces on the engine: weight $Mg = 120 \times 9.8 = \SI{1176}{N}$ downwards, tension $T$ upwards. Forces on the counterweight: weight $mg = 90 \times 9.8 = \SI{882}{N}$ downwards, tension $T$ upwards.

\begin{popfigure}
\begin{center}
\begin{tikzpicture}[scale=1.0, every node/.style={font=\small}]
% engine FBD
\draw[fill=poporange!60, draw=popdark] (-3.4,0) rectangle (-2.2,0.8);
\node at (-2.8,0.4) {$M$};
\draw[->, line width=1.2pt, popblue] (-2.8,0.8) -- (-2.8,2.1) node[above]{$T$};
\draw[->, line width=1.2pt, poppink] (-2.8,0) -- (-2.8,-1.4) node[below]{$Mg=1176$ N};
\draw[->, line width=1.2pt, popgreen!60!black] (-2.2,0.4) -- (-1.3,0.4) node[right]{$a$};
% counterweight FBD
\draw[fill=popgreenL, draw=popgreen!60!black] (1.8,0) rectangle (3.0,0.8);
\node at (2.4,0.4) {$m$};
\draw[->, line width=1.2pt, popblue] (2.4,0.8) -- (2.4,2.1) node[above]{$T$};
\draw[->, line width=1.2pt, poppink] (2.4,0) -- (2.4,-1.15) node[below]{$mg=882$ N};
\draw[->, line width=1.2pt, popgreen!60!black] (1.8,0.4) -- (0.9,0.4) node[left]{$a$};
\end{tikzpicture}
\end{center}
\legende{Free-body diagrams of the engine (left) and the counterweight (right).}
\end{popfigure}

\textbf{2. Acceleration and tension.}

The engine (heavier) descends; take the direction of motion as positive for each body.

Engine ($\downarrow$ positive): $\;Mg - T = Ma$, i.e.
\[ 1176 - T = 120a. \tag{1} \]
Counterweight ($\uparrow$ positive): $\;T - mg = ma$, i.e.
\[ T - 882 = 90a. \tag{2} \]

Adding (1) and (2) eliminates $T$:
\[ 1176 - 882 = (120+90)a \;\Longrightarrow\; a = \frac{294}{210} = \SI{1.4}{m.s^{-2}}. \]

Substituting in (2):
\[ T = 882 + 90(1.4) = 882 + 126 = \SI{1008}{N}. \]

Safety check: $T = \SI{1008}{N} < \SI{1500}{N}$. \textbf{The rope is safe}, with a safety margin of $\SI{492}{N}$ (a load factor of $\frac{1500}{1008} \approx 1.49$).

\textbf{3. Speed after descending \SI{1.5}{m}.}

\emph{(a) suvat.} With $u = 0$, $a = \SI{1.4}{m.s^{-2}}$, $s = \SI{1.5}{m}$:
\[ v^2 = u^2 + 2as = 0 + 2(1.4)(1.5) = 4.2 \;\Longrightarrow\; v = \sqrt{4.2} \approx \SI{2.05}{m.s^{-1}}. \]

\emph{(b) Energy method.} For the whole system, the tension does no net work (it is an internal force and the rope is inextensible, so the works done by $T$ on the two bodies cancel). The engine loses $Mgh$ of potential energy while the counterweight gains $mgh$; both bodies gain kinetic energy, a total of $\tfrac12 (M+m)v^2$:
\[ Mgh - mgh = \tfrac{1}{2}(M+m)v^2 \]
\[ (120-90)(9.8)(1.5) = \tfrac{1}{2}(210)v^2 \]
\[ 441 = 105\,v^2 \;\Longrightarrow\; v^2 = 4.2 \;\Longrightarrow\; v \approx \SI{2.05}{m.s^{-1}}. \]

The two methods agree exactly, as expected: the suvat result uses the acceleration obtained from Newton's second law, and for this conservative system Newton's second law and energy conservation are two statements of the same physics. Agreement of the two methods is a valuable \cle{independent check} in the examination.

\textbf{4. Second stage: the counterweight after the engine lands.}

When the engine touches the floor, the floor exerts a normal reaction on it; the engine can no longer pull the rope down. The rope goes \cle{slack} ($T = 0$), and the counterweight, moving upwards at $v = \SI{2.05}{m.s^{-1}}$, is now in free motion under gravity alone: its acceleration is $-g = \SI{-9.8}{m.s^{-2}}$ (taking upwards as positive).

Extra height $h'$ risen before momentary rest (final speed $0$):
\[ 0 = v^2 - 2gh' \;\Longrightarrow\; h' = \frac{4.2}{2 \times 9.8} = \frac{4.2}{19.6} \approx \SI{0.214}{m}. \]

During the first stage the counterweight rose $\SI{1.5}{m}$, so its greatest height above its release position is
\[ 1.5 + 0.214 \approx \SI{1.71}{m}. \]

Consistency check: the counterweight started about $\SI{2.5}{m}$ below the beam (the beam is $\SI{4.0}{m}$ above the floor and the engine hung $\SI{1.5}{m}$ above the floor), so even after rising $\SI{1.71}{m}$ it does not strike the beam: the model remains consistent.

\textbf{5. Extension: the rough $20^\circ$ ramp.}

Resolve the engine's weight along and perpendicular to the ramp:
\[ Mg\sin 20^\circ = 1176 \times 0.342 \approx \SI{402.2}{N}, \qquad R = Mg\cos 20^\circ = 1176 \times 0.940 \approx \SI{1105.0}{N}. \]

Limiting friction (acting \emph{down} the ramp, opposing the intended upward motion):
\[ F_{\max} = \mu R = 0.3 \times 1105.0 \approx \SI{331.5}{N}. \]

For the engine to be on the point of moving \emph{up} the ramp, the rope must pull it up with a tension at least equal to the sum of the downslope component of the weight and the limiting friction:
\[ T_{\text{needed}} = Mg\sin 20^\circ + \mu Mg\cos 20^\circ = 402.2 + 331.5 = \SI{733.7}{N}. \]

\emph{(a) Can the $90$~kg counterweight do the job?} For the counterweight to pull the engine up, the counterweight must descend, and the tension it can transmit to the rope is then \emph{less than} its weight $mg = \SI{882}{N}$ (since $mg - T = ma$ with $a > 0$). Motion is possible if and only if the driving weight exceeds the total resistance:
\[ mg = \SI{882}{N} > \SI{733.7}{N} = T_{\text{needed}}. \]
The condition is satisfied, so the $90$~kg counterweight \emph{can} pull the engine up the ramp. Applying Newton's second law to the whole system in the direction of motion,
\[ mg - Mg\sin 20^\circ - \mu Mg\cos 20^\circ = (M+m)a \]
\[ a = \frac{882 - 733.7}{210} = \frac{148.3}{210} \approx \SI{0.71}{m.s^{-2}}. \]
(The corresponding tension is $T = m(g - a) = 90(9.8 - 0.71) \approx \SI{818}{N}$, comfortably below the rope rating.)

\emph{(b) Minimum counterweight mass.} The engine is just on the point of moving up when the system is in limiting equilibrium, i.e. $m_{\min} g = Mg\sin 20^\circ + \mu Mg\cos 20^\circ$:
\[ m_{\min} = M(\sin 20^\circ + \mu\cos 20^\circ) = 120(0.342 + 0.3 \times 0.940) = 120(0.624) \approx \SI{74.9}{kg}. \]

So any counterweight heavier than about $75$~kg suffices; the existing $90$~kg of scrap is adequate, with margin.

\begin{attention}[A trap in Task 5]
Many students assume that because the $90$~kg counterweight is lighter than the $120$~kg engine, it cannot drag the engine up the ramp. Always compute the actual driving and resisting forces: on a $20^\circ$ ramp only the component $Mg\sin 20^\circ$ (plus friction) must be overcome, not the full weight $Mg$. Here the required force is only about $62\%$ of the engine's weight.
\end{attention}

\textbf{6. Critique of the model.}

\begin{itemize}
\item \emph{Rope mass.} A real rope of several metres has non-negligible mass; part of the driving force is used to accelerate the rope itself, so the true acceleration is slightly smaller than $\SI{1.4}{m.s^{-2}}$ and the tension varies along the rope (greatest at the top, near the beam).
\item \emph{Friction at the beam.} A rough axle makes the tensions on the two sides unequal; part of the weight difference is ``wasted'' overcoming this friction, reducing $a$. The tension on the engine side could then exceed our computed $\SI{1008}{N}$.
\item \emph{Sudden jerks.} If the rope is snatched, or the engine swings and is caught abruptly, impulsive tensions occur which can momentarily far exceed the steady value. For this reason real safety codes require the rated tension of lifting equipment to be a generous multiple of the computed working tension; our margin of only $1.49$ would be judged inadequate in professional practice.
\item \emph{Particle model.} The engine is bulky; if it swings, part of the motion is rotational and the simple model no longer describes it exactly.
\end{itemize}

\textbf{Conclusion.} The model predicts a safe, gentle descent ($a = \SI{1.4}{m.s^{-2}}$, $T = \SI{1008}{N}$, landing speed $\SI{2.05}{m.s^{-1}}$), a counterweight overshoot of about $\SI{0.21}{m}$ above its release position, and shows that the same scrap counterweight can haul the engine back up the rough ramp with acceleration about $\SI{0.71}{m.s^{-2}}$. The critique shows, however, that for real workshop use Mr.~Ekane should choose a rope rated well above $\SI{1500}{N}$, release the load smoothly, and keep the beam axle well greased.

\newpage

\coursec{Methods and worked examples}

\courssub{Skill 1: Setting Up Equations of Motion for Two Connected Particles on a Smooth Surface}

Two particles are said to be \cle{connected} when they are joined by a string, rod or tow-bar so that the motion of one constrains the motion of the other. Throughout this chapter we use the standard GCE modelling assumptions: the string is \cle{light} (its mass is negligible) and \cle{inextensible} (its length is constant), and every pulley is \cle{smooth} (no friction in its axle) and light. These assumptions have two crucial consequences:

\begin{propriete}[Consequences of the modelling assumptions]
For a light inextensible string passing over a smooth pulley:
\begin{itemize}
\item the \cle{tension} has the \textbf{same value} $T$ at every point of the string;
\item the two particles have accelerations of the \textbf{same magnitude} $a$ (and velocities of the same magnitude), because the length of the string is constant.
\end{itemize}
These facts are used in every problem of this chapter; proofs are not examinable, but you must be able to \emph{use} the results.
\end{propriete}

\begin{methode}[Two particles connected on a smooth horizontal surface]
To analyse two particles of masses $m_1$ and $m_2$ connected by a light inextensible string (possibly passing over a smooth pulley), with no friction:
\begin{enumerate}
\item \textbf{Draw} a clear diagram showing each particle \emph{separately}, with all forces acting \emph{on that particle only} (weight, normal reaction, tension).
\item \textbf{Choose} a positive direction of motion for each particle, consistent with the string: because the string is inextensible, both particles have the \cle{same magnitude of acceleration} $a$ and the same tension $T$ acts throughout (string light, pulley smooth).
\item \textbf{Apply} Newton's second law $F = ma$ to each particle \emph{along its direction of motion}.
\item \textbf{Solve} the two simultaneous equations: add them to eliminate $T$ and find $a$; substitute back to find $T$.
\item \textbf{Check}: $T$ must be positive and less than the weight of the descending mass; $a < g$ always.
\end{enumerate}
\textbf{Key reflex:} never write the weight of particle 1 in the equation of particle 2. Each particle has its own equation.
\end{methode}

\begin{popfigure}
\begin{center}
\begin{tikzpicture}[scale=1.0, >=stealth]
% table
\draw[thick, popdark] (-0.5,0) -- (4.5,0);
\draw[thick, popdark] (4.5,0) -- (4.5,-0.25);
\draw[popmuted] (0.3,0) -- (0.3,-1.6);
\draw[popmuted] (0.15,-1.6) -- (0.45,-1.6);
% block A
\draw[fill=popblueL] (1.2,0) rectangle (2.4,0.7);
\node at (1.8,0.35) {$A$};
\node[below] at (1.8,-0.05) {\small $m_1$};
% pulley
\draw[fill=white] (4.5,0) circle (0.22);
\fill (4.5,0) circle (0.03);
% string
\draw[thick, poporange] (2.4,0.35) -- (4.5,0.22);
\draw[thick, poporange] (4.72,0) -- (4.72,-1.5);
% block B
\draw[fill=popgreenL] (4.42,-1.5) rectangle (5.02,-2.2);
\node at (4.72,-1.85) {$B$};
\node[right] at (5.05,-1.85) {\small $m_2$};
% forces on A
\draw[->, thick, popblue] (1.8,0.35) -- (3.1,0.35) node[midway, above]{$T$};
\draw[->, thick, popdark] (1.8,0.9) -- (1.8,1.6) node[above]{$R$};
\draw[->, thick, popdark] (1.4,0.35) -- (1.4,-0.6) node[below]{$m_1 g$};
% forces on B
\draw[->, thick, popblue] (4.72,-1.5) -- (4.72,-0.7) node[midway, right]{$T$};
\draw[->, thick, popdark] (4.72,-2.2) -- (4.72,-3.0) node[below]{$m_2 g$};
% acceleration arrows
\draw[->, very thick, poppink] (2.6,1.0) -- (3.6,1.0) node[right]{$a$};
\draw[->, very thick, poppink] (5.35,-0.9) -- (5.35,-1.7) node[right]{$a$};
\end{tikzpicture}
\end{center}
\legende{Particle $A$ on a smooth table connected over a smooth pulley to the hanging particle $B$. Each particle carries its own forces; the tension $T$ is the same on both sides.}
\end{popfigure}

\begin{exoresolu}[Car pulling a trailer on a smooth surface]
A particle $A$ of mass $4$ kg rests on a smooth horizontal table. It is connected by a light inextensible string, passing over a small smooth pulley fixed at the edge of the table, to a particle $B$ of mass $2$ kg hanging freely. The system is released from rest. Taking $g = 9.8\ \mathrm{m\,s^{-2}}$, find (a) the acceleration of the system, (b) the tension in the string, (c) the speed of $B$ after it has fallen $0.9$ m.
\tcblower
\textbf{Solution.} The table is smooth, so the only horizontal force on $A$ is the tension $T$. The string is light and inextensible and the pulley smooth, so the tension is the same on both sides and both particles have acceleration of magnitude $a$.

\textbf{Equation for $A$ (horizontal motion):}
\[ T = 4a \quad (1) \]

\textbf{Equation for $B$ (vertical motion, downwards positive):}
\[ 2g - T = 2a \quad (2) \]

\textbf{(a) Acceleration.} Add (1) and (2) to eliminate $T$:
\[ 2g = 6a \quad\Longrightarrow\quad a = \frac{2 \times 9.8}{6} = \frac{19.6}{6} \approx 3.27\ \mathrm{m\,s^{-2}}. \]

\textbf{(b) Tension.} From (1):
\[ T = 4a = 4 \times 3.27 \approx 13.1\ \mathrm{N}. \]
Check: $T = 13.1$ N $< 2g = 19.6$ N, as expected since $B$ accelerates downwards.

\textbf{(c) Speed after falling $0.9$ m.} The acceleration is constant, so use $v^2 = u^2 + 2as$ with $u = 0$, $s = 0.9$ m:
\[ v^2 = 0 + 2 \times 3.27 \times 0.9 = 5.88 \quad\Longrightarrow\quad v \approx 2.42\ \mathrm{m\,s^{-1}}. \]
\end{exoresolu}

\begin{attention}[Frequent mistakes with two connected particles]
\begin{itemize}
\item Writing $T = m_2 g$ for the hanging mass. This is only true in \emph{equilibrium}; when the system accelerates, $T \neq m_2 g$.
\item Giving the two particles different accelerations. An inextensible string forces $|a_1| = |a_2|$.
\item Mixing the two equations: the mass in $F = ma$ must be the mass of \emph{that} particle, and the forces must be the forces acting \emph{on that particle}.
\end{itemize}
\end{attention}

\courssub{Skill 2: Motion on a Rough Horizontal Surface}

\begin{methode}[Connected particles with friction]
When one particle moves on a \cle{rough} surface:
\begin{enumerate}
\item \textbf{Identify} the normal reaction $R$ on the particle on the surface. On a horizontal surface with no vertical component of other forces, $R = mg$.
\item \textbf{Compute} the friction force $F = \mu R$ (motion occurs, so friction is limiting: $F = \mu R$), directed \emph{opposite} to the motion.
\item \textbf{Write} Newton's second law for each particle, including $F$ in the equation of the particle on the surface.
\item \textbf{Solve} simultaneously for $a$ and $T$.
\item \textbf{Check the motion assumption}: motion only occurs if the driving force exceeds $\mu R$; if $a$ comes out negative, the assumption of motion in that direction was wrong.
\end{enumerate}
\textbf{Key formula:} for mass $m_1$ on a rough table pulled by hanging mass $m_2$:
\[ a = \frac{(m_2 - \mu m_1)g}{m_1 + m_2}. \]
\end{methode}

\begin{exoresolu}[Rough table problem]
A block $P$ of mass $5$ kg lies on a rough horizontal table, the coefficient of friction between the block and the table being $0.2$. $P$ is connected by a light inextensible string passing over a smooth pulley at the edge of the table to a block $Q$ of mass $3$ kg hanging freely. The system is released from rest. Taking $g = 9.8\ \mathrm{m\,s^{-2}}$, find (a) the acceleration of the blocks, (b) the tension in the string, (c) the time taken for $Q$ to descend $1.2$ m.
\tcblower
\textbf{Solution.}

\textbf{Step 1: Reaction and friction on $P$.} Vertically, $P$ is in equilibrium:
\[ R = 5g = 5 \times 9.8 = 49\ \mathrm{N}. \]
Since $P$ slides, friction is limiting:
\[ F = \mu R = 0.2 \times 49 = 9.8\ \mathrm{N}, \]
acting opposite to the motion of $P$.

\textbf{Step 2: Equations of motion.} Let $a$ be the acceleration and $T$ the tension.

For $P$ (horizontal): $\quad T - F = 5a$, i.e.
\[ T - 9.8 = 5a \quad (1) \]

For $Q$ (vertical, downwards positive):
\[ 3g - T = 3a \quad\Longrightarrow\quad 29.4 - T = 3a \quad (2) \]

\textbf{Step 3: Solve.} Adding (1) and (2):
\[ 29.4 - 9.8 = 8a \quad\Longrightarrow\quad a = \frac{19.6}{8} = 2.45\ \mathrm{m\,s^{-2}}. \]
Since $a > 0$, motion does occur as assumed (indeed $3g = 29.4 > F = 9.8$).

From (1): $T = 5a + 9.8 = 12.25 + 9.8 = 22.05\ \mathrm{N}$.

\textbf{Step 4: Time to descend $1.2$ m.} Using $s = ut + \tfrac{1}{2}at^2$ with $u = 0$:
\[ 1.2 = \tfrac{1}{2} \times 2.45 \times t^2 \quad\Longrightarrow\quad t^2 = \frac{2.4}{2.45} = 0.9796 \quad\Longrightarrow\quad t \approx 0.99\ \mathrm{s}. \]
\end{exoresolu}

\begin{attention}[Friction opposes motion, not the tension]
Friction always acts \emph{opposite to the direction of sliding} (or of impending sliding). Do not automatically place $F$ on the left of your diagram: decide first which way the particle moves, then put $F$ the other way. Also remember $F = \mu R$ is valid only when sliding actually occurs; otherwise $F \leq \mu R$.
\end{attention}

\courssub{Skill 3: Motion on an Inclined Plane}

\begin{methode}[Particle on an inclined plane connected to a hanging mass]
\begin{enumerate}
\item \textbf{Resolve the weight} of the particle on the plane into components: $mg\sin\theta$ \emph{down the plane} and $mg\cos\theta$ \emph{into the plane}.
\item \textbf{Normal reaction}: $R = mg\cos\theta$ (no acceleration perpendicular to the plane).
\item \textbf{Friction} (if rough): $F = \mu R = \mu mg\cos\theta$, opposing the motion (i.e. acting down the plane if the particle moves up, and vice versa).
\item \textbf{Write} $F = ma$ along the plane for the particle on the plane, and vertically for the hanging particle.
\item \textbf{Solve} simultaneously. Decide the direction of motion first by comparing $m_2 g$ with $m_1 g\sin\theta + \mu m_1 g\cos\theta$.
\end{enumerate}
\textbf{Key reflex:} on a \emph{smooth} plane, $a = \dfrac{(m_2 - m_1\sin\theta)g}{m_1 + m_2}$ when $m_2$ descends.
\end{methode}

\begin{popfigure}
\begin{center}
\begin{tikzpicture}[scale=1.0, >=stealth]
% inclined plane
\draw[thick, popdark] (0,0) -- (6,0);
\draw[thick, popdark] (0,0) -- (5.2,2.6);
\draw[popdark] (1.2,0) arc (0:26.565:1.2);
\node at (1.65,0.22) {$\theta$};
% block A on plane (rotate)
\begin{scope}[rotate around={26.565:(2.2,1.1)}]
\draw[fill=popblueL] (1.7,1.1) rectangle (2.7,1.75);
\node at (2.2,1.425) {$A$};
\end{scope}
% pulley at top
\draw[fill=white] (5.2,2.6) circle (0.2);
\fill (5.2,2.6) circle (0.03);
% string parallel to slope then vertical
\draw[thick, poporange] (2.62,1.62) -- (5.2,2.8);
\draw[thick, poporange] (5.4,2.6) -- (5.4,1.0);
% hanging mass B
\draw[fill=popgreenL] (5.1,0.4) rectangle (5.7,1.0);
\node at (5.4,0.7) {$B$};
% weight components of A
\draw[->, thick, popdark] (2.2,1.55) -- (2.2,0.35) node[below left]{$m_1 g$};
\draw[->, thick, poppurple] (2.2,1.55) -- (1.31,1.11) node[below left]{$m_1 g\sin\theta$};
\draw[->, thick, poppurple] (2.2,1.55) -- (1.76,0.66) node[left]{$m_1 g\cos\theta$};
% tension on A
\draw[->, thick, popblue] (2.75,1.7) -- (3.75,2.2) node[above right]{$T$};
% forces on B
\draw[->, thick, popblue] (5.4,1.0) -- (5.4,1.8) node[right]{$T$};
\draw[->, thick, popdark] (5.4,0.4) -- (5.4,-0.4) node[below]{$m_2 g$};
\end{tikzpicture}
\end{center}
\legende{Resolving the weight on an inclined plane: component $m_1 g\sin\theta$ down the plane, $m_1 g\cos\theta$ perpendicular into the plane.}
\end{popfigure}

\begin{exoresolu}[Smooth inclined plane — motion does occur]
A particle $A$ of mass $3$ kg rests on a smooth plane inclined at $30°$ to the horizontal. $A$ is connected by a light inextensible string, parallel to the line of greatest slope and passing over a smooth pulley at the top of the plane, to a particle $B$ of mass $5$ kg hanging freely. The system is released from rest. Taking $g = 9.8\ \mathrm{m\,s^{-2}}$, find the acceleration and the tension in the string.
\tcblower
\textbf{Solution.} Component of $A$'s weight down the plane: $3g\sin 30° = 3 \times 9.8 \times 0.5 = 14.7$ N. Weight of $B$: $5g = 49$ N. Since $49 > 14.7$ and the plane is smooth, $B$ descends and $A$ moves up the plane.

\textbf{Equation for $A$ (up the plane positive):}
\[ T - 3g\sin 30° = 3a \quad\Longrightarrow\quad T - 14.7 = 3a \quad (1) \]

\textbf{Equation for $B$ (downwards positive):}
\[ 5g - T = 5a \quad\Longrightarrow\quad 49 - T = 5a \quad (2) \]

\textbf{Solve.} Adding (1) and (2):
\[ 49 - 14.7 = 8a \quad\Longrightarrow\quad a = \frac{34.3}{8} = 4.29\ \mathrm{m\,s^{-2}}. \]
From (1): $T = 14.7 + 3 \times 4.29 = 14.7 + 12.86 = 27.56\ \mathrm{N}$.

Check: $3g\sin 30° = 14.7 < T = 27.56 < 49 = 5g$. \checkmark
\end{exoresolu}

\begin{exoresolu}[Rough inclined plane — when nothing moves]
A particle $A$ of mass $6$ kg rests on a rough plane inclined at $30°$ to the horizontal, the coefficient of friction being $\tfrac{1}{4}$. $A$ is connected by a light inextensible string, parallel to the line of greatest slope and passing over a smooth pulley at the top of the plane, to a particle $B$ of mass $4$ kg hanging freely. The system is released from rest. Taking $g = 9.8\ \mathrm{m\,s^{-2}}$, find the acceleration of the system and the tension in the string.
\tcblower
\textbf{Solution.}

\textbf{Step 1: Which way does the system move?} Component of $A$'s weight down the plane:
\[ 6g\sin 30° = 6 \times 9.8 \times 0.5 = 29.4\ \mathrm{N}. \]
Weight of $B$: $4g = 39.2$ N. Since $39.2 > 29.4$ even before friction, we first test the hypothesis that $B$ descends and $A$ moves \emph{up} the plane; friction on $A$ would then act \emph{down} the plane.

\textbf{Step 2: Reaction and friction.}
\[ R = 6g\cos 30° = 6 \times 9.8 \times \frac{\sqrt{3}}{2} = 29.4\sqrt{3} \approx 50.9\ \mathrm{N}, \]
\[ F = \mu R = \tfrac{1}{4} \times 50.9 \approx 12.7\ \mathrm{N}. \]

\textbf{Step 3: Test the hypothesis.} For motion with $B$ descending we would need the driving force to exceed the resisting forces:
\[ 4g \;?\; 6g\sin 30° + F \quad\Longleftrightarrow\quad 39.2 \;?\; 29.4 + 12.7 = 42.1. \]
Since $39.2 < 42.1$, the hanging mass cannot pull $A$ up the plane.

\textbf{Step 4: Test the opposite direction.} Could $A$ slide down, pulling $B$ up? Then friction acts up the plane and we would need
\[ 6g\sin 30° \;?\; 4g + F \quad\Longleftrightarrow\quad 29.4 \;?\; 39.2 + 12.7 = 51.9, \]
which is also false.

\textbf{Conclusion:} the system remains at rest: $a = 0$. Equilibrium of $B$ gives $T = 4g = 39.2$ N, and equilibrium of $A$ along the plane gives friction $= T - 6g\sin 30° = 39.2 - 29.4 = 9.8$ N, acting down the plane. Since $9.8 < \mu R = 12.7$ N, friction is \emph{not} limiting, which is consistent with rest.

\textbf{Lesson:} always verify the direction of motion; if the computed acceleration comes out negative, the assumed motion does not occur and friction may not be limiting.
\end{exoresolu}

\courssub{Skill 4: Atwood's Machine}

\begin{methode}[Atwood's machine: two masses over a fixed pulley]
For masses $m_1 > m_2$ hanging on either side of a light inextensible string over a smooth fixed pulley:
\begin{enumerate}
\item \textbf{Take} downwards positive for the heavier mass $m_1$, upwards positive for $m_2$; both have acceleration of magnitude $a$.
\item \textbf{Equations:} $m_1 g - T = m_1 a$ and $T - m_2 g = m_2 a$.
\item \textbf{Add} to get the standard results:
\[ a = \frac{(m_1 - m_2)g}{m_1 + m_2}, \qquad T = \frac{2m_1 m_2 g}{m_1 + m_2}. \]
\item \textbf{Then} use constant-acceleration kinematics ($v = u + at$, $s = ut + \tfrac12 at^2$, $v^2 = u^2 + 2as$) for any further question.
\end{enumerate}
\textbf{Check:} $m_2 g < T < m_1 g$; if $m_1 = m_2$ then $a = 0$ and $T = m_1 g$.
\end{methode}

\begin{popfigure}
\begin{center}
\begin{tikzpicture}[scale=1.0, >=stealth]
% support and pulley
\draw[thick, popdark] (-1.6,2.6) -- (1.6,2.6);
\draw[fill=white, thick] (0,2.2) circle (0.4);
\fill (0,2.2) circle (0.04);
\draw[thick, popdark] (0,2.6) -- (0,2.2);
% string
\draw[thick, poporange] (-0.4,2.2) arc (180:0:0.4);
\draw[thick, poporange] (-0.4,2.2) -- (-0.4,0.6);
\draw[thick, poporange] (0.4,2.2) -- (0.4,1.2);
% masses
\draw[fill=popblueL] (-0.75,-0.2) rectangle (-0.05,0.6);
\node at (-0.4,0.2) {$m_1$};
\draw[fill=popgreenL] (0.1,0.4) rectangle (0.7,1.2);
\node at (0.4,0.8) {$m_2$};
% tensions
\draw[->, thick, popblue] (-0.4,0.6) -- (-0.4,1.5) node[midway, left]{$T$};
\draw[->, thick, popblue] (0.4,1.2) -- (0.4,1.95) node[midway, right]{$T$};
% weights
\draw[->, thick, popdark] (-0.4,-0.2) -- (-0.4,-1.1) node[below]{$m_1 g$};
\draw[->, thick, popdark] (0.4,0.4) -- (0.4,-0.5) node[below]{$m_2 g$};
% accelerations
\draw[->, very thick, poppink] (-1.15,0.5) -- (-1.15,-0.3) node[below]{$a$};
\draw[->, very thick, poppink] (1.1,0.3) -- (1.1,1.1) node[above]{$a$};
\end{tikzpicture}
\end{center}
\legende{Atwood's machine: the heavier mass $m_1$ accelerates downwards, the lighter mass $m_2$ upwards, with the same magnitude $a$; the tension $T$ is the same on both sides.}
\end{popfigure}

\begin{exoresolu}[Classic Atwood's machine]
Particles of masses $5$ kg and $3$ kg are attached to the ends of a light inextensible string passing over a smooth fixed pulley. The system is released from rest with the string taut and the $5$ kg particle $2$ m above the ground. Taking $g = 9.8\ \mathrm{m\,s^{-2}}$, find (a) the acceleration and the tension, (b) the speed of the $5$ kg mass when it hits the ground, (c) the time taken.
\tcblower
\textbf{Solution.}

\textbf{(a)} With $m_1 = 5$, $m_2 = 3$:
\[ a = \frac{(5-3)g}{5+3} = \frac{2 \times 9.8}{8} = 2.45\ \mathrm{m\,s^{-2}}, \]
\[ T = \frac{2 \times 5 \times 3 \times 9.8}{8} = \frac{294}{8} = 36.75\ \mathrm{N}. \]
Check: $3g = 29.4 < 36.75 < 49 = 5g$. \checkmark

\textbf{(b)} Using $v^2 = u^2 + 2as$ with $u = 0$, $s = 2$ m:
\[ v^2 = 2 \times 2.45 \times 2 = 9.8 \quad\Longrightarrow\quad v = \sqrt{9.8} \approx 3.13\ \mathrm{m\,s^{-1}}. \]

\textbf{(c)} Using $s = \tfrac12 a t^2$:
\[ t = \sqrt{\frac{2s}{a}} = \sqrt{\frac{4}{2.45}} = \sqrt{1.633} \approx 1.28\ \mathrm{s}. \]
\end{exoresolu}

\begin{savaistu}[Why Atwood built this machine]
George Atwood introduced his machine in 1784 to \emph{slow down} gravity: with nearly equal masses, $a = \dfrac{(m_1 - m_2)g}{m_1 + m_2}$ is a small, easily measured fraction of $g$, so the laws of accelerated motion could be checked accurately with the clocks of the time. The same trick — diluting $g$ — is used today in school laboratories with air tracks and ticker-timers.
\end{savaistu}

\courssub{Skill 5: Multi-Stage Motion (String Breaks or a Mass is Removed)}

\begin{methode}[Multi-stage problems]
When the string breaks, or a mass hits the ground and the string goes slack:
\begin{enumerate}
\item \textbf{Stage 1:} analyse the connected motion as usual; find the speed $v_1$ and positions at the instant the regime changes.
\item \textbf{Stage 2:} after the change, each remaining free particle moves \emph{under gravity alone} ($a = \pm g$) if the string is slack, or a new connected system forms with new equations. The \cle{initial velocity} of stage 2 is the final velocity of stage 1.
\item \textbf{Combine} the stages to answer questions on total time, greatest height, etc.
\end{enumerate}
\textbf{Trap:} when the string goes slack, the rising mass continues upwards, decelerating at $g$ — it does not stop immediately.
\end{methode}

\begin{exoresolu}[Mass hits the ground; how high does the other rise?]
In an Atwood's machine, masses of $4$ kg and $2$ kg hang at the same height, $1.5$ m above the floor, connected over a smooth fixed pulley. The system is released from rest. The $4$ kg mass falls to the floor and does not rebound. Taking $g = 9.8\ \mathrm{m\,s^{-2}}$, find the greatest height above the floor reached by the $2$ kg mass.
\tcblower
\textbf{Solution.}

\textbf{Stage 1: connected motion.} The acceleration is
\[ a = \frac{(4-2)g}{4+2} = \frac{2 \times 9.8}{6} = \frac{9.8}{3} \approx 3.267\ \mathrm{m\,s^{-2}}. \]
When the $4$ kg mass has fallen $1.5$ m, the common speed is given by $v^2 = 2as$:
\[ v_1^2 = 2 \times \frac{9.8}{3} \times 1.5 = 9.8 \quad\Longrightarrow\quad v_1 = \sqrt{9.8} \approx 3.13\ \mathrm{m\,s^{-1}}. \]
At this instant the $2$ kg mass is at height $1.5 + 1.5 = 3$ m, moving \emph{upwards} at $3.13\ \mathrm{m\,s^{-1}}$.

\textbf{Stage 2: string slack.} Once the $4$ kg mass hits the floor, the string goes slack and the $2$ kg mass moves freely under gravity, decelerating at $g$. The extra height gained is given by $v^2 = u^2 - 2gh$ with final $v = 0$:
\[ 0 = 9.8 - 2 \times 9.8 \times h \quad\Longrightarrow\quad h = \frac{9.8}{19.6} = 0.5\ \mathrm{m}. \]

\textbf{Greatest height:}
\[ H = 3 + 0.5 = 3.5\ \mathrm{m} \text{ above the floor.} \]
(One must of course assume the pulley is high enough for this to be possible.)
\end{exoresolu}

\begin{exercice}[Two-stage Atwood problem]
Two particles of masses $6$ kg and $4$ kg are connected by a light inextensible string over a smooth fixed pulley. Initially both are held at rest, $1.8$ m above the floor, and then released. The $6$ kg particle strikes the floor and does not rebound. Taking $g = 9.8\ \mathrm{m\,s^{-2}}$, find (a) the speed of the $4$ kg particle at the instant the $6$ kg particle hits the floor, (b) the greatest height reached by the $4$ kg particle, (c) the total time from release until the $4$ kg particle reaches its greatest height.
\end{exercice}

\begin{corrige}[Exercise 1]
\textbf{(a)} Acceleration: $a = \dfrac{(6-4)g}{6+4} = \dfrac{2 \times 9.8}{10} = 1.96\ \mathrm{m\,s^{-2}}$. Speed after falling $1.8$ m: $v_1^2 = 2 \times 1.96 \times 1.8 = 7.056$, so $v_1 \approx 2.66\ \mathrm{m\,s^{-1}}$.

\textbf{(b)} Height at that instant: $1.8 + 1.8 = 3.6$ m. Extra rise under gravity: $h = \dfrac{v_1^2}{2g} = \dfrac{7.056}{19.6} = 0.36$ m. Greatest height: $H = 3.6 + 0.36 = 3.96$ m.

\textbf{(c)} Stage 1 time: $t_1 = \dfrac{v_1}{a} = \dfrac{2.66}{1.96} \approx 1.36$ s. Stage 2 time: $t_2 = \dfrac{v_1}{g} = \dfrac{2.66}{9.8} \approx 0.27$ s. Total: $t \approx 1.63$ s.
\end{corrige}

\courssub{Skill 6: Pulley Systems with More Than Two Masses}

\begin{methode}[Three or more connected masses]
\begin{enumerate}
\item \textbf{Draw} each mass separately and mark the tension in \emph{each separate string}: different strings generally have \textbf{different} tensions $T_1, T_2, \dots$
\item \textbf{Use inextensibility} to relate the accelerations. In the common GCE case (masses in series on a table, or one mass on a table with two hanging masses on separate strings), all accelerations have the same magnitude $a$.
\item \textbf{Write} $F = ma$ for \emph{every} mass — one equation per mass.
\item \textbf{Solve} the system: adding all equations eliminates all internal tensions at once and gives $a$ directly; then back-substitute for each tension.
\end{enumerate}
\textbf{Key reflex:} internal tensions cancel when you add all the equations — the "driving force" is the sum of external forces along the direction of motion, and the "total mass" is the sum of all masses.
\end{methode}

\begin{exoresolu}[Two hanging masses, one on a smooth table]
A particle $A$ of mass $3$ kg lies on a smooth horizontal table. A light inextensible string connects $A$ over a smooth pulley at one edge of the table to a hanging mass $B$ of $2$ kg; a second light inextensible string connects $A$ over a smooth pulley at the opposite edge to a hanging mass $C$ of $5$ kg. The system is released from rest. Taking $g = 9.8\ \mathrm{m\,s^{-2}}$, find the acceleration of the system and the tension in each string.
\tcblower
\textbf{Solution.} Since $5 > 2$, mass $C$ descends, $B$ rises, and $A$ moves towards $C$'s edge. Let $T_1$ be the tension in the string joining $A$ to $B$, and $T_2$ the tension in the string joining $A$ to $C$. All accelerations have magnitude $a$ (strings inextensible).

\textbf{Equation for $C$ (downwards positive):}
\[ 5g - T_2 = 5a \quad (1) \]

\textbf{Equation for $A$ (towards $C$ positive):} the table is smooth, so
\[ T_2 - T_1 = 3a \quad (2) \]

\textbf{Equation for $B$ (upwards positive):}
\[ T_1 - 2g = 2a \quad (3) \]

\textbf{Solve.} Adding (1), (2) and (3) eliminates both tensions:
\[ 5g - 2g = (5 + 3 + 2)a \quad\Longrightarrow\quad 3g = 10a \quad\Longrightarrow\quad a = \frac{3 \times 9.8}{10} = 2.94\ \mathrm{m\,s^{-2}}. \]

From (3): $T_1 = 2g + 2a = 19.6 + 5.88 = 25.48\ \mathrm{N}$.

From (1): $T_2 = 5g - 5a = 49 - 14.7 = 34.3\ \mathrm{N}$.

\textbf{Check with (2):} $T_2 - T_1 = 34.3 - 25.48 = 8.82 = 3 \times 2.94 = 3a$. \checkmark

Note $T_1 \neq T_2$: the two strings are different, so their tensions differ.
\end{exoresolu}

\begin{attention}[One string, one tension — two strings, two tensions]
The rule "same tension throughout" applies to \emph{one} light string over smooth pulleys. As soon as the system contains two distinct strings, you must introduce two distinct unknowns $T_1$ and $T_2$. Writing a single $T$ everywhere in a two-string system is one of the most common errors at the GCE.
\end{attention}

\textbf{Variable acceleration in connected systems.}\par
In some problems the driving force varies with position or time (for example an elastic rope, or a force $F(t)$ applied to one mass). The method is unchanged in spirit: write Newton's second law for each particle, eliminate the tension, and obtain $a$ as a function of $x$ or $t$. Then, instead of the constant-acceleration formulae, use
\[ a = \frac{dv}{dt} \qquad \text{or} \qquad a = v\,\frac{dv}{dx}, \]
and integrate. For instance, if after eliminating $T$ you obtain $a = 4 - 0.5\,x$ (in SI units) for a system starting from rest at $x = 0$, then
\[ v\,\frac{dv}{dx} = 4 - 0.5x \quad\Longrightarrow\quad \tfrac12 v^2 = 4x - 0.25x^2, \]
which gives the speed at any position. The constant-acceleration formulae $v = u + at$, $s = ut + \tfrac12 at^2$ must \textbf{not} be used when $a$ is not constant.

\courssub{Skill 7: Energy Methods for Connected Particles}

\begin{methode}[Work–energy approach]
For connected systems, energy methods often give speeds much faster than Newton's second law:
\begin{enumerate}
\item \textbf{Identify} the displacement $s$ of each mass (equal magnitudes for a single inextensible string).
\item \textbf{Write the energy balance:}
\[ \text{loss of GPE} = \text{gain of KE} + \text{work done against friction}. \]
The tension does \textbf{no net work} on the whole system (it is internal), so it disappears — this is the great advantage of the method.
\item \textbf{KE of the system:} $E_k = \tfrac12(m_1 + m_2)v^2$ since both masses have the same speed $v$.
\item \textbf{Friction term:} $\mu R\, s = \mu m g\cos\theta \cdot s$ for a mass sliding $s$ along a rough plane.
\item \textbf{Solve} for $v$. If the acceleration is needed afterwards, use $v^2 = 2as$ (constant $a$) or differentiate the energy equation.
\end{enumerate}
\textbf{When to use it:} whenever the question asks for a \emph{speed} after a given \emph{distance}, and especially for real-life problems (lifts, vehicles, loads on ramps).
\end{methode}

\begin{exoresolu}[Energy on a rough slope — a loading ramp]
At a warehouse in Douala, a crate of mass $50$ kg is pulled up a rough ramp inclined at $20°$ to the horizontal (coefficient of friction $0.3$) by a light inextensible rope parallel to the ramp, passing over a smooth pulley at the top, with a counterweight of mass $40$ kg hanging vertically. The system is released from rest. Taking $g = 9.8\ \mathrm{m\,s^{-2}}$, use energy considerations to find the speed of the crate after it has moved $2$ m up the ramp, and verify the result by computing the acceleration from Newton's second law.
\tcblower
\textbf{Solution.}

\textbf{Step 1: Energy balance over $s = 2$ m.}

The counterweight descends $2$ m: loss of GPE $= 40 \times 9.8 \times 2 = 784$ J.

The crate rises vertically by $2\sin 20°$: gain of GPE
\[ = 50 \times 9.8 \times 2\sin 20° = 980 \times 0.342 = 335.2\ \mathrm{J}. \]

Work done against friction on the crate: $R = 50g\cos 20° = 490 \times 0.9397 = 460.4$ N, so
\[ W_F = \mu R s = 0.3 \times 460.4 \times 2 = 276.3\ \mathrm{J}. \]

Gain of kinetic energy of the whole system (both masses have speed $v$):
\[ E_k = \tfrac12(50 + 40)v^2 = 45v^2. \]

\textbf{Step 2: Apply the energy equation.}
\[ 784 = 335.2 + 276.3 + 45v^2 \]
\[ 45v^2 = 172.5 \quad\Longrightarrow\quad v^2 = 3.834 \quad\Longrightarrow\quad v \approx 1.96\ \mathrm{m\,s^{-1}}. \]

\textbf{Step 3: Verification by Newton's second law.}

For the counterweight: $40g - T = 40a$, i.e. $392 - T = 40a$.

For the crate along the ramp: $T - 50g\sin 20° - \mu(50g\cos 20°) = 50a$, i.e.
\[ T - 167.6 - 138.1 = 50a. \]

Adding: $392 - 305.7 = 90a$, so $a = \dfrac{86.3}{90} = 0.959\ \mathrm{m\,s^{-2}}$.

Then $v^2 = 2as = 2 \times 0.959 \times 2 = 3.836$, giving $v \approx 1.96\ \mathrm{m\,s^{-1}}$ — the same result. \checkmark

\textbf{Remark:} the energy method never required the tension; it is shorter and less error-prone when only a speed is required.
\end{exoresolu}

\begin{exercice}[Energy check on Atwood's machine]
In an Atwood's machine with masses $5$ kg and $3$ kg released from rest, use the conservation of mechanical energy to find the common speed after the $5$ kg mass has descended $2$ m. Take $g = 9.8\ \mathrm{m\,s^{-2}}$. Compare with the answer obtained from $a = \dfrac{(m_1 - m_2)g}{m_1 + m_2}$ and $v^2 = 2as$.
\end{exercice}

\begin{corrige}[Exercise 2]
Loss of GPE of the $5$ kg mass: $5 \times 9.8 \times 2 = 98$ J. Gain of GPE of the $3$ kg mass: $3 \times 9.8 \times 2 = 58.8$ J. Net loss $= 98 - 58.8 = 39.2$ J, converted into kinetic energy of both masses:
\[ \tfrac12(5+3)v^2 = 39.2 \quad\Longrightarrow\quad 4v^2 = 39.2 \quad\Longrightarrow\quad v^2 = 9.8 \quad\Longrightarrow\quad v \approx 3.13\ \mathrm{m\,s^{-1}}. \]
By dynamics: $a = \dfrac{2 \times 9.8}{8} = 2.45\ \mathrm{m\,s^{-2}}$, $v^2 = 2 \times 2.45 \times 2 = 9.8$. Same result. \checkmark
\end{corrige}

\begin{aretenir}[Chapter essentials]
\begin{itemize}
\item Light inextensible string over a smooth pulley: \textbf{same tension} throughout, \textbf{same magnitude of acceleration} for both particles.
\item One equation $F = ma$ \emph{per particle}; add the equations to eliminate the tension.
\item Rough surface: $R = mg$ (horizontal) or $R = mg\cos\theta$ (inclined plane); sliding friction $F = \mu R$ opposes the motion.
\item Atwood's machine: $a = \dfrac{(m_1 - m_2)g}{m_1 + m_2}$, $T = \dfrac{2m_1 m_2 g}{m_1 + m_2}$.
\item Always \textbf{check the direction of motion}: a negative acceleration means the assumed motion does not occur and friction may not be limiting.
\item Multi-stage motion: the final velocity of one stage is the initial velocity of the next; after the string goes slack, particles move under gravity alone.
\item Distinct strings have distinct tensions; adding all equations of motion eliminates every internal tension at once.
\item Energy method: loss of GPE $=$ gain of KE $+$ work against friction; the tension disappears, giving speeds directly.
\end{itemize}
\end{aretenir}

\newpage

\coursec{Exercises}

\courssub{I check my knowledge}

\begin{exercice}[Multiple choice questions]
For each question, choose the single correct answer. Take $g = \SI{9.8}{m.s^{-2}}$.
\begin{enumerate}
\item Two particles of masses $2\ \mathrm{kg}$ and $3\ \mathrm{kg}$ are connected by a light inextensible string passing over a smooth fixed pulley (Atwood's machine). The acceleration of the system is:
\begin{enumerate}
\item $\SI{1.96}{m.s^{-2}}$ \quad \item $\SI{4.9}{m.s^{-2}}$ \quad \item $\SI{9.8}{m.s^{-2}}$ \quad \item $\SI{0.98}{m.s^{-2}}$
\end{enumerate}
\item A light inextensible string connects two particles over a smooth pulley. Which statement is correct?
\begin{enumerate}
\item The tensions on the two sides of the pulley are different.
\item The tension is the same throughout the string.
\item The string can push as well as pull.
\item The acceleration of the two particles may differ in magnitude.
\end{enumerate}
\item A particle rests on a rough horizontal plane and is pulled by a horizontal string. Motion is about to begin. The friction force equals:
\begin{enumerate}
\item $\mu R$ \quad \item $\mu mg$ always \quad \item zero \quad \item $mg$
\end{enumerate}
\item A car tows a trailer along a horizontal road at constant speed. The tension in the tow-bar is:
\begin{enumerate}
\item zero \quad \item equal to the resistance on the trailer \quad \item equal to the driving force \quad \item equal to the total weight
\end{enumerate}
\end{enumerate}
\end{exercice}

\begin{exercice}[True or false?]
State whether each statement is true or false, justifying your answer in one or two sentences.
\begin{enumerate}
\item In a system of two particles joined by a light inextensible string over a smooth pulley, the magnitudes of the accelerations of the two particles are equal.
\item The tension in a string can be negative.
\item On a rough plane, the friction force is always equal to $\mu R$.
\item When a hanging particle in a connected system strikes the ground, the other particle stops immediately.
\end{enumerate}
\end{exercice}

\begin{exercice}[Key definitions]
\begin{enumerate}
\item Explain what is meant by a \emph{light inextensible string} and state two consequences of this model for the tension and the accelerations of the connected particles.
\item Explain what is meant by a \emph{smooth pulley} and state the consequence for the tensions on either side of it.
\item Define the \emph{coefficient of friction} $\mu$ between two surfaces and state the inequality satisfied by the friction force $F$ when the system is in equilibrium.
\end{enumerate}
\end{exercice}

\begin{exercice}[Direct calculations]
Take $g = \SI{9.8}{m.s^{-2}}$.
\begin{enumerate}
\item A particle of mass $5\ \mathrm{kg}$ hangs freely from a light inextensible string and moves upwards with acceleration $\SI{1.2}{m.s^{-2}}$. Find the tension in the string.
\item A block of mass $8\ \mathrm{kg}$ is pulled along a rough horizontal table ($\mu = 0.3$) by a horizontal force of $40\ \mathrm{N}$. Find the friction force and the acceleration of the block.
\item A particle of mass $4\ \mathrm{kg}$ rests on a smooth plane inclined at $30^{\circ}$ to the horizontal, held by a string parallel to the plane. Find the tension in the string and the normal reaction.
\end{enumerate}
\end{exercice}

\courssub{I practise}

\begin{exercice}[Particles on a smooth table]
A particle of mass $4\ \mathrm{kg}$ rests on a smooth horizontal table, connected by a light inextensible string passing over a smooth pulley at the edge of the table to a particle of mass $6\ \mathrm{kg}$ hanging freely. The system is released from rest. Take $g = \SI{9.8}{m.s^{-2}}$.
\begin{enumerate}
\item Draw a clear force diagram for each particle.
\item Find the acceleration of the system.
\item Find the tension in the string.
\item Find the magnitude of the resultant force exerted by the string on the pulley.
\end{enumerate}
\end{exercice}

\begin{exercice}[Particles on a rough table]
Repeat the previous exercise, but now the table is rough with coefficient of friction $\mu = 0.25$. Take $g = \SI{9.8}{m.s^{-2}}$.
\begin{enumerate}
\item Show that motion does occur when the system is released.
\item Find the new acceleration and the new tension in the string.
\item Find the speed of the $6\ \mathrm{kg}$ mass after it has fallen $0.8\ \mathrm{m}$ from rest.
\end{enumerate}
\end{exercice}

\begin{exercice}[Masses on an inclined plane]
A mass of $5\ \mathrm{kg}$ lies on a plane inclined at $30^{\circ}$ to the horizontal, connected by a light inextensible string parallel to the plane, passing over a smooth pulley at the top of the plane, to a hanging mass of $7\ \mathrm{kg}$. Take $g = \SI{9.8}{m.s^{-2}}$.
\begin{enumerate}
\item Given that the plane is smooth, find the acceleration of the system and the tension in the string.
\item Now suppose the plane is rough with $\mu = 0.2$. Verify the assumed direction of motion, then find the new acceleration and tension.
\end{enumerate}
\end{exercice}

\begin{exercice}[Two masses on a table and one hanging]
Two masses of $2\ \mathrm{kg}$ and $3\ \mathrm{kg}$ rest on a smooth horizontal table, joined by a light inextensible string. The $3\ \mathrm{kg}$ mass is connected by a second light inextensible string passing over a smooth pulley at the edge of the table to a hanging mass of $5\ \mathrm{kg}$. Take $g = \SI{9.8}{m.s^{-2}}$.
\begin{enumerate}
\item Draw the force diagram for each of the three masses.
\item Find the acceleration of the system.
\item Find the tensions $T_1$ (string between the $2\ \mathrm{kg}$ and $3\ \mathrm{kg}$ masses) and $T_2$ (string over the pulley).
\end{enumerate}
\end{exercice}

\courssub{I prepare for the GCE}

\begin{exercice}[Atwood's machine and multi-stage motion \hfill (14 marks)]
An Atwood's machine consists of particles of masses $3\ \mathrm{kg}$ and $5\ \mathrm{kg}$ connected by a light inextensible string passing over a smooth fixed pulley. The system is released from rest with both particles at the same height. Take $g = \SI{9.8}{m.s^{-2}}$.
\begin{enumerate}
\item[(a)] Find the acceleration of the system and the tension in the string. \hfill (4 marks)
\item[(b)] Find the reaction force exerted on the support of the pulley. \hfill (2 marks)
\item[(c)] After the $5\ \mathrm{kg}$ mass has fallen $2\ \mathrm{m}$, it strikes the ground and does not rebound. Find the speed of the $3\ \mathrm{kg}$ mass at this instant. \hfill (3 marks)
\item[(d)] Find the greatest height above its initial position reached by the $3\ \mathrm{kg}$ mass, assuming it does not reach the pulley. \hfill (5 marks)
\end{enumerate}
\end{exercice}

\begin{exercice}[Car and trailer on a hill \hfill (15 marks)]
A car of mass $1200\ \mathrm{kg}$ tows a trailer of mass $400\ \mathrm{kg}$ up a hill inclined at an angle $\sin^{-1}\!\left(\tfrac{1}{14}\right)$ to the horizontal. Resistance forces of $300\ \mathrm{N}$ on the car and $100\ \mathrm{N}$ on the trailer act against the motion. The driving force of the engine is $2600\ \mathrm{N}$. Take $g = \SI{9.8}{m.s^{-2}}$.
\begin{enumerate}
\item[(a)] Draw a force diagram for the car and for the trailer. \hfill (3 marks)
\item[(b)] By considering the car and the trailer separately (or the whole system), find the acceleration. \hfill (4 marks)
\item[(c)] Find the tension in the tow-bar. \hfill (3 marks)
\item[(d)] The car now descends the same hill and brakes with a braking force of $1800\ \mathrm{N}$, the resistances being unchanged. Find the thrust (compression) in the tow-bar. \hfill (5 marks)
\end{enumerate}
\end{exercice}

\begin{exercice}[Rough inclined plane, equilibrium range and energy \hfill (18 marks)]
A block $A$ of mass $6\ \mathrm{kg}$ rests on a rough plane inclined at $\alpha = 25^{\circ}$ to the horizontal, with coefficient of friction $\mu = 0.4$. The block is connected by a light inextensible string, parallel to the plane and passing over a smooth pulley at the top, to a particle $B$ of mass $m$ hanging freely. Take $g = \SI{9.8}{m.s^{-2}}$.
\begin{enumerate}
\item[(a)] Find the range of values of $m$ for which the system remains at rest. \hfill (6 marks)
\item[(b)] For $m = 10\ \mathrm{kg}$, find the acceleration of the system and the tension in the string. \hfill (4 marks)
\item[(c)] Using the work--energy principle, find the speed of the particles after $B$ has descended $1\ \mathrm{m}$ from rest. \hfill (5 marks)
\item[(d)] Discuss briefly how friction in the pulley bearing would affect your answers to parts (b) and (c). \hfill (3 marks)
\end{enumerate}
\end{exercice}

\begin{exercice}[A uniform chain sliding off a table \hfill (14 marks)]
A uniform chain of length $2\ \mathrm{m}$ and mass $4\ \mathrm{kg}$ lies in a straight line on a smooth horizontal table, perpendicular to the edge, with $0.5\ \mathrm{m}$ hanging over the edge. The chain is released from rest. Take $g = \SI{9.8}{m.s^{-2}}$.
\begin{enumerate}
\item[(a)] Show that, when a length $x$ of the chain hangs over the edge, the acceleration of the chain is $a = \dfrac{gx}{2}$. \hfill (4 marks)
\item[(b)] By writing $a = v\dfrac{dv}{dx}$ and integrating, find the speed of the chain at the instant it just leaves the table. \hfill (6 marks)
\item[(c)] Verify your answer to part (b) using the conservation of mechanical energy, taking the table top as the reference level for potential energy. \hfill (4 marks)
\end{enumerate}
\end{exercice}

\newpage

\coursec{Solutions to the exercises}

\begin{corrige}[Exercise 1 — Multiple choice questions]
\begin{enumerate}
\item For an Atwood's machine, applying Newton's second law to each mass and adding gives $a = \dfrac{m_2-m_1}{m_1+m_2}g = \dfrac{3-2}{2+3}\times 9.8 = \dfrac{1}{5}\times 9.8 = 1.96\,\mathrm{m\,s^{-2}}$. \textbf{Answer: (a)}.
\item For a \cle{light inextensible string} passing over a \cle{smooth pulley}, the tension is the same throughout the string. A string can only pull (never push), and because the string is inextensible the magnitudes of the accelerations of the two particles are equal. \textbf{Answer: (b)}.
\item When motion is about to begin (\cle{limiting equilibrium}), the friction force attains its maximum value $F_{\max} = \mu R$. On a horizontal plane with no vertical acceleration, $R = mg$, so $F_{\max} = \mu mg$. \textbf{Answer: (a)}.
\item At constant speed the acceleration is zero, so the resultant force on the trailer is zero. Horizontally the only forces on the trailer are the tension $T$ in the tow-bar and the resistance; hence $T$ equals the resistance force on the trailer. \textbf{Answer: (b)}.
\end{enumerate}
\end{corrige}

\begin{corrige}[Exercise 2 — True or false?]
\begin{enumerate}
\item \textbf{True.} The string is inextensible, so the length of string between the particles is fixed. Differentiating twice, the magnitudes of their velocities and accelerations along the string are equal (their directions in space may differ, e.g. one horizontal, one vertical).
\item \textbf{False.} Tension is a pulling force; a string cannot push. In the model of a light inextensible string the tension always satisfies $T \ge 0$. A negative value would correspond to compression, which a string cannot sustain — it would simply go slack.
\item \textbf{False.} The friction force satisfies $F \le \mu R$. Equality $F = \mu R$ holds only when sliding occurs (kinetic friction) or at limiting equilibrium (maximum static friction). In static equilibrium below the limit, $F$ takes exactly the value needed to prevent motion.
\item \textbf{False.} When the hanging particle strikes the ground, the string goes slack and the tension vanishes. The other particle continues with the velocity it had at that instant, decelerating only under the remaining forces (friction, or a component of gravity on a slope). It stops immediately only if an impulsive force acts on it.
\end{enumerate}
\end{corrige}

\begin{corrige}[Exercise 3 — Key definitions]
\begin{enumerate}
\item A \emph{light inextensible string} is a string whose mass is negligible compared with the masses it connects, and whose length does not change under tension.
Consequences:
\begin{itemize}
\item Because it is light, the tension is the same at every point of the string (provided any pulley it passes over is smooth).
\item Because it is inextensible, the magnitudes of the velocities and of the accelerations of the connected particles, measured along the string, are equal.
\end{itemize}
\item A \emph{smooth pulley} is a pulley with a frictionless axle and frictionless contact with the string. The consequence is that the tension in the string has the same magnitude on both sides of the pulley.
\item The \emph{coefficient of friction} $\mu$ between two surfaces is the ratio of the limiting friction force to the normal reaction: $\mu = \dfrac{F_{\max}}{R}$. While the surfaces are at rest relative to each other, the friction force satisfies $F \le \mu R$; during sliding, $F = \mu R$ and opposes the relative motion.
\end{enumerate}
\end{corrige}

\begin{corrige}[Exercise 4 — Direct calculations]
Take $g = 9.8\,\mathrm{m\,s^{-2}}$.
\begin{enumerate}
\item Forces on the $5\,\mathrm{kg}$ particle: tension $T$ upward, weight $5g$ downward. Acceleration $a = 1.2\,\mathrm{m\,s^{-2}}$ upward. Newton's second law (upward positive):
\[T - 5g = 5a \implies T = 5(g+a) = 5(9.8+1.2) = 5\times 11 = 55\,\mathrm{N}.\]
\item Block of mass $8\,\mathrm{kg}$, $\mu = 0.3$, horizontal pull $P = 40\,\mathrm{N}$.
Normal reaction: $R = 8g = 78.4\,\mathrm{N}$.
Maximum (limiting) friction: $F_{\max} = \mu R = 0.3\times 78.4 = 23.52\,\mathrm{N}$.
Since $P = 40\,\mathrm{N} > F_{\max}$, the block slides. Kinetic friction: $F = 23.52\,\mathrm{N}$.
Resultant force: $40 - 23.52 = 16.48\,\mathrm{N}$.
Acceleration: $a = \dfrac{16.48}{8} = 2.06\,\mathrm{m\,s^{-2}}$.
\item Mass $4\,\mathrm{kg}$ on a smooth plane inclined at $30^\circ$, held at rest by a string parallel to the plane.
Resolve perpendicular to the plane: $R = 4g\cos30^\circ = 4\times 9.8\times \dfrac{\sqrt{3}}{2} = 19.6\sqrt{3} \approx 33.9\,\mathrm{N}$.
Resolve parallel to the plane (equilibrium): $T - 4g\sin30^\circ = 0 \implies T = 4g\sin30^\circ = 4\times 9.8\times 0.5 = 19.6\,\mathrm{N}$.
\end{enumerate}
\end{corrige}

\begin{corrige}[Exercise 5 — Particles on a smooth table]
Masses: $m_1 = 4\,\mathrm{kg}$ on the table, $m_2 = 6\,\mathrm{kg}$ hanging. Smooth table, smooth pulley, light inextensible string. $g = 9.8\,\mathrm{m\,s^{-2}}$.
\begin{enumerate}
\item \textbf{Force diagrams:}
\begin{itemize}
\item $4\,\mathrm{kg}$ particle: weight $4g$ vertically downward, normal reaction $R$ vertically upward, tension $T$ horizontal towards the pulley.
\item $6\,\mathrm{kg}$ particle: weight $6g$ downward, tension $T$ upward.
\end{itemize}
\item Equations of motion (both particles have acceleration of magnitude $a$):
For the $4\,\mathrm{kg}$ particle (horizontal): $T = 4a$.
For the $6\,\mathrm{kg}$ particle (downward positive): $6g - T = 6a$.
Adding eliminates $T$: $6g = 10a \implies a = \dfrac{6\times 9.8}{10} = 5.88\,\mathrm{m\,s^{-2}}$.
\item Tension: $T = 4a = 4\times 5.88 = 23.52\,\mathrm{N}$.
\item The string exerts two forces of magnitude $T$ on the pulley: one horizontal (towards the table) and one vertically downward. These are perpendicular, so the resultant has magnitude
\[\sqrt{T^2 + T^2} = T\sqrt{2} = 23.52\sqrt{2} \approx 33.3\,\mathrm{N},\]
acting at $45^\circ$ below the horizontal, directed towards the edge of the table.
\end{enumerate}
\end{corrige}

\begin{corrige}[Exercise 6 — Particles on a rough table]
Same arrangement as Exercise 5, but the table is rough with $\mu = 0.25$. $g = 9.8\,\mathrm{m\,s^{-2}}$.
\begin{enumerate}
\item \textbf{Does motion occur?} Normal reaction on the $4\,\mathrm{kg}$ block: $R = 4g = 39.2\,\mathrm{N}$.
Maximum available friction: $F_{\max} = \mu R = 0.25\times 39.2 = 9.8\,\mathrm{N}$.
If the system were held in equilibrium, the string would have to transmit the full weight of the hanging mass, $6g = 58.8\,\mathrm{N}$, to the block. Since $58.8 > 9.8$, friction cannot prevent motion: \textbf{the system moves}.
\item With sliding, friction takes its limiting value $F = 9.8\,\mathrm{N}$, opposing the motion of the block.
$4\,\mathrm{kg}$ block: $T - 9.8 = 4a$.
$6\,\mathrm{kg}$ mass: $58.8 - T = 6a$.
Adding: $58.8 - 9.8 = 10a \implies 49 = 10a \implies a = 4.9\,\mathrm{m\,s^{-2}}$.
Tension: $T = 4a + 9.8 = 19.6 + 9.8 = 29.4\,\mathrm{N}$.
\item Speed after the hanging mass has fallen $0.8\,\mathrm{m}$ from rest:
\[v^2 = u^2 + 2as = 0 + 2\times 4.9\times 0.8 = 7.84 \implies v = \sqrt{7.84} = 2.8\,\mathrm{m\,s^{-1}}.\]
\end{enumerate}
\end{corrige}

\begin{corrige}[Exercise 7 — Masses on an inclined plane]
$m_A = 5\,\mathrm{kg}$ on a plane inclined at $\alpha = 30^\circ$, connected by a string parallel to the plane, over a smooth pulley at the top, to $m_B = 7\,\mathrm{kg}$ hanging freely. $g = 9.8\,\mathrm{m\,s^{-2}}$.
\begin{enumerate}
\item \textbf{Smooth plane.}
Component of $A$'s weight down the plane: $5g\sin30^\circ = 24.5\,\mathrm{N}$. Since $7g = 68.6 > 24.5$, assume $B$ descends and $A$ moves up the plane with acceleration of magnitude $a$.
$A$ (up the plane): $T - 24.5 = 5a$.
$B$ (downward): $68.6 - T = 7a$.
Adding: $68.6 - 24.5 = 12a \implies 44.1 = 12a \implies a = 3.675\,\mathrm{m\,s^{-2}}$.
Tension: $T = 24.5 + 5\times 3.675 = 24.5 + 18.375 = 42.875\,\mathrm{N}$.
\item \textbf{Rough plane, $\mu = 0.2$.}
Normal reaction on $A$: $R = 5g\cos30^\circ = 5\times 9.8\times \dfrac{\sqrt{3}}{2} \approx 42.44\,\mathrm{N}$.
Maximum friction: $F_{\max} = \mu R = 0.2\times 42.44 \approx 8.49\,\mathrm{N}$.
\textbf{Direction of motion:} if $A$ moves up the plane, the forces resisting this are $24.5 + 8.49 = 32.99\,\mathrm{N}$. Since $7g = 68.6 > 32.99$, $B$ does descend and $A$ moves up the plane; friction on $A$ acts down the plane.
Equations:
$A$: $T - 5g\sin30^\circ - \mu(5g\cos30^\circ) = 5a$, i.e. $T - 24.5 - 8.49 = 5a$.
$B$: $68.6 - T = 7a$.
Adding: $68.6 - 24.5 - 8.49 = 12a \implies 35.61 = 12a \implies a \approx 2.97\,\mathrm{m\,s^{-2}}$.
Tension: $T = 68.6 - 7\times 2.97 \approx 68.6 - 20.77 \approx 47.8\,\mathrm{N}$.
\end{enumerate}
\end{corrige}

\begin{corrige}[Exercise 8 — Two masses on a table and one hanging]
$m_1 = 2\,\mathrm{kg}$ and $m_2 = 3\,\mathrm{kg}$ on a smooth table, joined by a light inextensible string; $m_2$ is joined by a second string passing over a smooth pulley at the edge to $m_3 = 5\,\mathrm{kg}$ hanging freely. $g = 9.8\,\mathrm{m\,s^{-2}}$.
\begin{enumerate}
\item \textbf{Force diagrams (horizontal forces only for the table masses):}
\begin{itemize}
\item $2\,\mathrm{kg}$: tension $T_1$ pulling towards $m_2$ (the only horizontal force).
\item $3\,\mathrm{kg}$: tension $T_1$ pulling back towards $m_1$, tension $T_2$ pulling towards the pulley.
\item $5\,\mathrm{kg}$: weight $5g$ downward, tension $T_2$ upward.
\end{itemize}
Note that the two strings are different, so $T_1 \ne T_2$ in general.
\item All three masses move with acceleration of magnitude $a$ (both strings inextensible):
$2\,\mathrm{kg}$: $T_1 = 2a$.
$3\,\mathrm{kg}$: $T_2 - T_1 = 3a$.
$5\,\mathrm{kg}$: $5g - T_2 = 5a$.
Adding all three equations eliminates both tensions:
\[5g = (2+3+5)a = 10a \implies a = \frac{5\times 9.8}{10} = 4.9\,\mathrm{m\,s^{-2}}.\]
\item $T_1 = 2a = 9.8\,\mathrm{N}$; $\quad T_2 = 5g - 5a = 49 - 24.5 = 24.5\,\mathrm{N}$.
Check with the middle equation: $T_2 - T_1 = 24.5 - 9.8 = 14.7 = 3\times 4.9$ \checkmark.
\end{enumerate}
\end{corrige}

\begin{corrige}[Exercise 9 — Atwood's machine and multi-stage motion \hfill (14 marks)]
Masses $3\,\mathrm{kg}$ and $5\,\mathrm{kg}$ connected by a light inextensible string over a smooth pulley; released from rest at the same height. $g = 9.8\,\mathrm{m\,s^{-2}}$.
\begin{enumerate}
\item[(a)] \textbf{Acceleration and tension. (4 marks)}
Let $a$ be the acceleration of the $5\,\mathrm{kg}$ mass downward (and of the $3\,\mathrm{kg}$ mass upward).
$5\,\mathrm{kg}$: $5g - T = 5a$.
$3\,\mathrm{kg}$: $T - 3g = 3a$.
Adding: $2g = 8a \implies a = \dfrac{g}{4} = \dfrac{9.8}{4} = 2.45\,\mathrm{m\,s^{-2}}$.
$T = 3g + 3a = 3(g+a) = 3(9.8+2.45) = 3\times 12.25 = 36.75\,\mathrm{N}$.
\textit{Mark scheme: M1 for both equations of motion, A1 for $a = 2.45$, M1 for substituting back, A1 for $T = 36.75\,\mathrm{N}$.}
\item[(b)] \textbf{Reaction on the pulley support. (2 marks)}
The string exerts two downward forces, each of magnitude $T$, on the pulley. Since the pulley is light and in equilibrium, the support reaction is
\[R = 2T = 2\times 36.75 = 73.5\,\mathrm{N}.\]
\textit{Mark scheme: M1 for $R = 2T$, A1 for $73.5\,\mathrm{N}$.}
\item[(c)] \textbf{Speed when the $5\,\mathrm{kg}$ mass hits the ground after falling $2\,\mathrm{m}$. (3 marks)}
\[v^2 = u^2 + 2as = 0 + 2\times 2.45\times 2 = 9.8 \implies v = \sqrt{9.8} \approx 3.13\,\mathrm{m\,s^{-1}}.\]
Both particles have this speed at that instant (inextensible string).
\textit{Mark scheme: M1 for $v^2 = u^2 + 2as$, A1 for $v^2 = 9.8$, A1 for $v \approx 3.13\,\mathrm{m\,s^{-1}}$.}
\item[(d)] \textbf{Greatest height reached by the $3\,\mathrm{kg}$ mass above its initial position. (5 marks)}
While the string is taut, the $3\,\mathrm{kg}$ mass rises $2\,\mathrm{m}$ (the same distance the $5\,\mathrm{kg}$ mass falls).
When the $5\,\mathrm{kg}$ mass hits the ground, the string goes slack ($T = 0$). The $3\,\mathrm{kg}$ mass continues upward as a free projectile with initial speed $u = \sqrt{9.8}\,\mathrm{m\,s^{-1}}$ under gravity alone.
Extra height gained: $0 = u^2 - 2gh \implies h = \dfrac{u^2}{2g} = \dfrac{9.8}{2\times 9.8} = 0.5\,\mathrm{m}$.
Greatest height above the initial position: $2 + 0.5 = 2.5\,\mathrm{m}$.
\textit{Mark scheme: M1 for recognising the string goes slack, M1 for $v^2 = u^2 - 2gh$, A1 for $h = 0.5\,\mathrm{m}$, M1 for adding the $2\,\mathrm{m}$ already gained, A1 for $2.5\,\mathrm{m}$.}
\end{enumerate}
\end{corrige}

\begin{corrige}[Exercise 10 — Car and trailer on a hill \hfill (15 marks)]
Car $m_C = 1200\,\mathrm{kg}$, trailer $m_T = 400\,\mathrm{kg}$, hill inclined at $\theta$ with $\sin\theta = \dfrac{1}{14}$. Resistances: $300\,\mathrm{N}$ on the car, $100\,\mathrm{N}$ on the trailer. Driving force $F = 2600\,\mathrm{N}$ acting up the hill. $g = 9.8\,\mathrm{m\,s^{-2}}$.
Useful values: $1200g\sin\theta = 1200\times 9.8\times \dfrac{1}{14} = 840\,\mathrm{N}$; $\quad 400g\sin\theta = 400\times 9.8\times \dfrac{1}{14} = 280\,\mathrm{N}$.
\begin{enumerate}
\item[(a)] \textbf{Force diagrams. (3 marks)}
\begin{itemize}
\item Car: weight $1200g$ vertically down, normal reaction perpendicular to the plane, driving force $2600\,\mathrm{N}$ up the plane, resistance $300\,\mathrm{N}$ down the plane, tension $T$ in the tow-bar down the plane (the trailer pulls back on the car).
\item Trailer: weight $400g$ vertically down, normal reaction perpendicular to the plane, resistance $100\,\mathrm{N}$ down the plane, tension $T$ up the plane (the car pulls the trailer forward).
\end{itemize}
\item[(b)] \textbf{Acceleration. (4 marks)}
Treat the car and trailer as a single system of mass $1600\,\mathrm{kg}$ (the tension is internal and cancels):
\[F - 300 - 100 - (1200+400)g\sin\theta = 1600a\]
\[2600 - 400 - 1120 = 1600a \implies 1080 = 1600a \implies a = 0.675\,\mathrm{m\,s^{-2}}.\]
\textit{Mark scheme: M1 for resolving the weights along the plane, M1 for the resultant-force equation, A1 for $1080\,\mathrm{N}$, A1 for $a = 0.675\,\mathrm{m\,s^{-2}}$.}
\item[(c)] \textbf{Tension in the tow-bar. (3 marks)}
Apply Newton's second law to the trailer alone (up the plane positive):
\[T - 100 - 400g\sin\theta = 400a \implies T - 100 - 280 = 400\times 0.675\]
\[T = 380 + 270 = 650\,\mathrm{N}.\]
\textit{Mark scheme: M1 for the trailer equation, A1 for $400g\sin\theta = 280\,\mathrm{N}$, A1 for $T = 650\,\mathrm{N}$.}
\item[(d)] \textbf{Thrust in the tow-bar when descending with braking force $1800\,\mathrm{N}$. (5 marks)}
The vehicle now moves \emph{down} the plane, so the braking force ($1800\,\mathrm{N}$ on the car) and both resistances act \emph{up} the plane, while the weight components act down the plane. Let $P$ be the thrust (compression) in the tow-bar: the car pushes the trailer down the plane and the trailer pushes the car up the plane. Take down the plane as positive; let the acceleration be $a$.
Car: $840 - 300 - 1800 + P = 1200a \implies P - 1260 = 1200a$.
Trailer: $280 - 100 - P = 400a \implies 180 - P = 400a$.
Adding: $-1080 = 1600a \implies a = -0.675\,\mathrm{m\,s^{-2}}$.
The negative sign means the acceleration is directed up the plane: the vehicle is decelerating as it descends, which is consistent with braking.
Then $P = 180 - 400a = 180 - 400(-0.675) = 180 + 270 = 450\,\mathrm{N}$.
The thrust in the tow-bar is $450\,\mathrm{N}$.
\textit{Mark scheme: M1 for correct force directions, M1 for both equations, A1 for $a = -0.675$, M1 for solving for $P$, A1 for $P = 450\,\mathrm{N}$.}
\end{enumerate}
\end{corrige}

\begin{corrige}[Exercise 11 — Rough inclined plane, equilibrium range and energy \hfill (18 marks)]
Block $A$ of mass $6\,\mathrm{kg}$ on a rough plane inclined at $\alpha = 25^\circ$, $\mu = 0.4$; string parallel to the plane over a smooth pulley to a hanging mass $B$ of mass $m\,\mathrm{kg}$. $g = 9.8\,\mathrm{m\,s^{-2}}$.
Preliminary values: $\sin25^\circ \approx 0.4226$, $\cos25^\circ \approx 0.9063$.
$6g\sin25^\circ \approx 6\times 9.8\times 0.4226 \approx 24.85\,\mathrm{N}$; $\quad R = 6g\cos25^\circ \approx 53.29\,\mathrm{N}$; $\quad F_{\max} = \mu R = 0.4\times 53.29 \approx 21.32\,\mathrm{N}$.
\begin{enumerate}
\item[(a)] \textbf{Range of $m$ for equilibrium. (6 marks)}
In equilibrium the tension equals the weight of $B$: $T = mg$. Friction on $A$ can act either way along the plane, up to $F_{\max}$. The two limiting cases are:
\begin{itemize}
\item $A$ on the point of sliding \emph{down} the plane ($m$ smallest): friction acts up the plane at its maximum.
$T + F_{\max} = 6g\sin25^\circ \implies m_{\min}g = 6g\sin25^\circ - \mu(6g\cos25^\circ)$.
$m_{\min} = 6(\sin25^\circ - 0.4\cos25^\circ) = 6(0.4226 - 0.3625) = 6\times 0.0601 \approx 0.361\,\mathrm{kg}$.
\item $A$ on the point of sliding \emph{up} the plane ($m$ largest): friction acts down the plane at its maximum.
$T = 6g\sin25^\circ + F_{\max} \implies m_{\max}g = 6g\sin25^\circ + \mu(6g\cos25^\circ)$.
$m_{\max} = 6(\sin25^\circ + 0.4\cos25^\circ) = 6(0.4226 + 0.3625) = 6\times 0.7851 \approx 4.71\,\mathrm{kg}$.
\end{itemize}
Equilibrium is possible for $0.361\,\mathrm{kg} \le m \le 4.71\,\mathrm{kg}$.
\textit{Mark scheme: B1 for $R = 6g\cos25^\circ$, B1 for $F_{\max} = \mu R$, M1 for the two limiting equations, A1 for $m_{\min}$, A1 for $m_{\max}$, A1 for the range.}
\item[(b)] \textbf{For $m = 10\,\mathrm{kg}$, acceleration and tension. (4 marks)}
Since $10 > m_{\max}$, $B$ descends and $A$ slides up the plane; kinetic friction $F = 21.32\,\mathrm{N}$ acts down the plane on $A$.
$A$ (up the plane): $T - 24.85 - 21.32 = 6a$.
$B$ (downward): $10g - T = 10a$, i.e. $98 - T = 10a$.
Adding: $98 - 24.85 - 21.32 = 16a \implies 51.83 = 16a \implies a \approx 3.24\,\mathrm{m\,s^{-2}}$.
$T = 98 - 10a = 98 - 32.40 \approx 65.6\,\mathrm{N}$.
\textit{Mark scheme: M1 for both equations, A1 for $a \approx 3.24\,\mathrm{m\,s^{-2}}$, M1 for substituting back, A1 for $T \approx 65.6\,\mathrm{N}$.}
\item[(c)] \textbf{Work–energy principle: speed after $B$ has descended $1\,\mathrm{m}$. (5 marks)}
While $B$ descends $1\,\mathrm{m}$, $A$ moves $1\,\mathrm{m}$ up the plane (inextensible string). Both masses have the same speed $v$.
Work done by gravity on $B$: $+10g\times 1 = +98\,\mathrm{J}$.
Work done by gravity on $A$: $-6g\sin25^\circ \times 1 = -24.85\,\mathrm{J}$.
Work done by friction on $A$: $-21.32\times 1 = -21.32\,\mathrm{J}$.
(The tension does no net work on the whole system: it does $+T$ on $A$ and $-T$ on $B$.)
Total work = gain in kinetic energy of both masses:
\[98 - 24.85 - 21.32 = \tfrac{1}{2}(6+10)v^2 \implies 51.83 = 8v^2\]
\[v^2 = 6.479 \implies v \approx 2.55\,\mathrm{m\,s^{-1}}.\]
Check with constant-acceleration formula: $v^2 = 2as = 2\times 3.24\times 1 = 6.48$ \checkmark.
\textit{Mark scheme: M1 for the work–energy equation, A1 for each correct work term, A1 for the KE expression $\frac{1}{2}(16)v^2$, A1 for $v \approx 2.55\,\mathrm{m\,s^{-1}}$.}
\item[(d)] \textbf{Effect of friction in the pulley bearing. (3 marks)}
If the pulley bearing is rough, the tensions on the two sides of the pulley are no longer equal: the tension on the side of the descending mass $B$ exceeds that on the side of $A$. Consequently:
\begin{itemize}
\item the driving effect on $A$ is reduced, so the acceleration found in (b) decreases;
\item part of the work done by gravity is dissipated as heat in the bearing, so the speed found in (c) decreases.
\end{itemize}
\textit{Mark scheme: B1 for stating the tensions differ, B1 for reduced acceleration, B1 for reduced speed / energy dissipated.}
\end{enumerate}
\end{corrige}

\begin{corrige}[Exercise 12 — A uniform chain sliding off a table \hfill (14 marks)]
Uniform chain of length $L = 2\,\mathrm{m}$ and mass $M = 4\,\mathrm{kg}$ on a smooth table; initially a length $x_0 = 0.5\,\mathrm{m}$ hangs over the edge and the chain is released from rest. $g = 9.8\,\mathrm{m\,s^{-2}}$.
Linear density: $\lambda = \dfrac{M}{L} = \dfrac{4}{2} = 2\,\mathrm{kg\,m^{-1}}$.
\begin{enumerate}
\item[(a)] \textbf{Show that $a = \dfrac{gx}{2}$ when a length $x$ hangs over the edge. (4 marks)}
Mass of the hanging part: $\lambda x$; its weight $\lambda x g$ is the only external force driving the motion (the table is smooth, and the tension at the bend is internal to the chain).
The whole chain moves with the same acceleration $a$ (the chain is inextensible), and the total mass is $\lambda L = M$.
Newton's second law for the whole chain:
\[\lambda x g = (\lambda L)\,a \implies a = \frac{x}{L}g = \frac{x}{2}g = \frac{gx}{2}.\]
Note that the acceleration is \emph{not constant}: it increases as more of the chain hangs over the edge.
\textit{Mark scheme: M1 for identifying the driving force $\lambda x g$, M1 for applying $F = ma$ to the whole chain of mass $M$, A1 for the result, A1 for correct use of $\lambda = M/L$.}
\item[(b)] \textbf{Speed when the chain just leaves the table ($x = L = 2\,\mathrm{m}$). (6 marks)}
Since $a$ depends on position, use $a = v\dfrac{dv}{dx}$:
\[v\frac{dv}{dx} = \frac{gx}{2} \implies \int_0^v v\,dv = \int_{0.5}^{2} \frac{gx}{2}\,dx\]
\[\frac{1}{2}v^2 = \frac{g}{2}\left[\frac{x^2}{2}\right]_{0.5}^{2} = \frac{g}{4}\left(4 - 0.25\right) = \frac{g}{4}\times 3.75 = 0.9375g\]
\[v^2 = 1.875g = 1.875\times 9.8 = 18.375 \implies v = \sqrt{18.375} \approx 4.29\,\mathrm{m\,s^{-1}}.\]
\textit{Mark scheme: M1 for $a = v\,\frac{dv}{dx}$, M1 for separating variables with correct limits, A1 for the integrated expression, A1 for $v^2 = 18.375$, A1 for $v \approx 4.29\,\mathrm{m\,s^{-1}}$, B1 for a clear final statement.}
\item[(c)] \textbf{Verification by conservation of energy. (4 marks)}
Take the table top as the reference level for potential energy. The table is smooth, so mechanical energy is conserved.
Initial PE: the hanging part has length $0.5\,\mathrm{m}$, mass $\lambda\times 0.5 = 1\,\mathrm{kg}$, and its centre of mass is $0.25\,\mathrm{m}$ below the table:
$E_{P,i} = -1\times g\times 0.25 = -0.25g$.
Final PE (chain just vertical, length $2\,\mathrm{m}$, mass $4\,\mathrm{kg}$, centre of mass $1\,\mathrm{m}$ below the table):
$E_{P,f} = -4\times g\times 1 = -4g$.
Loss of PE: $E_{P,i} - E_{P,f} = -0.25g - (-4g) = 3.75g$.
Gain of KE: $\dfrac{1}{2}Mv^2 = \dfrac{1}{2}\times 4\times v^2 = 2v^2$.
Conservation of energy: $2v^2 = 3.75g \implies v^2 = 1.875g = 18.375 \implies v \approx 4.29\,\mathrm{m\,s^{-1}}$, in agreement with part (b).
\textit{Mark scheme: M1 for the initial PE (centre of mass of hanging part), M1 for the final PE, A1 for the loss $3.75g$, A1 for $v \approx 4.29\,\mathrm{m\,s^{-1}}$.}
\end{enumerate}
\end{corrige}

\newpage

\coursec{Summary sheet}

\begin{aretenir}[Summary Sheet --- Connected]
\begin{itemize}
\item \textbf{Key definitions:}
\begin{itemize}
\item \textbf{Connected particles}: bodies linked by a light inextensible string.
\item \textbf{Tension} $T$: force transmitted by the string.
\item \textbf{Modelling}: light inextensible string, smooth surfaces, light smooth pulley.
\end{itemize}
\item \textbf{Key formulas:}
\begin{itemize}
\item Same tension $T$ throughout a light string passing over a smooth pulley.
\item Same magnitude of acceleration for particles linked by an inextensible string.
\item Newton's second law applied to each particle: $F = ma$.
\end{itemize}
\item \textbf{Methods:}
\begin{itemize}
\item Draw a separate force diagram for each particle.
\item Choose a positive direction of motion for the whole system.
\item Write $F = ma$ for each particle, then solve the simultaneous equations for $a$ and $T$.
\end{itemize}
\item \textbf{Common mistakes:}
\begin{itemize}
\item Forgetting that the tension acts on \emph{both} particles, in opposite directions.
\item Using different positive directions inconsistently for the two particles.
\item Assuming the tension equals the weight of one of the particles.
\end{itemize}
\item \textbf{GCE tips:}
\begin{itemize}
\item Always state your modelling assumptions (light, inextensible, smooth).
\item Check your answer: the heavier side must accelerate downwards.
\item Give the acceleration and tension to a sensible number of significant figures, with units.
\end{itemize}
\end{itemize}
\end{aretenir}

\findecours

\end{document}
